% Lesson 37 — Listening: Texts on volunteering in gender equality promotion activities — Anglais Tle A — Cours signé OnBuch+
% Programme officiel MINESEC. Compiler avec : tectonic EN37-listening-texts-on-volunteering-in-gender-equality-prom.tex
\newcommand{\DOCMATIERE}{Anglais}
\newcommand{\DOCNIVEAU}{Tle A}
\newcommand{\DOCMODULE}{Chapter 4 : Citizenship and Human Rights}
\newcommand{\DOCLECON}{EN37}
\newcommand{\DOCTITRE}{Lesson 37 — Listening: Texts on volunteering in gender equality promotion activities}
% OnBuch+ — préambule « cours signés » ENRICHI (compilé avec tectonic / XeLaTeX)
% Système visuel « Pop » d'OnBuch+ : fond crème (jamais blanc), bordures encre
% épaisses, ombres « dures » décalées, titres Archivo Black en capitales.
% Version MATHÉMATIQUES : graphes (pgfplots/tikz), boîtes pédagogiques
% (définition, propriété, méthode, attention, le savais-tu, exercice résolu,
% exercice, TP, bilan), page de garde et signature OnBuch+ ancrée en bas.
%
% Macros attendues dans main.tex AVANT \input{preamble} :
%   \DOCMATIERE  \DOCNIVEAU  \DOCMODULE  \DOCLECON (numéro)  \DOCTITRE
\documentclass[11pt]{article}

\usepackage{fontspec}
\usepackage[english,french]{babel} % français = langue principale ; l'anglais via \eng{..} / textebox
\usepackage[a4paper,margin=2cm,top=3.4cm,bottom=2.8cm,headheight=1.4cm,footskip=1.3cm]{geometry}
\usepackage{amsmath,amssymb,amsfonts}
\usepackage{mathtools} % \xmapsto, \coloneqq... souvent utilisés par les modèles
\usepackage{cancel} % simplifications barrées (bilans de forces)
% \abs{z} (module, valeur absolue) et \norm{u} (norme), utilisés par les modèles
\providecommand{\abs}{}\let\abs\relax\DeclarePairedDelimiter{\abs}{\lvert}{\rvert}
\providecommand{\norm}{}\let\norm\relax\DeclarePairedDelimiter{\norm}{\lVert}{\rVert}
\usepackage[table]{xcolor} % [table] : fournit aussi \rowcolors (alternance de lignes)
\usepackage{pagecolor}
\usepackage{tikz}
\usetikzlibrary{arrows.meta,positioning,calc,shapes.geometric,shapes.misc,shapes.arrows,decorations.pathmorphing,decorations.pathreplacing,decorations.markings,patterns,shadows,mindmap,trees,fit,backgrounds,matrix,intersections,angles,quotes}
\usepackage{pgfplots}
\usepackage[version=4]{mhchem} % équations nucléaires \ce{^{235}_{92}U}
\usepackage[european,siunitx]{circuitikz} % schémas électriques
\pgfplotsset{compat=1.18}
\usepgfplotslibrary{fillbetween,groupplots} % aires entre courbes (calcul intégral)
\usepackage{enumitem}
\usepackage{fancyhdr}
% [most] charge toutes les bibliothèques tcolorbox (listings, minted, raster,
% documentation...) et gonfle inutilement la mémoire TeX sur un document long
% (~300 boîtes) : on ne charge que ce qui est réellement utilisé.
\usepackage[skins,breakable]{tcolorbox}
\usepackage{graphicx}
\usepackage{colortbl}
\usepackage{booktabs}
\usepackage{tabularx}
\usepackage{diagbox} % case d'en-tête diagonale des tableaux à double entrée
% Colonnes extensibles centrée / gauche / droite (C, L, R), souvent utilisées par les modèles
\newcolumntype{C}{>{\centering\arraybackslash}X}
\newcolumntype{L}{>{\raggedright\arraybackslash}X}
\newcolumntype{R}{>{\raggedleft\arraybackslash}X}
\usepackage{array}
\usepackage{multicol}
\usepackage{multirow}
\usepackage{siunitx}
% parse-numbers=false : accepte aussi \SI{2,5 \times 10^{-3}}{...}
\sisetup{locale=FR,output-decimal-marker={,},group-minimum-digits=4,parse-numbers=false}
\DeclareSIUnit\ohm{\ensuremath{\symup{\Omega}}} % U+2126 absent de la police
\DeclareSIUnit\division{div} % réglages d'oscilloscope (ms/div, V/div)
\DeclareSIUnit\tr{tr} % tours (vitesse de rotation, ex. \SI{1500}{\tr\per\minute})
\DeclareSIUnit\molar{M} % concentration molaire (ex. \SI{3}{\molar}, \SI{1}{\micro\molar})
\DeclareSIUnit\an{an} % année (le modèle confond parfois avec \year, un compteur TeX) — ex. \SI{5}{\kilo\meter\per\an}
\DeclareSIUnit\FCFA{FCFA} % franc CFA (ex. \SI{500}{\FCFA\per\kilo\gram})
\DeclareSIUnit\degreeBrix{\textdegree Brix} % teneur en sucre d'un sirop/jus (réfractométrie)
\DeclareSIUnit\month{mois} % le modèle utilise parfois \month, un compteur TeX, comme unité de durée
\usepackage{qrcode}
% \degree : commande fréquemment utilisée par les modèles pour un angle ou une
% température en dehors de \SI{}{} (non fournie par les paquets chargés).
\providecommand{\degree}{^{\circ}}
% \qed (fin de démonstration), utilisé par les modèles sans amsthm
\providecommand{\qed}{\hfill$\square$}
% \barre : double barre des valeurs interdites dans un tableau de variations (tkz-tab)
\providecommand{\barre}{\ensuremath{\|}}
% environnement proof (amsthm non chargé) utilisé par les modèles
\ifdefined\proof\else\newenvironment{proof}[1][Démonstration]{\par\noindent\textit{#1.}\ }{\qed\par}\fi
\usepackage{lastpage}
\usepackage{hyperref}
\hypersetup{hidelinks}

% ============================================================
% Palette « Pop » OnBuch+ — mêmes hex que la classe P (plus_ui.dart)
% ============================================================
\definecolor{poporange}{HTML}{F59321}
\definecolor{popdark}{HTML}{E07A0C}
\definecolor{popcream}{HTML}{FFF6E8}
\definecolor{popink}{HTML}{1C1714}
\definecolor{poppurple}{HTML}{7A5AE0}
\definecolor{popgreen}{HTML}{2E9E6A}
\definecolor{poppink}{HTML}{E0457B}
\definecolor{popblue}{HTML}{2F7DE1}
\definecolor{popgold}{HTML}{C98A12}
\definecolor{popmuted}{HTML}{8A7F76}
\definecolor{popline}{HTML}{EADFD2}
\definecolor{popgray}{HTML}{8A7F76}
\colorlet{popgrey}{popgray}
% Alias tolérants (le modèle nomme parfois les couleurs différemment)
\colorlet{popwhite}{white}
\colorlet{popblack}{popink}
\colorlet{popred}{poppink}
\colorlet{popyellow}{popgold}
% Teintes claires (fonds de tableaux/encadrés) — le modèle nomme parfois
% les couleurs avec un suffixe L (ex. popblueL) sans les avoir définies.
\colorlet{poporangeL}{poporange!20}
\colorlet{popdarkL}{popdark!20}
\colorlet{poppurpleL}{poppurple!20}
\colorlet{popgreenL}{popgreen!20}
\colorlet{poppinkL}{poppink!20}
\colorlet{popblueL}{popblue!20}
\colorlet{popgoldL}{popgold!20}
\colorlet{popredL}{popred!20}
\colorlet{popgrayL}{popgray!20}
% Teintes claires pour fonds de tableaux / zones
\colorlet{popblueL}{popblue!10!white}
\colorlet{poporangeL}{poporange!14!white}
\colorlet{popgreenL}{popgreen!12!white}
\colorlet{poppurpleL}{poppurple!12!white}
\colorlet{poppinkL}{poppink!12!white}
\colorlet{popgoldL}{popgold!14!white}

\pagecolor{popcream}

% ============================================================
% Typographie
% ============================================================
\defaultfontfeatures{Ligatures=TeX}
\setmainfont{PlusJakartaSans-Medium.ttf}[Path=../../../fonts/,BoldFont=PlusJakartaSans-Bold.ttf]
% Symboles, API et emojis absents de la police principale : repli sur DejaVu Sans (caractères actifs)
\newfontface\symfont{DejaVuSans.ttf}[Path=../../../fonts/]
\makeatletter
\newcommand\symactive[1]{\catcode#1=13 \begingroup\lccode`\~=#1 \lowercase{\endgroup\def~}{{\symfont\char#1}}}
\newcommand\symblank[1]{\catcode#1=13 \begingroup\lccode`\~=#1 \lowercase{\endgroup\def~}{}}
\newcommand\symhyph[1]{\catcode#1=13 \begingroup\lccode`\~=#1 \lowercase{\endgroup\def~}{\mbox{-}}}
\newcount\symc
\newcommand\symrange[2]{\symc=#1 \@whilenum\symc<#2 \do{\expandafter\symactive\expandafter{\the\symc}\advance\symc by 1 }}
\newcommand\symblankrange[2]{\symc=#1 \@whilenum\symc<#2 \do{\expandafter\symblank\expandafter{\the\symc}\advance\symc by 1 }}
\makeatother
\symrange{"0250}{"02FF}\symactive{"2713}\symactive{"2714}\symactive{"2717}\symactive{"2718}\symactive{"2610}\symactive{"2611}\symactive{"2612}
\symhyph{"2011}\symhyph{"2E1A}\symblankrange{"1F300}{"1FAFF}\symblank{"FE0F}
% Maths en sans-serif (Fira Math) avec chiffres et lettres droites de Plus
% Jakarta, pour que les formules se fondent dans le texte.
\usepackage[math-style=ISO,bold-style=ISO]{unicode-math}
\setmathfont{FiraMath-Regular.otf}[Path=../../../fonts/]
\setmathfont{PlusJakartaSans-Medium.ttf}[Path=../../../fonts/,range={up/num,up/{Latin,latin},\mathup}]
\usepackage{icomma} % virgule décimale sans espace en mode math
% Caractères mathématiques Unicode absents des polices de texte (Plus Jakarta,
% Archivo Black) : rendus par leur équivalent mathématique, en texte comme en
% formule — sinon ils s'affichent en carré vide (ex. « Arithmétique dans ℤ »).
\usepackage{newunicodechar}
\newunicodechar{ℝ}{\ensuremath{\mathbb{R}}}
\newunicodechar{ℤ}{\ensuremath{\mathbb{Z}}}
\newunicodechar{ℂ}{\ensuremath{\mathbb{C}}}
\newunicodechar{ℕ}{\ensuremath{\mathbb{N}}}
\newunicodechar{ℚ}{\ensuremath{\mathbb{Q}}}
\newunicodechar{𝔻}{\ensuremath{\mathbb{D}}}
\newunicodechar{α}{\ensuremath{\alpha}}
\newunicodechar{β}{\ensuremath{\beta}}
\newunicodechar{γ}{\ensuremath{\gamma}}
\newunicodechar{δ}{\ensuremath{\delta}}
\newunicodechar{ε}{\ensuremath{\varepsilon}}
\newunicodechar{θ}{\ensuremath{\theta}}
\newunicodechar{λ}{\ensuremath{\lambda}}
\newunicodechar{μ}{\ensuremath{\mu}}
\newunicodechar{σ}{\ensuremath{\sigma}}
\newunicodechar{τ}{\ensuremath{\tau}}
\newunicodechar{φ}{\ensuremath{\varphi}}
\newunicodechar{ψ}{\ensuremath{\psi}}
\newunicodechar{ω}{\ensuremath{\omega}}
\newunicodechar{Σ}{\ensuremath{\Sigma}}
% majuscules produites par \MakeUppercase dans les titres (α-aminés -> Α)
\newunicodechar{Α}{\ensuremath{\alpha}}
\newunicodechar{Β}{\ensuremath{\beta}}
\newunicodechar{Γ}{\ensuremath{\Gamma}}
\newunicodechar{Δ}{\ensuremath{\Delta}}
\newunicodechar{Ε}{\ensuremath{\varepsilon}}
\newunicodechar{Θ}{\ensuremath{\Theta}}
\newunicodechar{Λ}{\ensuremath{\Lambda}}
\newunicodechar{Μ}{\ensuremath{\mu}}
\newunicodechar{Τ}{\ensuremath{\tau}}
\newunicodechar{Φ}{\ensuremath{\Phi}}
\newunicodechar{Ψ}{\ensuremath{\Psi}}
\newunicodechar{Ω}{\ensuremath{\Omega}}
\newunicodechar{↦}{\ensuremath{\mapsto}}
\newunicodechar{∈}{\ensuremath{\in}}
\newunicodechar{∉}{\ensuremath{\notin}}
\newunicodechar{∧}{\ensuremath{\wedge}}
\newunicodechar{⇔}{\ensuremath{\Leftrightarrow}}
\newunicodechar{⟺}{\ensuremath{\Longleftrightarrow}}
\newunicodechar{⇒}{\ensuremath{\Rightarrow}}
\newunicodechar{⟹}{\ensuremath{\Longrightarrow}}
\newunicodechar{∘}{\ensuremath{\circ}}
\newunicodechar{∪}{\ensuremath{\cup}}
\newunicodechar{∩}{\ensuremath{\cap}}
\newunicodechar{≡}{\ensuremath{\equiv}}
\newunicodechar{⊂}{\ensuremath{\subset}}
\newunicodechar{⊥}{\ensuremath{\perp}}
\newunicodechar{∃}{\ensuremath{\exists}}
\newunicodechar{∀}{\ensuremath{\forall}}
\newunicodechar{∛}{\ensuremath{\sqrt[3]{\,}}}
\newunicodechar{∜}{\ensuremath{\sqrt[4]{\,}}}
\newunicodechar{⃗}{\ensuremath{{}^{\to}}}
\newunicodechar{ⁿ}{\ifmmode^{n}\else\textsuperscript{\ensuremath{n}}\fi}
\newunicodechar{ˣ}{\ifmmode^{x}\else\textsuperscript{\ensuremath{x}}\fi}
\newunicodechar{ᵇ}{\ifmmode^{b}\else\textsuperscript{\ensuremath{b}}\fi}
\newunicodechar{ʳ}{\ifmmode^{r}\else\textsuperscript{\ensuremath{r}}\fi}
\newunicodechar{ᵘ}{\ifmmode^{u}\else\textsuperscript{\ensuremath{u}}\fi}
\newunicodechar{ʸ}{\ifmmode^{y}\else\textsuperscript{\ensuremath{y}}\fi}
\newunicodechar{ᵃ}{\ifmmode^{a}\else\textsuperscript{\ensuremath{a}}\fi}
\newunicodechar{ᵉ}{\ifmmode^{e}\else\textsuperscript{\ensuremath{e}}\fi}
\newunicodechar{ᵏ}{\ifmmode^{k}\else\textsuperscript{\ensuremath{k}}\fi}
\newunicodechar{ᵖ}{\ifmmode^{p}\else\textsuperscript{\ensuremath{p}}\fi}
\newunicodechar{ᴺ}{\ifmmode^{N}\else\textsuperscript{\ensuremath{N}}\fi}
\newunicodechar{ᶜ}{\ifmmode^{c}\else\textsuperscript{\ensuremath{c}}\fi}
\newunicodechar{ʰ}{\ifmmode^{h}\else\textsuperscript{\ensuremath{h}}\fi}
\newunicodechar{ᵠ}{\ifmmode^{q}\else\textsuperscript{\ensuremath{q}}\fi}
\newunicodechar{ₙ}{\ifmmode_{n}\else\textsubscript{\ensuremath{n}}\fi}
\newunicodechar{ₐ}{\ifmmode_{a}\else\textsubscript{\ensuremath{a}}\fi}
\newunicodechar{ᵢ}{\ifmmode_{i}\else\textsubscript{\ensuremath{i}}\fi}
\newunicodechar{ᵣ}{\ifmmode_{r}\else\textsubscript{\ensuremath{r}}\fi}
\newunicodechar{ₖ}{\ifmmode_{k}\else\textsubscript{\ensuremath{k}}\fi}
\newunicodechar{ᵦ}{\ifmmode_{\beta}\else\textsubscript{\ensuremath{\beta}}\fi}
\newunicodechar{ₓ}{\ifmmode_{x}\else\textsubscript{\ensuremath{x}}\fi}
\newfontfamily\popdisplay{ArchivoBlack-Regular.ttf}[Path=../../../fonts/]
\newfontfamily\poplogo{SpaceGrotesk-Bold.ttf}[Path=../../../fonts/]
\newfontfamily\popsemi{PlusJakartaSans-SemiBold.ttf}[Path=../../../fonts/]
\newfontfamily\popxbold{PlusJakartaSans-ExtraBold.ttf}[Path=../../../fonts/]
\newfontfamily\popxboldls{PlusJakartaSans-ExtraBold.ttf}[Path=../../../fonts/,LetterSpace=3.0]

\setlist[enumerate]{leftmargin=1.7em,itemsep=5pt,topsep=4pt}
\setlist[itemize]{leftmargin=1.7em,itemsep=4pt,topsep=4pt,label={\color{poporange}\textbullet}}
\setlength{\parindent}{0pt}
\setlength{\parskip}{5pt}
\renewcommand{\arraystretch}{1.25}

% pgfplots : style Pop par défaut
\pgfplotsset{
  every axis/.append style={axis line style={popink,line width=1pt},tick style={popink},
    grid style={popline,line width=0.5pt},label style={font=\small},tick label style={font=\footnotesize}},
  popcurve/.style={poporange,line width=1.8pt,no markers,smooth},
}
% Même style disponible dans un \draw TikZ simple (hors pgfplots)
\tikzset{popcurve/.style={poporange,line width=1.8pt}}
\tikzset{popgrid/.style={popline,very thin}}
\colorlet{popcurve}{poporange} % le nom du style est parfois utilisé comme couleur

% ============================================================
% Pill / squiggle
% ============================================================
\newcommand{\poppill}[3]{%
  \tikz[baseline=-0.55ex]\node[draw=popink,line width=1.2pt,rounded corners=2.6pt,%
    fill=#2,inner xsep=6.5pt,inner ysep=3pt]{%
    \popxboldls\fontsize{7}{7}\selectfont\color{#3}\MakeUppercase{#1}};%
}
\newcommand{\popsquiggle}[1]{%
  \tikz[baseline=-0.4ex]{
    \draw[#1,line width=2.6pt,line cap=round]
      plot [smooth] coordinates {(0,0.09) (0.34,0.01) (0.68,0.21) (1.02,0.01) (1.36,0.10)};
  }%
}
% Mot-clé surligné (marqueur orange sous le texte)
\newcommand{\cle}[1]{{\popxbold\color{popdark}#1}}

% ============================================================
% En-tête / pied
% ============================================================
\pagestyle{fancy}
\fancyhf{}
\renewcommand{\headrulewidth}{0pt}
\renewcommand{\footrulewidth}{0pt}
\lhead{\raisebox{-0.15ex}{\poplogo\fontsize{13}{13}\selectfont\color{popink}OnBuch\color{poporange}+}}
\rhead{\poppill{\DOCMATIERE\ \textbullet\ \DOCNIVEAU\ \textbullet\ Leçon \DOCLECON}{white}{popink}}
\lfoot{\popxbold\fontsize{7.6}{7.6}\selectfont\color{popmuted}\MakeUppercase{Cours signé OnBuch+}\ \textbullet\ {\popsemi\DOCTITRE}}
\rfoot{\tikz[baseline=-0.6ex]\node[rounded corners=8.5pt,fill=popink,minimum height=17pt,minimum width=17pt,inner xsep=5pt,inner sep=0]{\popxbold\fontsize{8}{8}\selectfont\color{popcream}\thepage\,/\,\pageref*{LastPage}};}

% ============================================================
% Titres
% ============================================================
\newcommand{\chaptitre}[2]{%
  \begin{tcolorbox}[enhanced,colback=poporange,colframe=popink,boxrule=2.6pt,arc=11pt,
      drop shadow=popink,left=15pt,right=15pt,top=12pt,bottom=11pt,before skip=3pt,after skip=18pt]
    \poppill{#2}{popink}{popcream}\par\vspace{9pt}
    {\raggedright\hyphenpenalty=10000\exhyphenpenalty=10000\popdisplay\fontsize{22}{24}\selectfont\color{popcream}\MakeUppercase{#1}\par}
  \end{tcolorbox}
}
% Section numérotée : pastille orange avec le numéro + titre Archivo
\newcounter{popsec}
\newcommand{\coursec}[1]{%
  \stepcounter{popsec}\setcounter{popsub}{0}%
  \par\vspace{16pt}\noindent
  \begin{minipage}{\linewidth}
  \tikz[baseline=-0.7ex]\node[draw=popink,line width=1.6pt,fill=poporange,rounded corners=4pt,minimum size=22pt,inner sep=0]{\popdisplay\fontsize{12}{12}\selectfont\color{popcream}\thepopsec};%
  \hspace{8pt}{\popdisplay\fontsize{15}{17}\selectfont\color{popink}\MakeUppercase{#1}}\par
  \vspace{4pt}\noindent{\color{poporange}\rule{36pt}{3.6pt}}
  \end{minipage}\par\nopagebreak\vspace{8pt}
}
\newcounter{popsub}
\newcommand{\courssub}[1]{%
  \stepcounter{popsub}%
  \par\vspace{9pt}\noindent{\popxbold\fontsize{11.5}{13}\selectfont\color{popink}{\color{poporange}\thepopsec.\thepopsub}\hspace{6pt}#1}\par\nopagebreak\vspace{4pt}
}

% ============================================================
% Boîtes pédagogiques — même recette : fond blanc, bordure épaisse colorée,
% ombre dure encre, badge plein. Titre optionnel [#1].
% ============================================================
\tcbset{popbase/.style={breakable,enhanced,colback=white,boxrule=2.3pt,arc=9pt,
  shadow={3pt}{-3pt}{0pt}{popink},left=14pt,right=14pt,top=11pt,bottom=11pt,before skip=11pt,after skip=15pt,
  fontupper=\popsemi\fontsize{10.6}{14.5}\selectfont\color{popink},
  fontlower=\popsemi\fontsize{10.6}{14.5}\selectfont\color{popink}}}
% Coche dessinée (le glyphe ✓ n'existe pas dans les polices du document)
\newcommand{\popcheck}[1]{\tikz[baseline=-0.5ex]\draw[#1,line width=1.6pt,line cap=round,line join=round](-0.13,0.0)--(-0.04,-0.09)--(0.13,0.1);}
\providecommand{\coche}{\popcheck{popgreen}}
\providecommand{\resultat}[1]{\par\smallskip\noindent\fcolorbox{popgreen}{popgreenL}{\parbox{\dimexpr\linewidth-2\fboxsep-2\fboxrule\relax}{\textbf{Résultat.} #1}}\par\smallskip}
\newcommand{\popboxhead}[3]{\poppill{#1}{#2}{white}\ifx\relax#3\relax\else\hspace{6pt}{\popxbold\fontsize{10.6}{12}\selectfont\color{popink}#3}\fi\par\vspace{7pt}}

\newtcolorbox{aretenir}[1][]{popbase,colframe=popgreen,before upper={\popboxhead{À retenir}{popgreen}{#1}}}
% Texte en anglais (document de lecture, dialogue, extrait) : cadre
% d'étude. Un titre optionnel (source, auteur) s'affiche dans l'en-tête.
\newtcolorbox{textebox}[1][]{popbase,colframe=popblue,before upper={\popboxhead{Text}{popblue}{#1}\begin{otherlanguage}{english}},after upper={\end{otherlanguage}}}
% Marques ✓ / ✗ (forme correcte / fautive) souvent utilisées par les modèles
\providecommand{\cross}{\textcolor{poppink}{\ensuremath{\times}}}
\providecommand{\xmark}{\cross}\providecommand{\crossmark}{\cross}\providecommand{\cmark}{\ensuremath{\checkmark}}
% Ligne de réponse / justification à compléter (inventée par les modèles)
\providecommand{\justification}[1]{\par\noindent\textit{Justification~:} \dotfill\par}
% \textipa (package tipa non chargé) : la police ne couvre pas l'API -> simple passage
\providecommand{\textipa}[1]{#1}
% Soulignement (ulem non chargé) : \uline, \ul
\providecommand{\uline}[1]{\underline{#1}}\providecommand{\ul}[1]{\underline{#1}}
% \glossaire{\item ...} inventée par les modèles : liste de glossaire
\providecommand{\glossaire}[1]{\begin{itemize}#1\end{itemize}}
% \alert (beamer) inventée par les modèles : texte mis en évidence
\providecommand{\alert}[1]{\textbf{\textcolor{poppink}{#1}}}
% variantes ASCII d'ordinaux écrites en macro
\providecommand{\iere}{\textsuperscript{re}}\providecommand{\ieme}{\textsuperscript{e}}\providecommand{\ere}{\textsuperscript{re}}\providecommand{\eme}{\textsuperscript{e}}\providecommand{\er}{\textsuperscript{er}}\providecommand{\ie}{\textsuperscript{e}}
% Ordinaux français écrits en macro par les modèles : 1\ière, 3\ième
\newcommand{\ière}{\textsuperscript{re}}\newcommand{\ième}{\textsuperscript{e}}
% \exemple (inventée par les modèles) : étiquette « Exemple : »
\providecommand{\exemple}{\textbf{Exemple~:}\ }
% \square utilisable aussi hors mode math (cases à cocher)
\AtBeginDocument{\let\squareMath\square\renewcommand{\square}{\ensuremath{\squareMath}}}
% Mot ou phrase en anglais à l'intérieur d'un paragraphe français
\newcommand{\eng}[1]{\ifmmode\text{\textit{#1}}\else\begingroup\selectlanguage{english}\itshape #1\endgroup\fi}
\let\esp\eng % alias toléré
% flèches d'intonation inventées par les modèles
\providecommand{\textsearrow}{}\renewcommand{\textsearrow}{\ensuremath{\searrow}}\providecommand{\textnearrow}{}\renewcommand{\textnearrow}{\ensuremath{\nearrow}}
% \fr{..} : mot français (inventé par les modèles)
\providecommand{\fr}[1]{\textit{#1}}
\newtcolorbox{exemplebox}[1][]{popbase,colframe=poporange,before upper={\popboxhead{Exemple}{poporange}{#1}}}
\newtcolorbox{definition}[1][]{popbase,colframe=popblue,before upper={\popboxhead{Définition}{popblue}{#1}}}
\newtcolorbox{propriete}[1][]{popbase,colframe=poppurple,before upper={\popboxhead{Propriété}{poppurple}{#1}}}
\newtcolorbox{methode}[1][]{popbase,colframe=popgold,before upper={\popboxhead{Méthode}{popgold}{#1}}}
\newtcolorbox{attention}[1][]{popbase,colframe=poppink,before upper={\popboxhead{Attention}{poppink}{#1}}}
\newtcolorbox{experience}[1][]{popbase,colframe=popblue,colback=popblueL,before upper={\popboxhead{Expérience}{popblue}{#1}}}
% « Le savais-tu ? » : carte violette pleine (contexte, culture, Cameroun)
\newtcolorbox{savaistu}[1][]{popbase,colback=poppurple,colframe=popink,
  fontupper=\popsemi\fontsize{10.4}{14.2}\selectfont\color{white},
  before upper={\poppill{Le savais-tu\,?}{white}{poppurple}\ifx\relax#1\relax\else\hspace{6pt}{\popxbold\color{white}#1}\fi\par\vspace{7pt}}}
% Exercice résolu : énoncé en haut, \tcblower puis la solution sur fond crème
\newcounter{popexo}\newcounter{popexr}
\newtcolorbox{exoresolu}[1][]{popbase,colframe=poporange,colbacklower=popcream,
  segmentation style={popink,line width=1.2pt,dashed},
  before upper={\stepcounter{popexr}\popboxhead{Exercice résolu \thepopexr}{poporange}{#1}},
  before lower={\poppill{Solution}{popink}{popcream}\par\vspace{6pt}}}
% Exercice (non résolu en place) — la correction est dans la partie Corrigés
\newtcolorbox{exercice}[1][]{popbase,colframe=popink,
  before upper={\stepcounter{popexo}\popboxhead{Exercice \thepopexo}{popink}{#1}}}
% Corrigé
\newtcolorbox{corrige}[1][]{popbase,colframe=popgreen,colback=popgreenL,
  before upper={\popboxhead{Corrigé}{popgreen}{#1}}}
% Objectifs / prérequis / bilan
\newtcolorbox{objectifs}{popbase,colframe=popink,colback=poporangeL,
  before upper={\popboxhead{Objectifs de la leçon}{popink}{}}}
\newtcolorbox{prerequis}{popbase,colframe=popink,colback=popblueL,
  before upper={\popboxhead{Prérequis}{popblue}{}}}
\newtcolorbox{bilan}{popbase,colframe=popink,colback=white,boxrule=2.8pt,
  before upper={\popboxhead{Fiche bilan}{poporange}{L'essentiel à connaître pour l'examen}}}
% Figure encadrée avec légende
\newtcolorbox{popfigure}[1][]{enhanced,breakable=false,colback=white,colframe=popline,boxrule=1.4pt,arc=8pt,
  left=10pt,right=10pt,top=10pt,bottom=8pt,before skip=10pt,after skip=12pt,halign=center,
  lower separated=false,halign lower=center,
  fontlower=\popsemi\fontsize{9}{11.5}\selectfont\color{popmuted},
  #1}
\newcommand{\legende}[1]{\par\vspace{6pt}{\popsemi\fontsize{9}{11.5}\selectfont\color{popmuted}#1}}

% ============================================================
% Signature OnBuch+
% ============================================================
\newcommand{\poplogoinline}[1]{{\poplogo\fontsize{#1}{#1}\selectfont\color{popink}OnBuch\color{poporange}+}}
\newcommand{\signatureonbuch}{%
  \begin{tcolorbox}[enhanced,colback=popink,colframe=popink,boxrule=2.2pt,arc=10pt,
      drop shadow=poporange,left=14pt,right=14pt,top=10pt,bottom=10pt,before skip=0pt,after skip=0pt]
    \tikz[baseline=-0.7ex]\node[circle,fill=popgreen,draw=popcream,line width=1.3pt,minimum size=22pt,inner sep=0]{\popcheck{white}};%
    \hspace{9pt}\begin{minipage}[c]{0.62\linewidth}
      {\popxbold\fontsize{12}{13}\selectfont\color{popcream}Signé\ \textbullet\ {\poplogo OnBuch\color{poporange}+}}\par\vspace{2pt}
      {\raggedright\popsemi\fontsize{8}{10.5}\selectfont\color{popline}\MakeUppercase{Équipe pédagogique OnBuch+}\\\MakeUppercase{Conforme au programme officiel MINESEC}\par}
    \end{minipage}\hfill
    {\poplogo\fontsize{22}{22}\selectfont\color{popcream}OnBuch\color{poporange}+}
  \end{tcolorbox}%
}

% ============================================================
% Page de garde
% #1 titre · #2 matière · #3 niveau · #4 module · #5 n° leçon · #6 durée ·
% #7 description / objectifs (texte ou itemize)
% ============================================================
\newcommand{\pagegarde}[7]{%
  \thispagestyle{empty}
  \newgeometry{a4paper,margin=1.7cm,top=1.6cm,bottom=1.4cm}%
  \begin{tikzpicture}[remember picture,overlay]
    \fill[poporange!18] ([xshift=-3.2cm,yshift=-3.6cm]current page.north east) circle (4.6cm);
    \fill[poppurple!14] ([xshift=2.4cm,yshift=7.6cm]current page.south west) circle (3.8cm);
  \end{tikzpicture}%
  {\poplogo\fontsize{30}{30}\selectfont\color{popink}OnBuch\color{poporange}+}\hfill
  \raisebox{0.6ex}{\poppill{Cours signé \textbullet\ Premium}{popink}{popcream}}\par
  \vspace{1pt}\popsquiggle{poppurple}\par
  \vspace{8pt}
  \begin{tcolorbox}[enhanced,colback=poporange,colframe=popink,boxrule=2.8pt,arc=13pt,
      drop shadow=popink,left=18pt,right=18pt,top=12pt,bottom=13pt,after skip=10pt]
    \poppill{#2 \textbullet\ #3}{popink}{popcream}\hspace{5pt}\poppill{#4}{white}{popink}\par\vspace{8pt}
    {\popdisplay\fontsize{14}{14}\selectfont\color{popink}LEÇON #5}\par\vspace{4pt}
    {\raggedright\hyphenpenalty=10000\exhyphenpenalty=10000\popdisplay\fontsize{25}{27}\selectfont\color{popcream}\MakeUppercase{#1}\par}
    \vspace{8pt}
    {\popxbold\fontsize{9.5}{9.5}\selectfont\color{popink}\MakeUppercase{Durée conseillée : #6}}
  \end{tcolorbox}
  \begin{tcolorbox}[enhanced,colback=white,colframe=popink,boxrule=2.2pt,arc=10pt,
      drop shadow=popink,left=16pt,right=16pt,top=10pt,bottom=10pt,after skip=8pt]
    \poppill{Dans cette leçon}{poppurple}{white}\par\vspace{5pt}
    \popsemi\fontsize{9.3}{12.6}\selectfont\color{popink}#7
  \end{tcolorbox}
  \noindent\begin{minipage}[t]{0.487\linewidth}
    \begin{tcolorbox}[enhanced,colback=popgreen,colframe=popink,boxrule=2.2pt,arc=10pt,
        drop shadow=popink,left=11pt,right=11pt,top=10pt,bottom=10pt,height=3.1cm,valign=center]
      \tcbox[colback=white,colframe=popink,boxrule=1.3pt,arc=5pt,boxsep=0pt,nobeforeafter,
        left=3pt,right=3pt,top=3pt,bottom=3pt,valign=center]{\qrcode[height=1.9cm]{https://play.google.com/store/apps/details?id=com.luvvix.onbuch&hl=fr}}%
      \hspace{8pt}\begin{minipage}[c]{0.5\linewidth}
        {\popdisplay\fontsize{11}{12.5}\selectfont\color{white}\MakeUppercase{Télécharge OnBuch}}\par\vspace{4pt}
        {\raggedright\popsemi\fontsize{8.4}{11}\selectfont\color{white}Gratuit \textbullet\ Google Play\\Tuteur IA, annales, concours, résultats\par}
      \end{minipage}
    \end{tcolorbox}
  \end{minipage}\hfill
  \begin{minipage}[t]{0.487\linewidth}
    \begin{tcolorbox}[enhanced,colback=poppurple,colframe=popink,boxrule=2.2pt,arc=10pt,
        drop shadow=popink,left=11pt,right=11pt,top=10pt,bottom=10pt,height=3.1cm,valign=center]
      \tikz[baseline=-0.7ex]\node[circle,fill=white,draw=popink,line width=1.3pt,minimum size=1.6cm,inner sep=0]{\popdisplay\fontsize{24}{24}\selectfont\color{poppurple}+};%
      \hspace{8pt}\begin{minipage}[c]{0.6\linewidth}
        {\popdisplay\fontsize{11}{12.5}\selectfont\color{white}\MakeUppercase{Passe à OnBuch+}}\par\vspace{4pt}
        {\raggedright\popsemi\fontsize{8.4}{11}\selectfont\color{white}Tous les cours signés, Tuteur IA illimité, fiches PDF — onglet OnBuch+ de l'app\par}
      \end{minipage}
    \end{tcolorbox}
  \end{minipage}
  \vspace{8pt}
  \popcontenu
  \vfill
  \signatureonbuch
  \restoregeometry
  \newpage
}

% Rangée « Ce cours contient » (page de garde)
\newcommand{\poptile}[3]{% #1 couleur #2 gros texte #3 légende
  \begin{tcolorbox}[enhanced,colback=#1,colframe=popink,boxrule=2pt,arc=9pt,shadow={2.5pt}{-2.5pt}{0pt}{popink},
      left=6pt,right=6pt,top=9pt,bottom=9pt,halign=center,valign=center,height=1.75cm,width=\linewidth,nobeforeafter]
    {\popdisplay\fontsize{12.5}{13}\selectfont\color{white}\MakeUppercase{#2}}\par\vspace{3pt}
    {\centering\popsemi\fontsize{7.8}{9.5}\selectfont\color{white}#3\par}
  \end{tcolorbox}}
\newcommand{\popcontenu}{%
  \par\noindent{\popxboldls\fontsize{7.5}{7.5}\selectfont\color{popmuted}CE COURS SIGNÉ CONTIENT}\par\vspace{7pt}
  \noindent\begin{minipage}[t]{0.235\linewidth}\poptile{popblue}{Cours}{complet, illustré, pas à pas}\end{minipage}\hfill
  \begin{minipage}[t]{0.235\linewidth}\poptile{poporange}{Méthodes}{et exercices résolus}\end{minipage}\hfill
  \begin{minipage}[t]{0.235\linewidth}\poptile{poppink}{Exercices}{type examen + corrigés}\end{minipage}\hfill
  \begin{minipage}[t]{0.235\linewidth}\poptile{popgreen}{Fiche bilan}{l'essentiel à retenir}\end{minipage}\par}

% Sommaire
\newtcolorbox{sommaire}{enhanced,colback=white,colframe=popink,boxrule=2.2pt,arc=10pt,
  drop shadow=popink,left=18pt,right=18pt,top=13pt,bottom=13pt,before skip=6pt,after skip=20pt,
  before upper={\poppill{Sommaire}{poppurple}{white}\par\vspace{9pt}\popsemi\fontsize{10.6}{14.5}\selectfont\color{popink}}}

% Signature de fin de cours — poussée tout en bas de la dernière page
\newcommand{\findecours}{%
  \par\vfill\nopagebreak
  \begin{center}\popsquiggle{poporange}\end{center}\vspace{4pt}
  \signatureonbuch
}
\AtBeginDocument{\def\llbracket{\mathopen{[\mkern-3.5mu[}}\def\rrbracket{\mathclose{]\mkern-3.5mu]}}} % intervalles d'entiers (glyphe ⟦ absent de la police)
% Glyphes absents de Fira Math : redéfinis à partir de glyphes disponibles
\AtBeginDocument{%
  \def\setminus{\mathbin{\backslash}}\let\smallsetminus\setminus
  \def\vdots{\mathord{\vbox{\baselineskip3.2pt\lineskiplimit0pt\kern4pt\hbox{.}\hbox{.}\hbox{.}}}}%
  \def\ddots{\mathinner{\mkern1mu\raise6.5pt\hbox{.}\mkern2mu\raise3.6pt\hbox{.}\mkern2mu\raise0.7pt\hbox{.}\mkern1mu}}%
  \def\triangle{\mathord{\scalebox{0.9}{$\symup{\Delta}$}}}%
  \def\nabla{\mathord{\raisebox{\depth}{\scalebox{1}[-1]{$\symup{\Delta}$}}}}%
  \def\checkmark{\popcheck{popdark}}%
  \def\star{\mathbin{\ast}}\def\bigstar{\mathord{\ast}}%
  \def\diamond{\mathbin{\scalebox{0.6}{$\Diamond$}}}%
  \def\bigcup{\mathop{\mathchoice{\raisebox{-0.3em}{\scalebox{1.7}{$\cup$}}}{\scalebox{1.25}{$\cup$}}{\cup}{\cup}}\displaylimits}%
  \def\bigcap{\mathop{\mathchoice{\raisebox{-0.3em}{\scalebox{1.7}{$\cap$}}}{\scalebox{1.25}{$\cap$}}{\cap}{\cap}}\displaylimits}%
  \def\lesseqqgtr{\mathrel{\lessgtr}}\def\bigoplus{\mathop{\oplus}\displaylimits}%
}
\providecommand{\ssi}{si et seulement si} % abréviation fréquente
\AtBeginDocument{\providecommand{\wideparen}{\overparen}} % arc de cercle

%% FIGFIX-BEGIN v3 — correctifs globaux des figures (docs/onbuch-generation/tools/figfix)
\usepackage{adjustbox}\usepackage{etoolbox}
% toute figure TikZ/pgfplots plus large que la ligne est réduite à la largeur disponible
\BeforeBeginEnvironment{tikzpicture}{\begin{adjustbox}{max width=\linewidth}}
\AfterEndEnvironment{tikzpicture}{\end{adjustbox}}
% légendes pgfplots lisibles par-dessus les courbes
\pgfplotsset{every axis/.append style={legend style={fill=white,fill opacity=0.92,text opacity=1,draw=black!15,inner sep=3pt}}}
% étiquettes d'arêtes (syntaxe edge["..."]) avec fond blanc
\tikzset{every edge quotes/.append style={fill=white,inner sep=1.2pt}}
% bulles de cartes mentales : texte à la ligne dans la bulle, bulle un peu plus grande
\tikzset{every concept/.append style={text width=2.0cm,align=center,minimum size=2.3cm,inner sep=2pt}}
%% FIGFIX-END
\begin{document}

\pagegarde{Lesson 37 — Listening: Texts on volunteering in gender equality promotion activities}{Anglais}{Tle A}{Chapter 4 : Citizenship and Human Rights}{EN37}{2 h}{In this listening lesson, students work on an original audio script about volunteers promoting gender equality in Cameroon and abroad. They learn to identify the theme, keywords, figures and proper nouns, build thematic vocabulary, and answer varied comprehension questions before reusing the lexicon in short guided productions.
\vspace{4pt}
\begin{multicols}{2}\small
\begin{itemize}
\item À la fin de cette leçon, je sais repérer le thème général d'un enregistrement dès la première écoute.
\item À la fin de cette leçon, je sais identifier les mots-clés, les chiffres et les noms propres dans un texte écouté.
\item À la fin de cette leçon, je sais prendre des notes efficaces pendant l'écoute (tableaux, abréviations).
\item À la fin de cette leçon, je sais répondre à des questions de compréhension de types vrai/faux, QCM et réponses courtes.
\item À la fin de cette leçon, je sais utiliser le lexique du volontariat et de l'égalité des sexes dans des phrases simples.
\item À la fin de cette leçon, je sais distinguer les faux amis et les pièges de prononciation liés au thème.
\end{itemize}\end{multicols}}

\begin{sommaire}
\begin{multicols}{2}
\begin{enumerate}
\item Warm-up: entering the theme
\item Listening script (support text)
\item Comprehension questions
\item Thematic vocabulary
\item Listening strategies and worked example
\item Guided production: reuse the lexicon
\item Activité d'intégration
\item Méthodes et exercices résolus
\item Exercices
\item Corrigés des exercices
\item Fiche bilan
\end{enumerate}
\end{multicols}
\end{sommaire}

\begin{objectifs}
\begin{itemize}
\item À la fin de cette leçon, je sais repérer le thème général d'un enregistrement dès la première écoute.
\item À la fin de cette leçon, je sais identifier les mots-clés, les chiffres et les noms propres dans un texte écouté.
\item À la fin de cette leçon, je sais prendre des notes efficaces pendant l'écoute (tableaux, abréviations).
\item À la fin de cette leçon, je sais répondre à des questions de compréhension de types vrai/faux, QCM et réponses courtes.
\item À la fin de cette leçon, je sais utiliser le lexique du volontariat et de l'égalité des sexes dans des phrases simples.
\item À la fin de cette leçon, je sais distinguer les faux amis et les pièges de prononciation liés au thème.
\item À la fin de cette leçon, je sais résumer oralement ou par écrit les informations essentielles d'un enregistrement.
\item À la fin de cette leçon, je sais produire un court message de sensibilisation au volontariat pour l'égalité des sexes.
\end{itemize}
\end{objectifs}

\begin{prerequis}
\begin{itemize}
\item Vocabulaire de base des droits de l'homme vu en Première (rights, equality, discrimination)
\item Les nombres, dates et pourcentages en anglais
\item Les questions en WH- (who, what, where, when, why, how many)
\item Le présent simple et le present perfect pour parler d'actions et d'expériences
\item Notions sur les ONG et la citoyenneté (Chapter 4, lessons précédentes)
\end{itemize}
\end{prerequis}

\newpage

\coursec{Warm-up: entering the theme}

\courssub{Opening situation: a radio report from Bamenda}

\begin{textebox}[Every Saturday morning in Bamenda]
Every Saturday morning in Bamenda, young volunteers join the “Girls Can Lead” association to mentor teenage girls and convince parents to keep their daughters in school. Last year, a British NGO reported that over 5,000 Cameroonian volunteers took part in gender equality campaigns across the country.

Imagine you are listening to a radio report about one of these volunteers: can you catch her name, her age, the number of girls she has helped, and the villages she visits? That is exactly the skill this lesson will train — listening for key facts in texts about volunteering for gender equality.
\end{textebox}

\textbf{Problématique.} How can we understand, from a spoken text, the essential facts about people who volunteer to promote gender equality — who they are, where they work, what they do, and why? Pour répondre à cette question, nous allons d'abord construire le vocabulaire du thème, apprendre à prédire le contenu d'un enregistrement avant de l'écouter, puis échanger nos opinions sur le volontariat. À la fin de la leçon, tu rédigeras un \cle{short appeal for volunteers} (un court appel au volontariat).

\courssub{Brainstorming: what is volunteering?}

Avant toute définition, observons des exemples concrets tirés de la vie quotidienne au Cameroun :

\begin{itemize}
\item \eng{A nurse in Yaoundé spends her Sundays teaching young mothers about hygiene — for free.} « Une infirmière à Yaoundé passe ses dimanches à enseigner l'hygiène aux jeunes mères — gratuitement. »
\item \eng{Students in Douala collect books and send them to village schools without receiving any money.} « Des élèves de Douala collectent des livres et les envoient à des écoles de village sans recevoir d'argent. »
\item \eng{In Buea, a retired teacher helps girls prepare for the GCE examination in her living room.} « À Buea, une enseignante retraitée aide des filles à préparer le GCE dans son salon. »
\end{itemize}

Que remarques-tu ? Dans les trois cas, une personne \cle{gives her time} (donne de son temps), \cle{helps a community} (aide une communauté) et \cle{is not paid} (n'est pas payée). C'est exactement le sens du mot \cle{volunteering}.

\begin{definition}[Volunteering]
\eng{Volunteering means giving your time and energy freely to help a cause or a community, without being paid.} « Le volontariat, c'est donner librement de son temps et de son énergie pour aider une cause ou une communauté, sans être payé. »
\begin{itemize}
\item \eng{a volunteer} (nom) : un volontaire, une volontaire — \eng{She is a volunteer in a health centre.} « Elle est volontaire dans un centre de santé. »
\item \eng{to volunteer} (verbe) : se porter volontaire — \eng{He volunteered to coach the girls' football team.} « Il s'est porté volontaire pour entraîner l'équipe de football des filles. »
\item \eng{voluntary} (adjectif) : volontaire, bénévole — \eng{voluntary work} = travail bénévole.
\end{itemize}
\end{definition}

\begin{definition}[Gender equality]
\eng{Gender equality is the state in which men and women enjoy the same rights, opportunities and respect.} « L'égalité des genres (ou égalité hommes-femmes) est la situation dans laquelle les hommes et les femmes jouissent des mêmes droits, des mêmes opportunités et du même respect. »
\begin{itemize}
\item \eng{gender} (nom) : le genre (social), par opposition au sexe biologique.
\item \eng{equality} (nom) : l'égalité ; adjectif : \eng{equal} (égal) ; verbe : \eng{to treat equally} (traiter de façon égale).
\item Exemple : \eng{Gender equality means that girls and boys have the same right to education.} « L'égalité des genres signifie que les filles et les garçons ont le même droit à l'éducation. »
\end{itemize}
\end{definition}

\begin{attention}[Faux amis : volunteer et voluntary]
Attention, francophones ! Le verbe anglais \eng{to volunteer} ne signifie pas seulement « faire du volontariat » : il veut dire \textbf{se porter volontaire, proposer spontanément son aide}. On peut donc dire \eng{I volunteered to carry her bag} (« je me suis proposé pour porter son sac ») sans être membre d'une association. De même, \eng{voluntary} signifie « volontaire, bénévole », mais attention à l'ordre des mots : on dit \eng{a volunteer teacher} ou \eng{a voluntary worker}, jamais \eng{a voluntarist}. Enfin, \eng{gender} ne se traduit pas par « genre » au sens grammatical : \eng{gender equality} = « l'égalité hommes-femmes », pas « l'égalité du genre ».
\end{attention}

\courssub{Listening script: Two texts on volunteering for gender equality}

\begin{textebox}[Text 1: Radio Report – “A Voice for Girls in the Far North”]
Good morning listeners. This is CRTV Radio Far North. Today we meet Amina, a twenty-four-year-old volunteer from Maroua. For two years, Amina has worked with the association “Education for All”. Every week, she visits three villages: Kousseri, Waza, and Mora. She mentors fifteen teenage girls who risk dropping out of school. Last month, her campaign “Girls Stay in School” reached five hundred parents. Amina says: “I volunteer because I want to make a difference. Gender equality starts with education.”
\end{textebox}

\begin{textebox}[Text 2: Interview – “Building Change in Bamenda”]
\eng{Interviewer:} Good afternoon, Mr. Nfor. You coordinate volunteers in Bamenda. How many volunteers do you have?
\eng{Mr. Nfor:} We have thirty active volunteers. Most are university students.
\eng{Interviewer:} What do they do?
\eng{Mr. Nfor:} They run awareness campaigns in five villages around Bamenda. They talk to parents about early marriage. They also give school supplies to girls. Last year, we helped two hundred girls return to school.
\eng{Interviewer:} Who can volunteer?
\eng{Mr. Nfor:} Anyone over eighteen with a passion for gender equality. Just call six double-five, double-three, double-two.
\end{textebox}

\courssub{Comprehension questions}

\begin{exercice}[Text 1: Radio Report – “A Voice for Girls in the Far North”]
\textbf{A. True or False? Justify with a phrase from the text.}
\begin{enumerate}
\item Amina is twenty-five years old.
\item Amina visits four villages every week.
\item She mentors fifteen teenage girls.
\item Her campaign reached five hundred parents last month.
\end{enumerate}

\textbf{B. Multiple Choice. Choose the correct answer.}
\begin{enumerate}
\item What is the name of Amina's association?
      \begin{itemize}
      \item[a)] Girls Can Lead
      \item[b)] Education for All
      \item[c)] Volunteers for Change
      \end{itemize}
\item How long has Amina been a volunteer?
      \begin{itemize}
      \item[a)] One year
      \item[b)] Two years
      \item[c)] Three years
      \end{itemize}
\item What is the name of Amina's campaign?
      \begin{itemize}
      \item[a)] Girls Can Lead
      \item[b)] Education for All
      \item[c)] Girls Stay in School
      \end{itemize}
\end{enumerate}

\textbf{C. Short Answers. Answer in complete sentences.}
\begin{enumerate}
\item Where does Amina come from?
\item Why does Amina volunteer?
\item How many girls does she mentor?
\end{enumerate}
\end{exercice}

\begin{corrige}[Exercice 1]
\textbf{A. True or False?}
\begin{enumerate}
\item \textbf{False.} \eng{“Amina, a twenty-four-year-old volunteer…”}
\item \textbf{False.} \eng{“Every week, she visits three villages: Kousseri, Waza, and Mora.”}
\item \textbf{True.} \eng{“She mentors fifteen teenage girls…”}
\item \textbf{True.} \eng{“Last month, her campaign ‘Girls Stay in School’ reached five hundred parents.”}
\end{enumerate}

\textbf{B. Multiple Choice.}
\begin{enumerate}
\item \textbf{b)} Education for All. (\eng{“works with the association ‘Education for All’”})
\item \textbf{b)} Two years. (\eng{“For two years, Amina has worked…”})
\item \textbf{c)} Girls Stay in School. (\eng{“her campaign ‘Girls Stay in School’”})
\end{enumerate}

\textbf{C. Short Answers.}
\begin{enumerate}
\item \eng{Amina comes from Maroua.} (Ou : \eng{She is from Maroua.})
\item \eng{She volunteers because she wants to make a difference.} (Ou : \eng{Because gender equality starts with education.})
\item \eng{She mentors fifteen girls.}
\end{enumerate}
\end{corrige}

\begin{exercice}[Text 2: Interview – “Building Change in Bamenda”]
\textbf{A. True or False? Justify with a phrase from the text.}
\begin{enumerate}
\item Mr. Nfor has thirty active volunteers.
\item Most volunteers are secondary school students.
\item They run campaigns in five villages around Bamenda.
\item They helped two hundred girls return to school last year.
\end{enumerate}

\textbf{B. Multiple Choice. Choose the correct answer.}
\begin{enumerate}
\item What do the volunteers talk to parents about?
      \begin{itemize}
      \item[a)] Hygiene
      \item[b)] Early marriage
      \item[c)] Football
      \end{itemize}
\item What do they give to girls?
      \begin{itemize}
      \item[a)] Money
      \item[b)] School supplies
      \item[c)] Clothes
      \end{itemize}
\item What is the minimum age to volunteer?
      \begin{itemize}
      \item[a)] Sixteen
      \item[b)] Eighteen
      \item[c)] Twenty-one
      \end{itemize}
\end{enumerate}

\textbf{C. Short Answers. Answer in complete sentences.}
\begin{enumerate}
\item Who are the volunteers? (Occupation)
\item What are the two main activities mentioned in the interview?
\item What is the phone number to call? (Write in digits)
\end{enumerate}
\end{exercice}

\begin{corrige}[Exercice 2]
\textbf{A. True or False?}
\begin{enumerate}
\item \textbf{True.} \eng{“We have thirty active volunteers.”}
\item \textbf{False.} \eng{“Most are university students.”} (Pas secondary school students).
\item \textbf{True.} \eng{“They run awareness campaigns in five villages around Bamenda.”}
\item \textbf{True.} \eng{“Last year, we helped two hundred girls return to school.”}
\end{enumerate}

\textbf{B. Multiple Choice.}
\begin{enumerate}
\item \textbf{b)} Early marriage. (\eng{“They talk to parents about early marriage.”})
\item \textbf{b)} School supplies. (\eng{“They also give school supplies to girls.”})
\item \textbf{b)} Eighteen. (\eng{“Anyone over eighteen…”})
\end{enumerate}

\textbf{C. Short Answers.}
\begin{enumerate}
\item \eng{The volunteers are university students.} (Ou : \eng{Most of them are university students.})
\item \eng{They run awareness campaigns about early marriage and give school supplies to girls.}
\item \eng{655 33 22} (Ou \eng{6 55 33 22}).
\end{enumerate}
\end{corrige}

\courssub{Thematic vocabulary: key words and expressions}

\courssub{Trigger vocabulary: the five key words}

Voici les cinq mots déclencheurs (\cle{trigger words}) qui te permettront d'entrer dans le thème et de prédire le contenu de l'enregistrement. Prononciation indiquée en lettres françaises.

\begin{center}
\begin{tabularx}{\linewidth}{|l|X|X|}
\hline
\rowcolor{popblueL} English (prononciation) & Français & Exemple en contexte \\
\hline
\eng{volunteer} « vo-lon-tir » & volontaire ; se porter volontaire & \eng{Twenty volunteers work in the association.} \\
\hline
\eng{gender equality} « djen-deur i-kwo-li-ti » & égalité hommes-femmes & \eng{They fight for gender equality in schools.} \\
\hline
\eng{campaign} « kam-pein » & campagne (de sensibilisation) & \eng{The campaign reached ten villages.} \\
\hline
\eng{community} « ko-miou-ni-ti » & communauté & \eng{The whole community supports the project.} \\
\hline
\eng{awareness} « a-wèur-ness » & sensibilisation, prise de conscience & \eng{They raise awareness about girls' education.} \\
\hline
\end{tabularx}
\end{center}

\begin{propriete}[Expressions utiles autour du thème]
Quelques collocations (mots qui vont ensemble) à mémoriser telles quelles :
\begin{itemize}
\item \eng{to raise awareness about...} = sensibiliser à... \eng{The NGO raises awareness about early marriage.}
\item \eng{to take part in a campaign} = participer à une campagne.
\item \eng{to keep girls in school} = maintenir les filles à l'école.
\item \eng{to drop out of school} = abandonner l'école, déscolariser.
\item \eng{to mentor somebody} = encadrer, parrainer quelqu'un.
\item \eng{to make a difference} = changer les choses, avoir un impact.
\end{itemize}
\end{propriete}

La carte mentale ci-dessous organise tout le champ lexical du thème. Recopie-la dans ton cahier et ajoute une branche de ton choix.

\begin{popfigure}
\begin{center}
\begin{tikzpicture}[mindmap, grow cyclic, every node/.style={concept, align=center, font=\small},
root concept/.append style={concept color=popblue, font=\small\bfseries, text=white},
level 1/.append style={level distance=3.4cm, sibling angle=90},
level 2/.append style={level distance=2.2cm, sibling angle=50, font=\footnotesize}]
\node[root concept]{VOLUNTEERING FOR GENDER EQUALITY}
  child[concept color=poporange] { node {who?}
    child { node {volunteers} }
    child { node {NGOs} }
  }
  child[concept color=popgreen] { node {where?}
    child { node {villages} }
    child { node {schools} }
  }
  child[concept color=poppurple] { node {what?}
    child { node {campaigns} }
    child { node {mentoring} }
  }
  child[concept color=poppink] { node {why?}
    child { node {equal rights} }
    child { node {girls' education} }
  };
\end{tikzpicture}
\end{center}
\legende{Carte mentale lexicale : le champ du volontariat pour l'égalité des genres.}
\end{popfigure}

\begin{savaistu}[Volunteering in Cameroon]
Au Cameroun, des journées nationales du volontariat mobilisent chaque année des milliers de jeunes, notamment pour la scolarisation des filles dans les régions du Nord et de l'Extrême-Nord. De nombreuses associations, souvent soutenues par des ONG internationales, envoient des volontaires dans les villages pour discuter avec les parents, offrir des fournitures scolaires et suivre les filles en difficulté. Le mot anglais \eng{volunteer} vient du latin \emph{voluntas}, « la volonté, le bon vouloir » : un volontaire est littéralement quelqu'un qui agit « de son plein gré ».
\end{savaistu}

\courssub{Listening strategies and worked example}

\courssub{Predicting from the title: what will the recording be about?}

Avant d'écouter un enregistrement, un bon auditeur ne reste pas passif : il \cle{prédit} le contenu. Le titre de notre leçon est : \eng{Texts on volunteering in gender equality promotion activities}. Décomposons-le :

\begin{itemize}
\item \eng{volunteering} $\rightarrow$ on va entendre parler de personnes qui aident gratuitement ;
\item \eng{gender equality promotion} $\rightarrow$ le but est de promouvoir l'égalité hommes-femmes ;
\item \eng{activities} $\rightarrow$ on attend des actions concrètes : campagnes, visites de villages, mentorat.
\end{itemize}

\begin{methode}[How to predict before listening]
\begin{enumerate}
\item \textbf{Read the title and the questions BEFORE listening.} Lis le titre et toutes les questions de compréhension avant d'appuyer sur « play » : ton cerveau saura quoi chercher.
\item \textbf{Underline the important words.} Souligne les mots porteurs de sens : noms propres, chiffres, verbes d'action (\eng{help, train, visit, build}).
\item \textbf{Imagine two or three possible answers.} Pour chaque question, imagine deux ou trois réponses possibles : si la question est \eng{How many girls has she helped?}, prépare-toi à entendre un nombre, et révise les nombres en anglais (\eng{fifteen} vs \eng{fifty} !).
\item \textbf{Anticipate the vocabulary.} Fais la liste des mots que tu t'attends à entendre : ici, probablement \eng{volunteer, village, school, parents, campaign, girls}.
\end{enumerate}
\end{methode}

\begin{exoresolu}[Prédire à partir du titre]
Énoncé. Voici le titre d'un enregistrement : \eng{A young volunteer changes lives in the Far North}. Sans écouter, réponds en anglais : (a) Who will probably speak? (b) Where does the story take place? (c) Write two words you expect to hear.
\tcblower
Solution détaillée. (a) \eng{A young volunteer will probably speak, or a journalist interviewing her.} Le titre annonce « a young volunteer », donc soit elle parle, soit un journaliste parle d'elle. (b) \eng{The story takes place in the Far North region of Cameroon.} Le titre dit « in the Far North ». (c) Mots attendus : par exemple \eng{girls, school, village, parents, education, campaign}. Toute réponse appartenant au champ lexical du volontariat et de l'éducation des filles est correcte : c'est exactement cette anticipation qui rendra l'écoute plus facile.
\end{exoresolu}

\courssub{Guided production: reuse the lexicon}

\begin{experience}[Two-minute pair discussion]
Travaille en binôme. Pose la question à ton voisin et réponds avec une raison :
\eng{Would you volunteer to promote girls' education? Why or why not?}
Utilise ces structures de réponse :
\begin{itemize}
\item \eng{Yes, I would, because every girl deserves an education.}
\item \eng{Yes, I would. I believe volunteering can change my community.}
\item \eng{I am not sure, because I do not have much free time, but I could help during the holidays.}
\item \eng{No, I would not, because... but I could give books or money instead.}
\end{itemize}
Chaque élève parle au moins trente secondes. Puis deux binômes présentent leur conclusion à la classe.
\end{experience}

\begin{attention}[Attention aux structures de réponse]
Ne traduis pas mot à mot le français !
\begin{itemize}
\item « Je suis d'accord pour être volontaire » ne se dit pas \eng{I am agree to volunteer} mais \eng{I agree to volunteer} ou \eng{I am willing to volunteer} (le verbe \eng{to agree} n'a pas besoin de \eng{be}).
\item « Pourquoi pas ? » = \eng{Why not?} (sans verbe).
\item Après \eng{would}, toujours la base verbale : \eng{I would volunteer}, jamais \eng{I would to volunteer}.
\end{itemize}
\end{attention}

\courssub{Final task: a short appeal for volunteers}

\begin{methode}[Your final task]
À la fin de la leçon, tu rédigeras un \cle{short appeal for volunteers} (environ 60 à 80 mots) pour une association imaginaire qui promeut l'égalité des genres dans ta région. Ton texte devra :
\begin{enumerate}
\item nommer l'association et sa cause ;
\item expliquer ce que font les volontaires (au moins deux activités) ;
\item dire qui peut se porter volontaire et comment contacter l'association ;
\item utiliser au moins quatre mots du lexique du jour : \eng{volunteer, gender equality, campaign, community, awareness}.
\end{enumerate}
Modèle de début : \eng{Do you want to make a difference? Join “Girls Can Lead”! Our volunteers visit villages, mentor teenage girls and raise awareness about gender equality...}
\end{methode}

Tu es maintenant prêt à entrer dans l'écoute proprement dite : garde en tête la carte mentale, les cinq mots-clés et la stratégie de prédiction — ils seront tes meilleurs outils pour capter les faits essentiels de l'enregistrement.

\coursec{Warm-up: Entering the Theme}

Avant d'aborder l'écoute, activons vos connaissances sur le bénévolat et l'égalité des sexes. Cette étape de \emph{prédiction} est cruciale : elle prépare votre cerveau à reconnaître le vocabulaire et la structure du message oral.

\begin{experience}[Brainstorming et Prédictions — Travail en binôme]
\begin{enumerate}
    \item \textbf{Word Cloud :} En 2 minutes, notez tous les mots anglais qui vous viennent à l'esprit sur le thème \eng{Volunteering} et \eng{Gender Equality}. Partagez avec la classe. (Mots attendus : \eng{volunteer, NGO, rights, school, girl, boy, rural, workshop, diploma, convince, community}).
    \item \textbf{Photo Talk :} Observez l'image projetée (une jeune femme s'adressant à un groupe de parents dans un village, banderole \eng{Equal Chances}). Répondez en anglais :
    \begin{itemize}
        \item \eng{Who do you see?}
        \item \eng{Where are they?}
        \item \eng{What is happening?}
        \item \eng{What problem is being addressed?}
    \end{itemize}
    \item \textbf{Prediction :} D'après le titre \eng{“Community Voices”} et l'image, quels \textbf{chiffres} et \textbf{noms propres} vous attendez-vous à entendre ? Notez 3 prédictions.
\end{enumerate}
\end{experience}

\begin{methode}[Comment aborder le Warm-up ?]
\begin{enumerate}
    \item Ne cherchez pas la phrase parfaite : des \textbf{mots-clés} suffisent.
    \item Utilisez le contexte camerounais : \eng{Far North, Garoua, CRTV, GCE, bendskin} (pour le contexte sonore).
    \item Partagez vos idées : la compréhension de l'oral est collaborative avant d'être individuelle.
\end{enumerate}
\end{methode}

\coursec{Listening Script (Support Text)}

Après l'activité de \emph{warm-up} qui a permis d'activer vos connaissances (vocabulaire clé, brainstorming, prédictions à partir d'images), nous abordons maintenant le cœur de la leçon : le support d'écoute authentique. La réussite en compréhension de l'oral repose sur une démarche structurée en plusieurs écoutes, chacune ayant un objectif précis. Nous découvrons d'abord le script intégral pour l'analyse linguistique, puis nous mettrons en œuvre la stratégie des trois écoutes.

\courssub{The Radio Script : “Community Voices”}

Le texte ci-dessous est le script exact de l'émission de radio que votre enseignant lira (ou diffusera). Il s'agit d'un reportage au format “magazine communautaire”, typique des radios locales anglophones (CRTV Radio, BBC Africa, RFI English). Observez la structure : accroche (\emph{hook}), présentation du témoin, actions concrètes, chiffres, témoignage direct (\emph{quote}), informations pratiques (numéro de téléphone), slogan final.

\begin{textebox}[RADIO SCRIPT — “Community Voices”]
Good morning, listeners! Today on Community Voices, we meet Amina Bello, a 24-year-old volunteer from Garoua. For the past three years, Amina has worked with the association Equal Chances, which promotes gender equality in the Far North Region. Every weekend, she travels to rural villages to talk to parents about sending their daughters to school. So far, her team of fifteen volunteers has convinced more than 200 families to enrol their girls in secondary school. “It is not always easy,” Amina says. “Some parents think education is only for boys. But when they see successful women from their own village, they change their minds.” The association also organises workshops on women's rights and trains young girls in leadership. Last month, 50 girls received certificates after completing a three-month training programme. Amina's dream? “To see every girl in my region holding a diploma.” If you want to join Equal Chances, call 677-45-32-10. Volunteering changes lives — starting with yours!
\end{textebox}

\courssub{Glossary and Key Vocabulary}

Avant d'appliquer la stratégie d'écoute, assurez-vous de maîtriser le lexique suivant, extrait du texte. Ces mots sont fréquents dans les discours sur le développement, les ONG et les droits humains. \textbf{La prononciation est notée en approximation française entre guillemets} (ex. « in-rol »).

\begin{textebox}[GLOSSARY — Vocabulaire clé du reportage]
\textbf{to enrol} (UK) / \eng{enroll} (US) « in-rol » (verbe) : inscrire (à l'école, sur une liste). \eng{They enrolled their daughter in secondary school.} « Ils ont inscrit leur fille au secondaire. » \\
\textbf{a workshop} « wor-cheup » (nom) : un atelier (de formation, de travail). \eng{The association organises workshops on women's rights.} « L'association organise des ateliers sur les droits des femmes. » \\
\textbf{a diploma} « di-plo-me » (nom) : un diplôme (souvent niveau secondaire ou supérieur court). \eng{Every girl holding a diploma.} « Chaque fille tenant un diplôme. » \\
\textbf{to convince} « keun-vinss » (verbe) : convaincre. \eng{They convinced 200 families.} « Ils ont convaincu 200 familles. » \\
\textbf{rural} « rou-ral » (adjectif) : rural, campagnard. \eng{She travels to rural villages.} « Elle se rend dans des villages ruraux. » \\
\textbf{to train} « treyn » (verbe) : former, entraîner. \eng{The programme trains young girls in leadership.} « Le programme forme les jeunes filles au leadership. » \\
\textbf{a certificate} « ser-ti-fi-ket » (nom) : une attestation, un certificat. \eng{50 girls received certificates.} « 50 filles ont reçu des attestations. » \\
\textbf{to change one's mind} (locution) : changer d'avis. \eng{They change their minds.} « Ils changent d'avis. » \\
\textbf{gender equality} (nom) : l'égalité des sexes / entre genres. \\
\textbf{volunteer} « vo-lun-ti-eu » (nom/verbe) : volontaire, bénévole / faire du bénévolat. \\
\textbf{so far} (locution) : jusqu'ici, jusqu'à présent (marque le présent perfect).
\end{textebox}

\courssub{Listening Strategy : The Three-Phase Approach}

En compréhension de l'oral (Listening), on ne comprend pas tout du premier coup. La méthode universitaire et celle préconisée par le MINESEC reposent sur \textbf{trois écoutes successives} aux objectifs distincts. Ne cherchez pas à tout écrire dès la première écoute.

\begin{methode}[Stratégie des 3 écoutes — Mode d'emploi]
\begin{enumerate}
    \item \textbf{1re écoute : Global Comprehension (Saisir le thème général).} \\
    \textit{Objectif :} Identifier le type de document (reportage, dialogue, discours), le thème principal, le ton, les intervenants principaux. \\
    \textit{Action :} Écoutez sans écrire. Répondez mentalement : \eng{Who? What? Where?} (Qui ? Quoi ? Où ?).
    \item \textbf{2e écoute : Selective Listening (Repérer les détails clés).} \\
    \textit{Objectif :} Relever les \textbf{noms propres} (Amina, Garoua, Equal Chances), les \textbf{chiffres} (24, 3, 15, 200, 50, 3, 677-45-32-10), les \textbf{lieux} (Far North Region, villages ruraux). \\
    \textit{Action :} Prenez des notes sous forme de mots-clés ou de schémas (mind-map). Ne rédigez pas de phrases complètes.
    \item \textbf{3e écoute : Checking and Completing (Vérifier et compléter).} \\
    \textit{Objectif :} Valider vos réponses aux questions de détail (Vrai/Faux, QCM, réponses courtes), combler les trous, vérifier les connecteurs logiques (\eng{but, so far, also}). \\
    \textit{Action :} Relisez vos notes, répondez aux questions écrites, puis relisez le script (si fourni) pour validation finale.
\end{enumerate}
\end{methode}

\begin{popfigure}
\begin{tikzpicture}[
    node distance=1.2cm and 2cm,
    box/.style={draw=popblue, thick, rounded corners, fill=popblueL, text width=4.5cm, align=center, minimum height=2.2cm, font=\small},
    arrow/.style={-Stealth, thick, popdark}
]
\node[box, align=center] (ec1){\textbf{1\textsuperscript{re} ÉCOUTE}\\\emph{Global Comprehension}\\\eng{Who? What? Where?}\\\textit{Thème : Volunteering \& Gender Equality}};
\node[box, right=of ec1, align=center] (ec2){\textbf{2\textsuperscript{e} ÉCOUTE}\\\emph{Selective Listening}\\\textbf{Key Details :}\\\textit{Names: Amina, Equal Chances}\\\textit{Numbers: 24, 3, 15, 200, 50, 3}\\\textit{Places: Garoua, Far North}};
\node[box, right=of ec2, align=center] (ec3){\textbf{3\textsuperscript{e} ÉCOUTE}\\\emph{Checking \& Completing}\\\textit{Validate True/False}\\\textit{Complete short answers}\\\textit{Check connectors (but, so far)}};
\draw[arrow] (ec1) -- node[above, font=\tiny, text=popdark] {Focus on details} (ec2);
\draw[arrow] (ec2) -- node[above, font=\tiny, text=popdark] {Validate answers} (ec3);
\end{tikzpicture}
\legende{Organigramme de la stratégie des trois écoutes successives}
\end{popfigure}

\courssub{Language Focus : Grammar and Vocabulary in Context}

Le script est une mine d'or pour réviser des points grammaticaux essentiels au programme de Terminale (Present Perfect, quantificateurs, discours rapporté, conditionnel zéro). Observons d'abord des phrases extraites du texte, puis formulons les règles.

\begin{exemplebox}[Analyse de phrases du script — Observation guidée]
\begin{enumerate}
    \item \eng{For the past three years, Amina \textbf{has worked} with the association Equal Chances.}\\
    « Depuis trois ans, Amina travaille / travaille avec l'association... » \\
    \textit{Observation :} \textbf{Present Perfect Simple} (\eng{has worked}) + durée (\eng{for three years}) + point de départ (\eng{since / for the past}). Action commencée dans le passé et qui continue \textbf{au présent}.
    \item \eng{Her team \textbf{has convinced} more than 200 families to enrol their girls.}\\
    « Son équipe a convaincu plus de 200 familles... » \\
    \textit{Observation :} \textbf{Present Perfect Simple} pour un résultat acquis, une expérience achevée \textbf{à ce jour} (\eng{so far}).
    \item \eng{“Some parents \textbf{think} education \textbf{is} only for boys.”}\\
    « Certains parents pensent que l'éducation est seulement pour les garçons. » \\
    \textit{Observation :} \textbf{Present Simple} pour des opinions, des vérités générales (dans l'esprit des parents).
    \item \eng{“But when they \textbf{see} successful women..., they \textbf{change} their minds.”}\\
    « Mais quand ils voient des femmes qui réussissent..., ils changent d'avis. » \\
    \textit{Observation :} \textbf{Zero Conditional} (si présent $\rightarrow$ présent) : vérité générale, cause-effet systématique. \eng{When} = \eng{If} ici.
    \item \eng{Last month, 50 girls \textbf{received} certificates after \textbf{completing} a three-month training programme.}\\
    « Le mois dernier, 50 filles ont reçu des certificats après avoir terminé... » \\
    \textit{Observation :} \textbf{Past Simple} (\eng{received}) pour une action datée (\eng{Last month}). \textbf{Gérondif présent} (\eng{after completing}) avec valeur d'antériorité (après avoir terminé).
\end{enumerate}
\end{exemplebox}

\begin{propriete}[Rappel : Present Perfect Simple vs Past Simple — Le piège francophone]
\begin{center}
\begin{tabularx}{\linewidth}{|X|X|X|}
\hline
\rowcolor{popblueL}
\textbf{Critère} & \textbf{Present Perfect Simple} & \textbf{Past Simple} \\
\hline
Forme & \eng{has/have + V-ed / 3\textsuperscript{e} col.} & \eng{V-ed / 2\textsuperscript{e} col.} \\
\hline
Lien avec le présent & \textbf{OUI} : action inachevée, résultat visible \textbf{maintenant}. & \textbf{NON} : action terminée, datée dans le passé. \\
\hline
Marqueurs temporels & \eng{for, since, just, already, yet, so far, ever, never, recently.} & \eng{yesterday, last week/month/year, in 2020, ago, when?} \\
\hline
\textbf{Exemple du texte} & \eng{Amina \textbf{has worked}... for three years.} (Elle y est encore) & \eng{50 girls \textbf{received} certificates \textbf{last month}.} (Fini) \\
\hline
\textbf{Erreur fréquente} & \eng{*Amina works here for three years.} (Présent simple + for = FAUX) & \eng{*Amina has worked here last year.} (Present Perfect + date passée = FAUX) \\
\hline
\end{tabularx}
\end{center}
\end{propriete}

\begin{propriete}[Quantificateurs du texte : \eng{more than}, \eng{some}, \eng{every}, \eng{fifteen}]
\begin{center}
\begin{tabular}{|l|l|l|}
\hline
\rowcolor{popgreenL}
\textbf{Quantificateur} & \textbf{Exemple du texte} & \textbf{Traduction / Nuance} \\
\hline
\eng{more than} + nombre & \eng{more than 200 families} & \textbf{Plus de} 200 familles (supérieur strict). \\
\hline
\eng{some} + pluriel & \eng{Some parents think...} & \textbf{Certains} parents (indéfini, partie d'un tout). \\
\hline
\eng{every} + singulier & \eng{every girl in my region} & \textbf{Chaque} fille / \textbf{Toutes les} filles (distributif, verbe au singulier). \\
\hline
\textbf{Nombre cardinal} & \eng{fifteen volunteers} & \textbf{Quinze} bénévoles (précis, adjectif invariable). \\
\hline
\end{tabular}
\end{center}
\end{propriete}

\courssub{Focus on Pronunciation and Common Pitfalls}

\begin{attention}[Faux-amis et pièges de prononciation — Spécial Francophones]
\begin{itemize}
    \item \textbf{Volunteer (nom/verbe) « vo-lun-ti-eu » vs Voluntary (adjectif) « vo-lun-te-ri »} \\
    \textit{Erreur :} “He is a voluntary.” (FAUX). \\
    \textit{Correct :} \eng{He is a \textbf{volunteer}.} (Nom) / \eng{He does \textbf{voluntary} work.} (Adjectif). \\
    \textit{Astuce :} \eng{Volunteer} rime avec \eng{engineer} (son final \eng{/ɪə/} $\approx$ « i-eu »). \eng{Voluntary} a 4 syllabes distinctes.
    \item \textbf{Actually « ak-chou-eu-li » $\neq$ Actuellement} \\
    \textit{Faux ami majeur.} \eng{Actually} = \textbf{En fait, en réalité, vraiment}. \\
    \eng{“Actually, I don't agree.”} = « En fait, je ne suis pas d'accord. » \\
    Pour dire \textbf{actuellement} : \eng{currently}, \eng{at the moment}, \eng{nowadays}.
    \item \textbf{Diploma vs Diplôme} \\
    En anglais, \eng{a diploma} désigne souvent un diplôme technique, universitaire court (BTS, DUT, Bachelor) ou de fin d'études secondaires. Pour le \eng{Baccalauréat}, on dit \eng{High School Diploma} (US) ou \eng{A-levels} (UK). Au Cameroun anglophone : \eng{GCE Advanced Level}.
    \item \textbf{Prononciation de \eng{rural} « rou-ral »} : Deux syllabes seulement. Le \eng{u} est bref (comme dans \eng{put} « pout »). Ne dites pas *\eng{ru-ral} (3 syllabes).
    \item \textbf{\eng{Enrol} (UK) / \eng{Enroll} (US)} : Double \eng{l} aux USA, simple \eng{l} au UK/Commonwealth (Cameroun inclus). Participe passé : \eng{enrolled} (un seul \eng{l} partout).
\end{itemize}
\end{attention}

\courssub{Comprehension Questions (Practice)}

Appliquez maintenant la stratégie des 3 écoutes. Votre enseignant lira le texte (ou le diffusera) trois fois. Répondez aux questions ci-dessous \textbf{après la 2e et la 3e écoute}.

\begin{exercice}[Compréhension détaillée — Types Baccalauréat]
\textbf{A. True / False + Justification} (Citez le segment exact du texte, 3 à 6 mots maximum).
\begin{enumerate}
    \item Amina Bello works for an association based in Yaoundé.
    \item The association has convinced exactly 200 families.
    \item The training programme mentioned lasted three months.
    \item Amina works with a team of fifty volunteers.
    \item The phone number to join the association is 677-45-32-10.
\end{enumerate}

\textbf{B. Multiple Choice Questions.} Choisissez la bonne réponse (a, b ou c).
\begin{enumerate}
    \item What is the main goal of the association “Equal Chances”?
    \begin{itemize}
        \item[a)] To build schools in Garoua.
        \item[b)] To promote gender equality in the Far North Region.
        \item[c)] To train boys in leadership.
    \end{itemize}
    \item How often does Amina travel to villages?
    \begin{itemize}
        \item[a)] Every day.
        \item[b)] Every week.
        \item[c)] Every weekend.
    \end{itemize}
    \item What does Amina say makes parents change their minds?
    \begin{itemize}
        \item[a)] Seeing successful women from their own village.
        \item[b)] Receiving money from the association.
        \item[c)] Listening to the radio programme.
    \end{itemize}
\end{enumerate}

\textbf{C. Short Answers.} Répondez par une phrase courte et complète.
\begin{enumerate}
    \item How old is Amina Bello?
    \item Where does Amina come from?
    \item How many girls received certificates last month?
    \item What is Amina's dream?
\end{enumerate}
\end{exercice}

\begin{corrige}[Exercice — Comprehension Questions]
\textbf{A. True / False + Justification}
\begin{enumerate}
    \item \textbf{FALSE.} Justification : \eng{“...volunteer from \textbf{Garoua}.”} (Pas Yaoundé).
    \item \textbf{FALSE.} Justification : \eng{“...convinced \textbf{more than 200 families}...”} (Pas exactement 200).
    \item \textbf{TRUE.} Justification : \eng{“...a \textbf{three-month} training programme.”}
    \item \textbf{FALSE.} Justification : \eng{“...her \textbf{team of fifteen volunteers}...”} (15, pas 50).
    \item \textbf{TRUE.} Justification : \eng{“...call \textbf{677-45-32-10}.”}
\end{enumerate}

\textbf{B. Multiple Choice}
\begin{enumerate}
    \item \textbf{b)} (L. 3 : \eng{promotes gender equality in the Far North Region}).
    \item \textbf{c)} (L. 4 : \eng{Every \textbf{weekend}}).
    \item \textbf{a)} (L. 6-7 : \eng{when they \textbf{see successful women from their own village}}).
\end{enumerate}

\textbf{C. Short Answers}
\begin{enumerate}
    \item She is 24 years old. / Amina Bello is 24.
    \item She comes from Garoua. / She is from Garoua.
    \item 50 girls received certificates.
    \item Her dream is to see every girl in her region holding a diploma.
\end{enumerate}
\end{corrige}

\courssub{Worked Example : Detailed Comprehension Question (Exam Technique)}

Voici comment traiter une question de détail type examen (Baccalauréat / Probatoire) en utilisant la 2e et 3e écoute, avec la méthode du \emph{scanning} (repérage visuel dans le script).

\begin{exoresolu}[Question de détail — Vrai ou Faux + Justification Modèle]
\textbf{Énoncé :} \textit{Listen to the text and say whether the following statements are True (T) or False (F). Justify each answer with a specific phrase from the text.}
\begin{enumerate}
    \item Amina Bello is 25 years old.
    \item The association “Equal Chances” operates in the Adamawa Region.
    \item Amina works alone in the villages.
    \item More than 200 families have been convinced to send their girls to school.
    \item The training programme lasted two months.
\end{enumerate}

\tcblower
\textbf{Solution détaillée (Méthode : Repérage de mots-clés / \emph{Scanning})}

\begin{enumerate}
    \item \textbf{FALSE.} \\
    \textit{Justification :} \eng{“...we meet Amina Bello, a \textbf{24-year-old} volunteer...”} (L. 2). Le texte dit 24 ans, pas 25.
    \item \textbf{FALSE.} \\
    \textit{Justification :} \eng{“...promotes gender equality in the \textbf{Far North Region}.”} (L. 3). C'est l'Extrême-Nord, pas l'Adamaoua.
    \item \textbf{FALSE.} \\
    \textit{Justification :} \eng{“...her \textbf{team of fifteen volunteers} has convinced...”} (L. 4-5). Elle travaille en équipe (15 bénévoles), pas seule.
    \item \textbf{TRUE.} \\
    \textit{Justification :} \eng{“...convinced \textbf{more than 200 families} to enrol their girls...”} (L. 5). Correspondance exacte.
    \item \textbf{FALSE.} \\
    \textit{Justification :} \eng{“...after completing a \textbf{three-month} training programme.”} (L. 9). Durée : 3 mois, pas 2.
\end{enumerate}

\textbf{Conseil de pro :} Pour la justification, copiez \textbf{exactement} le segment de phrase (3 à 6 mots max) qui contient la réponse. Ne paraphrasez pas en anglais approximatif.
\end{exoresolu}

\courssub{Guided Practice : Reusing the Lexicon}

Avant de passer à l'expression écrite libre, entraînez-vous à réemployer le vocabulaire et les structures dans des phrases personnelles. C'est l'étape d'\emph{appropriation} (transfer).

\begin{exercice}[Mise en œuvre lexicale et grammaticale]
Complétez les phrases suivantes en utilisant le vocabulaire du glossaire et la bonne forme verbale (Present Perfect, Past Simple, Present Simple, Zero Conditional).
\begin{enumerate}
    \item The NGO \_\_\_\_\_\_\_\_ (organise) free \_\_\_\_\_\_\_\_ (workshop) on leadership every summer.
    \item Since 2021, they \_\_\_\_\_\_\_\_ (convince) over 500 parents to \_\_\_\_\_\_\_\_ (enrol) their daughters.
    \item Last year, 30 girls \_\_\_\_\_\_\_\_ (receive) a \_\_\_\_\_\_\_\_ (certificate) after the training.
    \item \_\_\_\_\_\_\_\_ (Every) child has the right to education, but in \_\_\_\_\_\_\_\_ (rural) areas, access is difficult.
    \item My brother is a \_\_\_\_\_\_\_\_ (volunteer); he does \_\_\_\_\_\_\_\_ (voluntary) work for the Red Cross.
    \item \_\_\_\_\_\_\_\_ (Actually), the situation \_\_\_\_\_\_\_\_ (improve) slowly but surely.
    \item When parents \_\_\_\_\_\_\_\_ (see) the benefits, they usually \_\_\_\_\_\_\_\_ (support) their daughters.
\end{enumerate}
\end{exercice}

\begin{corrige}[Exercice n°1 — Mise en œuvre]
\begin{enumerate}
    \item \textbf{organises} (Present Simple, habitude) / \textbf{workshops} (pluriel après \eng{free}).
    \item \textbf{have convinced} (Present Perfect, \eng{Since 2021} $\rightarrow$ durée non achevée) / \textbf{enrol} (infinitif après \eng{to}).
    \item \textbf{received} (Past Simple, \eng{Last year} = date passée) / \textbf{certificates} (pluriel naturel : chaque fille reçoit son certificat) ou \textbf{a certificate} (singulier après \eng{a}).
    \item \textbf{Every} (singulier : \eng{Every child has}) / \textbf{rural} (adjectif invariable).
    \item \textbf{volunteer} (nom) / \textbf{voluntary} (adjectif). \textit{Attention à la distinction nom/adjectif !}
    \item \textbf{Actually} (En fait) / \textbf{has improved} (Present Perfect, résultat présent) ou \textbf{is improving} (Present Continuous, évolution en cours).
    \item \textbf{see} / \textbf{support} (Zero Conditional : Present Simple dans les deux propositions pour vérité générale).
\end{enumerate}
\end{corrige}

\begin{savaistu}[Le bénévolat au Cameroun : Contexte réel]
Le Cameroun compte des milliers d'associations de jeunesse et d'ONG locales (ex. \eng{RECAJ}, \eng{AJVC}, \eng{Plan International Cameroon}, \eng{UNICEF} partners) qui mènent des actions de \eng{gender mainstreaming}. Le \textbf{Ministère de la Promotion de la Femme et de la Famille (MINPROFF)} pilote la politique nationale. Le volontariat est souvent valorisé dans le \eng{National Youth Service} (Service National de la Jeunesse) et les programmes \eng{PAJER-U} (Programme d'Appui aux Jeunes Ruraux et Urbains). Connaître ce vocabulaire vous permet de comprendre les reportages de CRTV Radio, BBC Africa \eng{Focus on Africa} ou RFI \eng{Appels sur l'actualité} sur ce thème.
\end{savaistu}

\begin{aretenir}[Bilan de la Leçon 37 — Ce qu'il faut retenir pour l'examen]
\begin{itemize}
    \item \textbf{Structure du reportage radio :} Accroche $\rightarrow$ Témoin + Contexte $\rightarrow$ Actions/Chiffres $\rightarrow$ Citation directe $\rightarrow$ Infos pratiques $\rightarrow$ Slogan.
    \item \textbf{Stratégie 3 écoutes :} 1. Global (Thème) $\rightarrow$ 2. Détails (Noms, Chiffres, Lieux) $\rightarrow$ 3. Validation (Questions).
    \item \textbf{Grammaire clé :} \eng{Present Perfect} (\eng{for/since/so far}) vs \eng{Past Simple} (\eng{last month/ago}). \eng{Zero Conditional} (\eng{When/If} + Present $\rightarrow$ Present). Gérondif après préposition (\eng{after completing}).
    \item \textbf{Lexique indispensable :} \eng{volunteer} (n/v) vs \eng{voluntary} (adj) ; \eng{enrol} (UK) ; \eng{workshop} ; \eng{diploma} ; \eng{convince} ; \eng{rural} ; \eng{actually} (= en fait) ; \eng{certificate}.
    \item \textbf{Prononciation (approx. française) :} \eng{rural} « rou-ral » ; \eng{volunteer} « vo-lun-ti-eu » ; \eng{certificate} « ser-ti-fi-ket » ; \eng{workshop} « wor-cheup ».
    \item \textbf{Compétence transférable :} Savoir remplir une grille de repérage (Who/What/Where/When/How many) en 2e écoute est la clé de la réussite au Bac.
\end{itemize}
\end{aretenir}

\coursec{Comprehension questions}

Après avoir écouté (ou lu) le script de la section précédente — l'interview d'Amina Bello, volontaire de l'association \eng{Equal Chances} dans la Région de l'Extrême-Nord — il est temps de vérifier ta compréhension. Au baccalauréat comme dans les évaluations de classe, trois types de questions reviennent toujours après un document audio : les \cle{vrai/faux justifiés}, les \cle{QCM} et les \cle{questions à réponses courtes}. Cette section t'entraîne aux trois, avec la méthode pour obtenir tous les points.

\courssub{Before you answer: the note-taking grid}

Pendant l'écoute, tu ne peux pas tout retenir. La stratégie des bons candidats consiste à remplir une grille de prise de notes avec les \cle{noms propres}, les \cle{chiffres} et les \cle{lieux}. Voici la grille vierge correspondant à notre enregistrement ; remplis-la pendant que le professeur lit le script une seconde fois.

\begin{popfigure}
\begin{tikzpicture}
\draw[thick, popblue] (0,0) rectangle (14.5,4.2);
\draw[thick, popblue] (0,3.4) -- (14.5,3.4);
\fill[popblueL] (0,3.4) rectangle (14.5,4.2);
\node[align=center, font=\bfseries] at (1.5,3.8) {Name};
\node[align=center, font=\bfseries] at (3.4,3.8) {Age};
\node[align=center, font=\bfseries] at (5.4,3.8) {Place};
\node[align=center, font=\bfseries] at (8.1,3.8) {Number of\\volunteers};
\node[align=center, font=\bfseries] at (11.1,3.8) {Families\\convinced};
\node[align=center, font=\bfseries] at (13.6,3.8) {Girls\\trained};
\draw[popblue] (2.6,0) -- (2.6,3.4);
\draw[popblue] (4.2,0) -- (4.2,3.4);
\draw[popblue] (6.6,0) -- (6.6,3.4);
\draw[popblue] (9.6,0) -- (9.6,3.4);
\draw[popblue] (12.5,0) -- (12.5,3.4);
\node[align=center, popmuted] at (7.2,1.7) {\emph{Listen and complete...}};
\end{tikzpicture}
\legende{Grille de prise de notes à compléter pendant l'écoute : les six informations-clés du script.}
\end{popfigure}

\begin{methode}[How to fill the grid while listening]
\begin{enumerate}
\item \textbf{Avant l'écoute} : lis les titres des colonnes et prédis le type d'information attendu (un nombre ? un lieu ? un nom ?).
\item \textbf{Première écoute} : note uniquement les chiffres et les noms propres, en abrégé (\eng{32}, \eng{Maroua}, \eng{15 vol.}).
\item \textbf{Deuxième écoute} : vérifie chaque case et complète les trous.
\item \textbf{Après l'écoute} : utilise la grille pour répondre aux questions sans réécouter.
\end{enumerate}
\end{methode}

\courssub{Thematic vocabulary: Volunteering \& Gender Equality}

Avant de traiter les exercices, assure-toi de maîtriser le lexique clé du script. Ces mots reviennent dans les questions et dans ta production.

\begin{definition}[Mots et expressions essentiels du script]
\begin{center}
\begin{tabularx}{\linewidth}{|l|X|l|}
\hline
\rowcolor{popblueL} \textbf{English} & \textbf{Français} & \textbf{Classe} \\
\hline
\eng{volunteer} & bénévole, volontaire & n. / v. \\
\hline
\eng{gender equality} & égalité des sexes / genre & n. phr. \\
\hline
\eng{promote} & promouvoir, encourager & v. \\
\hline
\eng{Far North Region} & Région de l'Extrême-Nord & n. phr. \\
\hline
\eng{village} & village & n. \\
\hline
\eng{convince} & convaincre & v. \\
\hline
\eng{workshop} & atelier (de formation) & n. \\
\hline
\eng{training programme} & programme de formation & n. phr. \\
\hline
\eng{diploma} & diplôme & n. \\
\hline
\eng{commitment} & engagement & n. \\
\hline
\eng{empowerment} & autonomisation, renforcement des capacités & n. \\
\hline
\eng{deserve} & mériter & v. \\
\hline
\eng{rights} & droits & n. pl. \\
\hline
\end{tabularx}
\end{center}
\emph{Note :} \eng{volunteer} peut être nom ou verbe. \eng{Gender equality} est un terme institutionnel fixe. \eng{Programme} s'écrit avec \eng{-mme} en anglais britannique (contexte camerounais).
\end{definition}

\courssub{Exercise 1: True or False, with justification}

\begin{textebox}[VRAI OU FAUX : justifiez par une citation du script.]
\textbf{1.} Amina is 34 years old.

\textbf{2.} She works in the Far North Region.

\textbf{3.} The association has fifty volunteers.

\textbf{4.} Some parents believe education is only for boys.

\textbf{5.} Amina wants every girl to have a diploma.
\end{textebox}

\begin{methode}[Répondre aux vrai/faux au Bac]
\begin{enumerate}
\item Lis l'affirmation et souligne le mot qui peut être faux (un chiffre, un lieu, un adverbe).
\item Retrouve dans le script la phrase qui parle du même point.
\item Écris \eng{TRUE} ou \eng{FALSE}, puis \textbf{cite toujours la phrase du texte} qui justifie ta réponse : une réponse sans justification perd des points, même si elle est correcte.
\item Formule-type : \eng{FALSE. The speaker says: “I am thirty-two years old.”}
\end{enumerate}
\end{methode}

\begin{exoresolu}[True or False — items 1 and 3 solved]
\textbf{Énoncé.} Justify your answers with a quotation. 1. Amina is 34 years old. 3. The association has fifty volunteers.
\tcblower
\textbf{Solution détaillée.}

\textbf{1. FALSE.} The speaker says: \eng{“My name is Amina Bello and I am thirty-two years old.”} Le chiffre 34 est un \cle{distracteur} : il ressemble à 32, mais le script dit bien \eng{thirty-two}. Attention à la confusion fréquente des francophones entre \eng{thirty-two} (32) et \eng{thirty-four} (34) à l'oreille.

\textbf{3. FALSE.} The speaker says: \eng{“We are only fifteen volunteers, but we work very hard.”} Le nombre \eng{fifty} (50) apparaît dans le script, mais il désigne les \emph{filles formées} (\eng{fifty girls}), pas les volontaires. C'est le piège classique du QCM et du vrai/faux : le bon chiffre au mauvais endroit.
\end{exoresolu}

\begin{corrige}[Exercice 1 — True or False]
\textbf{1. FALSE.} \eng{“I am thirty-two years old.”}

\textbf{2. TRUE.} \eng{“I work in the Far North Region, in the villages around Maroua.”}

\textbf{3. FALSE.} \eng{“We are only fifteen volunteers.”} (Fifty is the number of trained girls.)

\textbf{4. TRUE.} \eng{“Some parents still believe that education is only for boys.”}

\textbf{5. TRUE.} \eng{“My dream is that every girl in my region goes to school and gets a diploma.”}
\end{corrige}

\courssub{Exercise 2: Multiple choice — numbers and proper names}

\begin{textebox}[QCM : choisissez la bonne réponse.]
\textbf{1.} Equal Chances promotes \quad a) sports \quad b) gender equality \quad c) agriculture \quad d) tourism.

\textbf{2.} Amina travels to villages \quad a) every day \quad b) every weekend \quad c) once a year \quad d) never.

\textbf{3.} The training programme lasted \quad a) three weeks \quad b) three months \quad c) one year \quad d) two days.

\textbf{4.} To join the association, you can \quad a) write a letter \quad b) send an email \quad c) call a phone number \quad d) visit an office.
\end{textebox}

\begin{attention}[Les distracteurs reprennent les chiffres du texte]
Dans les QCM d'écoute, les mauvaises réponses (\cle{distracteurs}) contiennent presque toujours des mots ou des nombres \emph{réellement présents} dans le script, mais associés à une autre information. Ici, le script contient \eng{fifteen volunteers}, \eng{fifty girls} et \eng{two hundred families} : si tu entends un chiffre, demande-toi \textbf{à quoi il se réfère} avant de répondre. Écoute le \cle{contexte}, pas seulement le nombre. De même, \eng{three weeks} et \eng{three months} se ressemblent : la durée exacte du programme de formation est \eng{three months}.
\end{attention}

\begin{corrige}[Exercice 2 — Multiple choice]
\textbf{1. b)} gender equality — \eng{“Equal Chances is an association that promotes gender equality.”}

\textbf{2. b)} every weekend — \eng{“Every weekend, I travel to the villages to talk with families.”}

\textbf{3. b)} three months — \eng{“Our training programme lasted three months.”}

\textbf{4. c)} call a phone number — \eng{“If you want to join us, just call our phone number.”}
\end{corrige}

\courssub{Exercise 3: Short answers with WH- questions}

Les questions en \cle{WH-} (\eng{who}, \eng{where}, \eng{how many}, \eng{why}, \eng{what}) exigent une réponse courte mais complète, en phrase si possible. Rappel des correspondances :

\begin{center}
\begin{tabularx}{\linewidth}{|l|X|X|}
\hline
\rowcolor{popblueL} \textbf{Mot interrogatif} & \textbf{Question porte sur...} & \textbf{Type de réponse attendue} \\
\hline
\eng{Who} & une personne & a name, a person \\
\hline
\eng{Where} & un lieu & a place (in Maroua, in the Far North) \\
\hline
\eng{How many} & une quantité & a number (fifteen, two hundred) \\
\hline
\eng{How long} & une durée & for + duration (for six years) \\
\hline
\eng{Why} & une cause & because + clause \\
\hline
\eng{What} & une chose, une activité & a noun or an activity \\
\hline
\end{tabularx}
\end{center}

\begin{textebox}[QUESTIONS COURTES : répondez par des phrases complètes.]
\textbf{1.} What is the volunteer's full name?

\textbf{2.} How long has she worked with Equal Chances?

\textbf{3.} How many families were convinced?

\textbf{4.} What do the workshops focus on?

\textbf{5.} What is Amina's dream?
\end{textebox}

\begin{corrige}[Exercice 3 — Short answers]
\textbf{1.} Her full name is Amina Bello.

\textbf{2.} She has worked with Equal Chances for six years. (Present perfect + \eng{for} : action commencée dans le passé et qui continue.)

\textbf{3.} Two hundred families were convinced to send their daughters to school.

\textbf{4.} The workshops focus on girls' education and women's rights.

\textbf{5.} Her dream is that every girl in her region goes to school and gets a diploma.
\end{corrige}

\begin{attention}[Réponses courtes : les erreurs des francophones]
\begin{itemize}
\item Ne réponds pas seulement par un mot isolé si la consigne demande une phrase : écris \eng{Two hundred families were convinced.} et non \eng{200}.
\item \eng{How long} demande une \textbf{durée} (\eng{for six years}), pas une date. Ne confonds pas avec \eng{when}.
\item Ordre des mots dans la réponse : sujet + verbe, jamais l'inversion de la question. \eng{Her dream is...} et non \eng{Is her dream...}.
\item \eng{How many} est suivi d'un nom \textbf{pluriel} : \eng{How many families?} (faux ami : « combien de » ne prend pas de pluriel visible en français parlé).
\end{itemize}
\end{attention}

\courssub{Synthesis question: sum up Amina's message}

La dernière question, souvent la mieux notée, demande de \cle{résumer} le message principal en deux phrases. On ne recopie pas le script : on reformule avec ses propres mots en gardant l'idée essentielle (qui ? fait quoi ? pourquoi ?).

\begin{exoresolu}[Sum up Amina's message in two sentences]
\textbf{Énoncé.} In two sentences, sum up Amina's message about volunteering and gender equality.
\tcblower
\textbf{Solution détaillée.} Un bon résumé contient : (1) l'identité et l'action de la personne ; (2) le but ou le résultat.

\eng{Amina Bello is a volunteer who works with the association Equal Chances to promote girls' education in the Far North Region of Cameroon. Thanks to her commitment, many families have accepted to send their daughters to school, because she believes that every girl deserves a diploma.}

Deux phrases, aucune citation mot à mot, toutes les informations essentielles : c'est le modèle à reproduire.
\end{exoresolu}

\courssub{Collective correction: justifying orally}

La correction collective est un moment d'expression orale : chaque élève donne sa réponse \textbf{et sa preuve}. Utilise la formule rituelle :

\eng{I know this because the speaker says: “...”} — « Je le sais parce que la personne qui parle dit : … »

Variantes utiles pour enrichir ton oral :

\begin{center}
\begin{tabularx}{\linewidth}{|X|X|}
\hline
\rowcolor{popblueL} \textbf{English} & \textbf{Français} \\
\hline
\eng{I know this because the speaker says...} & Je le sais parce que le locuteur dit... \\
\hline
\eng{According to Amina, ...} & Selon Amina, ... \\
\hline
\eng{The text mentions that...} & Le texte mentionne que... \\
\hline
\eng{We can hear that...} & On peut entendre que... \\
\hline
\eng{She explains that...} & Elle explique que... \\
\hline
\end{tabularx}
\end{center}

\begin{experience}[Oral correction in pairs]
Travaillez par deux. L'élève A joue le rôle de l'examinateur et pose une question du vrai/faux ou du QCM ; l'élève B répond en phrase complète et justifie avec la formule \eng{I know this because the speaker says...}. Puis échangez les rôles. L'examinateur vérifie que la citation correspond bien au script et que le chiffre cité se rapporte à la bonne information.
\end{experience}

\begin{aretenir}[L'essentiel de la section]
\begin{itemize}
\item Vrai/faux : réponse + \textbf{citation obligatoire} du script, sinon perte de points.
\item QCM : méfie-toi des distracteurs qui reprennent les chiffres du texte (\eng{15 volunteers}, \eng{50 girls}, \eng{200 families}) — vérifie toujours à quoi le nombre se réfère.
\item Questions WH- : réponds par une phrase complète, avec le bon mot interrogatif (\eng{how long} = durée, \eng{how many} = quantité).
\item Synthèse : deux phrases, tes propres mots, l'idée essentielle.
\item À l'oral : \eng{I know this because the speaker says...}
\end{itemize}
\end{aretenir}

\coursec{Thematic vocabulary}

Après avoir travaillé la compréhension détaillée des textes d'écoute sur le volontariat et la promotion de l'égalité des sexes, nous allons maintenant consolider le lexique thématique indispensable pour comprendre, discuter et écrire sur ce sujet. Cette section structure le vocabulaire en quatre champs sémantiques majeurs, présente les collocations clés, analyse les familles de mots et attire votre attention sur les pièges classiques pour les francophones.

\courssub{Champ lexical 1 : Le volontariat et l'engagement associatif}

Observons d'abord quelques phrases extraites des scripts d'écoute de la leçon précédente pour mettre en contexte les termes fondamentaux.

\begin{textebox}[Extraits du script d'écoute — Volunteering context]
Many young people in Buea decide to \textbf{volunteer} for local \textbf{NGOs} during the long holidays. 
\textit{Beaucoup de jeunes à Buea décident de faire du bénévolat pour des ONG locales pendant les grandes vacances.}

\textbf{Joining} an \textbf{association} requires a strong sense of \textbf{commitment}. 
\textit{Rejoindre une association exige un fort sens de l'engagement.}

\textbf{Unpaid work} does not mean useless work; volunteers gain valuable \textbf{skills} and \textbf{experience}. 
\textit{Le travail non rémunéré ne signifie pas un travail inutile ; les bénévoles acquièrent des compétences et une expérience précieuses.}
\end{textebox}

\begin{definition}[Termes de base du volontariat]
\begin{itemize}
    \item \cle{volunteer} (n.) : un volontaire, un bénévole. Personne qui offre librement son temps et ses compétences sans rémunération.
    \item \cle{to volunteer} (v.) : faire du bénévolat, se porter volontaire. Verbe régulier : \eng{I volunteered last summer.} (J'ai fait du bénévolat l'été dernier.)
    \item \cle{NGO} (n., abbr. de \eng{Non-Governmental Organisation}) : une ONG (Organisation Non Gouvernementale). Structure privée, à but non lucratif, indépendante de l'État.
    \item \cle{association} (n.) : une association. Groupement de personnes unies pour un but commun (sportif, culturel, humanitaire).
    \item \cle{to join} (v.) : rejoindre, adhérer à. \eng{She joined the Red Cross.} (Elle a rejoint la Croix-Rouge.)
    \item \cle{commitment} (n.) : l'engagement (moral, contractuel), la dévotion. Nom issu du verbe \eng{to commit} (s'engager).
    \item \cle{unpaid work} / \cle{voluntary work} (n.) : le travail non rémunéré / le travail bénévole. Attention : \eng{voluntary} (adj.) = volontaire, bénévole ; \eng{voluntarily} (adv.) = volontairement.
\end{itemize}
\end{definition}

\begin{textebox}[TABLEAU 1 — Volunteering : Vocabulaire essentiel]
\begin{center}
\begin{tabularx}{\linewidth}{|X|X|X|}
\hline
\rowcolor{popblueL}
\textbf{English} & \textbf{Français} & \textbf{Note / Collocation} \\
\hline
a volunteer & un volontaire / un bénévole & \eng{a community volunteer} (un bénévole communautaire) \\
to volunteer (for / to do sth) & faire du bénévolat / se porter volontaire & \eng{to volunteer for an NGO} \\
an NGO (Non-Governmental Organisation) & une ONG & \eng{an international NGO} \\
an association & une association & \eng{a youth association} \\
to join & rejoindre / adhérer à & \eng{to join a campaign} \\
a commitment & un engagement & \eng{to make a commitment to} \\
unpaid work / voluntary work & travail non rémunéré / bénévolat & \eng{voluntary service} (service volontaire) \\
a volunteer coordinator & un coordinateur des bénévoles & poste clé dans une ONG \\
to sign up (for) & s'inscrire (à) & plus informel que \eng{to join} \\
\hline
\end{tabularx}
\end{center}
\end{textebox}

\begin{propriete}[Grammaire : \eng{Volunteer} — Famille de mots]
Le mot \cle{volunteer} fonctionne comme nom et verbe. Ses dérivés suivent la morphologie standard :
\begin{center}
\begin{tabular}{|l|l|l|}
\hline
\rowcolor{popgreenL}
\textbf{Forme} & \textbf{Catégorie} & \textbf{Exemple traduit} \\
\hline
volunteer & nom (n.) / verbe (v.) & \eng{She is a volunteer.} / \eng{He volunteers weekly.} \\
voluntary & adjectif (adj.) & \eng{It was a voluntary decision.} (C'était une décision volontaire.) \\
voluntarily & adverbe (adv.) & \eng{He left voluntarily.} (Il est parti volontairement.) \\
\hline
\end{tabular}
\end{center}
\emph{Attention :} En français, « volontaire » est adjectif ET nom. En anglais, le nom est \eng{volunteer} (pas *voluntary).
\end{propriete}

\courssub{Champ lexical 2 : L'égalité des sexes et les droits des femmes}

Les textes d'écoute utilisent un vocabulaire précis pour décrire les enjeux, les obstacles et les objectifs. Voici les termes incontournables.

\begin{textebox}[Extraits du script d'écoute — Gender Equality context]
Promoting \textbf{gender equality} is not only a \textbf{women's rights} issue; it benefits the whole society. 
\textit{Promouvoir l'égalité des sexes n'est pas seulement une question de droits des femmes ; cela profite à toute la société.}

\textbf{Discrimination} based on sex prevents girls from accessing \textbf{equal opportunities} in education and employment. 
\textit{La discrimination basée sur le sexe empêche les filles d'accéder à l'égalité des chances dans l'éducation et l'emploi.}

\textbf{Empowerment} allows women to become \textbf{role models} for the next generation. 
\textit{L'autonomisation permet aux femmes de devenir des modèles pour la génération suivante.}
\end{textebox}

\begin{textebox}[TABLEAU 2 — Gender Equality : Vocabulaire essentiel]
\begin{center}
\begin{tabularx}{\linewidth}{|X|X|X|}
\hline
\rowcolor{poppurpleL}
\textbf{English} & \textbf{Français} & \textbf{Contexte / Collocation} \\
\hline
gender equality & l'égalité des sexes & \eng{to promote gender equality} \\
women's rights & les droits des femmes & \eng{to defend women's rights} \\
discrimination (against) & la discrimination (envers) & \eng{gender-based discrimination} \\
empowerment & l'autonomisation / le renforcement des capacités & \eng{women's economic empowerment} \\
equal opportunities & l'égalité des chances & \eng{equal opportunities policy} \\
a role model & un modèle / un exemple à suivre & \eng{She is a role model for girls.} \\
gender bias & un préjugé sexiste / un biais de genre & \eng{unconscious gender bias} \\
to advocate for & plaider pour / défendre (une cause) & \eng{to advocate for gender parity} \\
parity & la parité & \eng{gender parity in politics} \\
\hline
\end{tabularx}
\end{center}
\end{textebox}

\begin{propriete}[Morphologie : La famille du mot \eng{Equal}]
C'est une famille cruciale pour l'examen (transformation de phrases, composition). Maîtrisez les suffixes :
\begin{center}
\begin{tabular}{|l|l|l|l|}
\hline
\rowcolor{poporangeL}
\textbf{Adjectif} & \textbf{Nom} & \textbf{Adverbe} & \textbf{Contraire (préfixe \eng{un-})} \\
\hline
equal & equality & equally & unequal / inequality \\
\hline
\end{tabular}
\end{center}
\begin{itemize}
    \item \eng{Equal} (adj.) : égal(e). \eng{Equal pay for equal work.} (À travail égal, salaire égal.)
    \item \eng{Equality} (n.) : l'égalité. \eng{Gender equality is a human right.}
    \item \eng{Equally} (adv.) : équitablement, également. \eng{Boys and girls must be treated equally.}
    \item \eng{Unequal} (adj.) / \eng{Inequality} (n.) : inégal / inégalité. \eng{Social inequalities persist.}
\end{itemize}
\end{propriete}

\courssub{Champ lexical 3 : Les actions de sensibilisation et de plaidoyer}

Les ONG et associations mènent des actions concrètes. Le vocabulaire suivant décrit ces activités opérationnelles.

\begin{textebox}[Extraits du script d'écoute — Action context]
The NGO launched a nationwide \textbf{awareness campaign} against early marriage. 
\textit{L'ONG a lancé une campagne de sensibilisation nationale contre le mariage précoce.}

They organise \textbf{workshops} and \textbf{training programmes} for community leaders. 
\textit{Ils organisent des ateliers et des programmes de formation pour les leaders communautaires.}

\textbf{Mentoring} and \textbf{advocacy} are key strategies to change mindsets. 
\textit{Le mentorat et le plaidoyer sont des stratégies clés pour changer les mentalités.}
\end{textebox}

\begin{textebox}[TABLEAU 3 — Actions : Vocabulaire opérationnel]
\begin{center}
\begin{tabularx}{\linewidth}{|X|X|X|}
\hline
\rowcolor{poppinkL}
\textbf{English} & \textbf{Français} & \textbf{Verbe associé / Collocation} \\
\hline
an awareness campaign & une campagne de sensibilisation & \eng{to launch / run a campaign} \\
to raise awareness (about) & sensibiliser (à / sur) & \eng{to raise awareness about girls' rights} \\
a workshop & un atelier (participatif) & \eng{to attend / conduct a workshop} \\
a training programme & un programme de formation & \eng{to provide training} \\
mentoring (n.) & le mentorat / le tutorat & \eng{a mentoring scheme} \\
advocacy (n.) & le plaidoyer (action politique/sociale) & \eng{to do advocacy work} \\
to advocate (for) & plaider pour / défendre & \eng{to advocate for policy change} \\
outreach (n.) & action de terrain / porte-à-porte & \eng{community outreach} \\
\hline
\end{tabularx}
\end{center}
\end{textebox}

\begin{attention}[Faux amis et pièges de traduction — ALERTE BAC]
\begin{enumerate}
    \item \textbf{Formation $\neq$ \eng{a formation}.} En anglais, \eng{formation} existe mais signifie « formation géologique » ou « arrangement ». Pour « une formation professionnelle/pédagogique », on utilise \cle{training} (n. indénombrable) ou \cle{a training programme / course / session}.
    
    \item \textbf{Sensible $\neq$ \eng{sensible}.}
    \begin{itemize}
        \item \eng{Sensible} = raisonnable, sensé, pratique. \eng{A sensible decision.} (Une décision raisonnable.)
        \item \cle{Sensitive} = sensible, délicat, qui réagit vite. \eng{A sensitive topic.} (Un sujet sensible/délicat.) \eng{Sensitive skin.} (Peau sensible.)
    \end{itemize}
    
    \item \textbf{Assister $\neq$ \eng{to assist}.}
    \begin{itemize}
        \item \eng{To assist} = aider, secourir. \eng{The volunteer assisted the victims.} (Le bénévole a aidé les victimes.)
        \item \cle{To attend} = assister à (un événement, une réunion, un atelier). \eng{She attended the workshop.} (Elle a assisté à l'atelier.)
    \end{itemize}
    
    \item \textbf{Une campagne $\neq$ \eng{a campaign} seulement.} Attention à la collocation : \eng{to \textbf{launch} a campaign} (lancer), \eng{to \textbf{run} a campaign} (mener), \eng{to \textbf{carry out} a campaign} (réaliser). Pas *\eng{to make a campaign}.
\end{enumerate}
\end{attention}

\begin{exemplebox}[Exemples]
    \eng{She completed a training course in leadership.} (Elle a suivi une formation en leadership.) \\
    \eng{Teacher training is essential.} (La formation des enseignants est essentielle.)
    \end{exemplebox}


\courssub{Champ lexical 4 : L'éducation des filles et la scolarisation}

L'éducation est le levier principal cité dans les textes pour l'autonomisation. Voici le lexique spécialisé.

\begin{textebox}[Extraits du script d'écoute — Girls' Education context]
Many families struggle to \textbf{enrol} their daughters in secondary school. 
\textit{De nombreuses familles peinent à inscrire leurs filles au secondaire.}

High \textbf{dropout} rates among girls are often linked to poverty and early pregnancy. 
\textit{Les taux de décrochage élevés chez les filles sont souvent liés à la pauvreté et aux grossesses précoces.}

A \textbf{scholarship} can change a girl's life and help her obtain a \textbf{diploma}. 
\textit{Une bourse peut changer la vie d'une fille et l'aider à obtenir un diplôme.}
\end{textebox}

\begin{textebox}[TABLEAU 4 — Girls' Education : Vocabulaire clé]
\begin{center}
\begin{tabularx}{\linewidth}{|X|X|X|}
\hline
\rowcolor{popgoldL}
\textbf{English} & \textbf{Français} & \textbf{Précision / Exemple} \\
\hline
to enrol (US: enroll) & inscrire (à l'école) & \eng{to enrol in school} / \eng{enrolment} (n. = inscription/effectif) \\
schooling (n.) & la scolarité / la scolarisation & \eng{access to schooling} (accès à la scolarité) \\
dropout (n.) & le décrochage (scolaire) & \eng{a high dropout rate} (un taux de décrochage élevé) \\
to drop out (of school) & abandonner / décrocher (l'école) & \eng{She dropped out in Form 3.} \\
scholarship (n.) & une bourse (d'études) & \eng{to award / grant a scholarship} \\
diploma (n.) & un diplôme (souvent fin secondaire) & \eng{to obtain a diploma} / \eng{degree} (univ.) \\
literacy (n.) & l'alphabétisation & \eng{literacy rate} (taux d'alphabétisation) \\
illiterate (adj.) & illettré / analphabète & \eng{illiterate women} \\
\hline
\end{tabularx}
\end{center}
\end{textebox}

\courssub{Collocations clés : Les « blocs de sens » à mémoriser}

En anglais, les mots ne s'associent pas au hasard. Pour parler naturellement et réussir l'expression écrite, vous devez connaître ces associations figées (collocations). Apprenez-les par cœur comme du vocabulaire unique.

\begin{textebox}[Collocations incontournables — Volunteering \& Gender Equality]
\begin{center}
\begin{tabularx}{\linewidth}{|X|X|}
\hline
\rowcolor{popblueL}
\textbf{Collocation (Bloc anglais)} & \textbf{Traduction française} \\
\hline
to \textbf{raise awareness} (about / on) & sensibiliser (à / sur) \\
to \textbf{promote equality} / \textbf{gender equality} & promouvoir l'égalité / l'égalité des sexes \\
to \textbf{fight discrimination} / \textbf{combat discrimination} & lutter contre la discrimination \\
to \textbf{empower women} / \textbf{girls} & autonomiser les femmes / les filles \\
to \textbf{make a difference} & faire la différence / changer les choses \\
to \textbf{advocate for} change / rights & plaider pour le changement / les droits \\
to \textbf{launch a campaign} & lancer une campagne \\
to \textbf{conduct a workshop} / \textbf{training} & animer un atelier / une formation \\
to \textbf{provide scholarships} & offrir des bourses \\
to \textbf{reduce dropout rates} & réduire les taux de décrochage \\
\hline
\end{tabularx}
\end{center}
\end{textebox}

\begin{methode}[Comment apprendre et réemployer les collocations]
\begin{enumerate}
    \item \textbf{Identifiez} le verbe noyau (raise, promote, fight, empower, make, advocate, launch, conduct, provide, reduce).
    \item \textbf{Mémorisez} le nom qui le suit systématiquement (awareness, equality, discrimination, women, a difference, change, a campaign, a workshop, scholarships, dropout rates).
    \item \textbf{Conjuguez} le verbe à différents temps : \eng{The NGO \textbf{raised} awareness last year.} / \eng{They \textbf{will launch} a campaign next month.}
    \item \textbf{Transformez} en nominalisation pour l'écrit formel : \eng{The \textbf{raising of awareness} is crucial.} / \eng{\textbf{Promoting equality} requires funding.}
\end{enumerate}
\end{methode}

\courssub{Synthèse visuelle : Carte mentale lexicale}

La figure ci-dessous résume les quatre champs lexicaux et leurs mots satellites. Utilisez-la pour réviser en visualisant les liens sémantiques.

\begin{popfigure}
\begin{tikzpicture}[
    mindmap,
    grow cyclic,
    every node/.style={concept, align=center},
    concept color=popblue,
    level 1/.append style={level distance=4.5cm, sibling angle=90},
    level 2/.append style={level distance=3cm, sibling angle=45},
    text=white,
    font=\small\bfseries
]
\node[align=center]{Vocabulary\\Volunteering \&\\Gender Equality}
    child[concept color=popgreen] {
        node {Volunteering}
        child { node {volunteer (n/v)} }
        child { node {NGO} }
        child { node {commitment} }
        child { node {unpaid work} }
    }
    child[concept color=poppurple] {
        node[align=center]{Gender\\Equality}
        child { node {women's rights} }
        child { node {empowerment} }
        child { node {equal opportunities} }
        child { node {role model} }
    }
    child[concept color=poporange] {
        node[align=center]{Awareness\\Actions}
        child { node {raise awareness} }
        child { node {workshop} }
        child { node {advocacy} }
        child { node {mentoring} }
    }
    child[concept color=poppink] {
        node[align=center]{Girls'\\Education}
        child { node {enrol} }
        child { node {scholarship} }
        child { node {dropout} }
        child { node {diploma} }
    };
\end{tikzpicture}
\legende{Carte mentale des 4 champs lexicaux majeurs de la leçon 37.}
\end{popfigure}

\courssub{Texte original d'entraînement : « The Buea Girls' Club Initiative »}

Voici un texte original (environ 250 mots) réemployant l'essentiel du vocabulaire ci-dessus. Il sert de support de révision globale (reading comprehension + lexical reuse).

\begin{textebox}[The Buea Girls' Club Initiative]
In the bustling city of Buea, at the foot of Mount Cameroon, a local NGO named “Hope for Her” has been running a remarkable \textbf{awareness campaign} since 2020. Their main \textbf{objective} is to \textbf{promote gender equality} by keeping girls in school. 

Every Saturday, \textbf{volunteers} gather at the community hall to \textbf{conduct workshops} on \textbf{women's rights}, reproductive health, and leadership. These sessions \textbf{raise awareness} about the dangers of early marriage and \textbf{gender-based discrimination}. “We want every girl to know she has \textbf{equal opportunities},” says Ms. Aminata, the \textbf{volunteer coordinator}. 

The association also runs a \textbf{mentoring} programme. Successful women from Buea — teachers, nurses, entrepreneurs — act as \textbf{role models} for the teenagers. They share their journeys, proving that \textbf{empowerment} starts with education. 

Last year, the NGO \textbf{advocated for} a municipal decree to support teenage mothers returning to school. Thanks to their \textbf{advocacy work}, three girls who had \textbf{dropped out} were able to \textbf{re-enrol}. Furthermore, “Hope for Her” secured five university \textbf{scholarships} for top-performing students. 

“\textbf{Volunteering} changed my life,” confides Sarah, a Form 5 student. “Before, I was shy. Now, I \textbf{volunteer} to teach younger girls. I want to \textbf{make a difference} in my community.” This \textbf{commitment} illustrates how \textbf{unpaid work} builds not only a fairer society but also confident future leaders.
\end{textebox}

\begin{definition}[Glossaire du texte (mots difficiles)]
\begin{itemize}
    \item \cle{bustling} (adj.) : animé, bouillonnant, plein d'activité.
    \item \cle{remarkable} (adj.) : remarquable, exceptionnel.
    \item \cle{reproductive health} (n.) : la santé reproductive.
    \item \cle{teenage mothers} (n.) : des mères adolescentes.
    \item \cle{municipal decree} (n.) : un arrêté municipal.
    \item \cle{top-performing} (adj.) : les meilleurs / les plus performants.
    \item \cle{confides} (v.) : confie (verbe \eng{to confide in sb}).
    \item \cle{illustrates} (v.) : illustre, montre par l'exemple.
\end{itemize}
\end{definition}

\begin{exercice}[Questions de compréhension et réemploi lexical — Texte « The Buea Girls' Club Initiative »]
\begin{enumerate}
    \item \textbf{True or False?} Justify with a quote from the text.
    \begin{itemize}
        \item[a)] The NGO is called “Hope for Him”.
        \item[b)] Workshops take place every Sunday.
        \item[c)] The NGO helped three teenage mothers go back to school.
        \item[d)] Sarah is a university student.
    \end{itemize}
    \item \textbf{Vocabulary search.} Find in the text the English equivalents for:
    \begin{itemize}
        \item[a) sensibiliser (verbe + nom) \quad b) plaider pour (verbe + préposition) \quad c) décrocher (l'école) (verbe + particule) \quad d) un modèle (nom) \quad e) l'autonomisation (nom)]
    \end{itemize}
    \item \textbf{Collocation completion.} Complete the sentences using a collocation from the table in Section 4.4 (put the verb in the correct tense).
    \begin{itemize}
        \item[a)] The association \_\_\_\_\_\_\_\_\_\_\_\_ a new campaign next month. (lancer)
        \item[b)] Education helps to \_\_\_\_\_\_\_\_\_\_\_\_ women economically. (autonomiser)
        \item[c)] We must \_\_\_\_\_\_\_\_\_\_\_\_ against all forms of discrimination. (lutter)
    \end{itemize}
\end{enumerate}
\end{exercice}

\begin{corrige}[Exercice — The Buea Girls' Club Initiative]
\begin{enumerate}
    \item \textbf{True/False + Justification}
    \begin{itemize}
        \item[a)] \textbf{False.} Quote: “a local NGO named “Hope for \textbf{Her}”.”
        \item[b)] \textbf{False.} Quote: “Every \textbf{Saturday}, volunteers gather...”
        \item[c)] \textbf{True.} Quote: “three girls who had dropped out were able to re-enrol.”
        \item[d)] \textbf{False.} Quote: “Sarah, a \textbf{Form 5 student}” (élève de 1ère/Terminale).
    \end{itemize}
    \item \textbf{Vocabulary search}
    \begin{itemize}
        \item[a)] \eng{raise awareness}
        \item[b)] \eng{advocate for}
        \item[c)] \eng{drop out}
        \item[d)] \eng{a role model}
        \item[e)] \eng{empowerment}
    \end{itemize}
    \item \textbf{Collocation completion}
    \begin{itemize}
        \item[a)] \eng{will launch} (future simple, prédiction/plan) ou \eng{is launching} (futur proche planifié).
        \item[b)] \eng{empower women} (infinitif après \eng{helps to}).
        \item[c)] \eng{fight discrimination} (infinitif après modal \eng{must}).
    \end{itemize}
\end{enumerate}
\end{corrige}

\courssub{Exercice résolu : Transformation de phrases (Grammaire + Lexique)}

Cet exercice type « Bac » teste votre capacité à manipuler les familles de mots (dérivation) et les collocations.

\begin{exoresolu}[Transformation — Derivation and Collocations]
\textbf{Énoncé :} Rewrite the sentences using the word in brackets. Do not change the meaning. 

1. The NGO \textbf{sensitises} communities \textbf{about} girls' education. (\eng{awareness}) \\
2. \textbf{Voluntary work} helps young people gain confidence. (\eng{volunteer}) \\
3. We must \textbf{fight} \textbf{discrimination} against women. (\eng{combat}) \\
4. She is a \textbf{model} for all young girls. (\eng{role}) \\
5. The \textbf{empowerment} of women is a priority. (\eng{empower}) \\
6. Many girls \textbf{leave school early} because of poverty. (\eng{dropout}) \\

\textbf{Solution détaillée :}

1. \eng{The NGO \textbf{raises awareness about} girls' education.} 
\begin{itemize}
    \item \textit{Analyse :} Le verbe \eng{sensitise} (UK) / \eng{sensitize} (US) est moins courant que la collocation nominale \eng{raise awareness}. Le nom \eng{awareness} impose le verbe \eng{raise}. Préposition conservée : \eng{about}.
\end{itemize}

2. \eng{\textbf{Volunteering} helps young people gain confidence.} 
\begin{itemize}
    \item \textit{Analyse :} \eng{Voluntary work} = travail bénévole (adj + nom). Le mot imposé \eng{volunteer} (verbe/nom de base) doit devenir le nom d'action (gérondif) \eng{volunteering} pour fonctionner comme sujet. \eng{To volunteer} aide aussi mais \eng{Volunteering} est plus naturel comme sujet général.
\end{itemize}

3. \eng{We must \textbf{combat discrimination} against women.} 
\begin{itemize}
    \item \textit{Analyse :} \eng{Fight} (verbe) $\rightarrow$ \eng{combat} (verbe transitif direct, plus formel/écrit). Attention : \eng{combat} ne prend pas de préposition ici (contrairement à \eng{fight against}). \eng{Combat discrimination} = collocation formelle.
\end{itemize}

4. \eng{She is a \textbf{role model} for all young girls.} 
\begin{itemize}
    \item \textit{Analyse :} \eng{Model} seul peut prêter à confusion (mannequin, maquette). L'expression figée pour « modèle à suivre » est \eng{role model} (nom composé).
\end{itemize}

5. \eng{\textbf{Empowering} women is a priority.} \quad OU \quad \eng{To \textbf{empower} women is a priority.}
\begin{itemize}
    \item \textit{Analyse :} \eng{Empowerment} (nom) $\rightarrow$ forme verbale. Le sujet peut être le gérondif \eng{Empowering} (action générale) ou l'infinitif \eng{To empower} (but). Les deux sont corrects.
\end{itemize}

6. \eng{Many girls \textbf{drop out} of school because of poverty.} 
\begin{itemize}
    \item \textit{Analyse :} \textbf{Leave school early} $\rightarrow$ verbe phrasal \eng{drop out} (intransitif). Nécessite la préposition \eng{of} devant le nom \eng{school}. \eng{Dropout} (nom) ne peut pas fonctionner comme verbe ici.
\end{itemize}
\end{exoresolu}

\courssub{Listening strategies and worked example}

Cette section explicite les stratégies pour réussir l'épreuve de compréhension de l'oral (Listening) sur ce thème, conformément aux attendus du Baccalauréat.

\begin{methode}[Stratégies d'écoute : Repérer l'essentiel en 4 étapes]
\begin{enumerate}
    \item \textbf{Anticipation (Avant l'écoute).} Lisez le titre, les questions, les images. Cochez le vocabulaire probable : \eng{NGO, volunteer, gender equality, workshop, scholarship, dropout, campaign}. Prédisez le thème (ex. : « Une ONG aide des filles à l'école »).
    \item \textbf{Écoute globale (1re écoute).} Identifiez : \textbf{Who?} (locuteurs : journaliste, bénévole, bénéficiaire), \textbf{Where?} (Buea, centre communautaire, école), \textbf{What?} (thème principal : éducation des filles, mariage précoce, mentorat), \textbf{Tone?} (optimiste, urgent, informatif).
    \item \textbf{Écoute détaillée (2e écoute).} Concentrez-vous sur les \textbf{mots-clés} et \textbf{chiffres} : dates (2020, last year), nombres (five scholarships, three girls), noms propres (Hope for Her, Ms. Aminata, Sarah, Mount Cameroon), verbes d'action (\eng{launched, conduct, raise, advocated, secured, dropped out, re-enrol}).
    \item \textbf{Vérification \& Inférence (3e écoute / pause).} Relisez vos notes. Complétez les trous par déduction logique (contexte, connecteurs logiques : \eng{however, therefore, because}). Attention aux \textbf{pièges phonétiques} : \eng{volunteer / voluntary}, \eng{advocate (verbe) / advocacy (nom)}, \eng{enrol / enrollment}.
\end{enumerate}
\end{methode}

\begin{exemplebox}[Exemple travaillé — Script simulé (extrait 1 min 30)]
\textbf{Script :} “Good morning. I'm James reporting for Cameroon Radio Television from Buea. Today, we focus on the ‘Hope for Her’ initiative. Founded in 2020, this NGO tackles gender inequality through education. Coordinator Ms. Aminata explains: ‘We conduct weekly workshops on rights and health. Our mentoring scheme connects girls with professional women. Last year, we lobbied the council for a decree allowing teen mothers back in class. Three girls re-enrolled. We also funded five university scholarships.’ Sarah, a Form 5 student, says: ‘Volunteering here built my confidence. I now teach younger girls.’ A powerful example of how unpaid work changes lives.”

\textbf{Questions types (QCM / Vrai-Faux / Complétion) :}
\begin{enumerate}
    \item What is the name of the NGO? (A) Hope for Him (B) Hope for Her (C) Hope for All \rightarrow \textbf{B}
    \item How many scholarships were secured? \rightarrow \textbf{Five}
    \item True or False: Workshops are held on Sundays. \rightarrow \textbf{False (Saturdays)}
    \item What phrasal verb means “return to school”? \rightarrow \textbf{re-enrol} (ou \eng{go back to school})
    \item What does Sarah say volunteering gave her? \rightarrow \textbf{Confidence / She became a teacher for younger girls.}
\end{enumerate}
\end{exemplebox}

\courssub{Guided production: Reusing the lexicon}

Mettez en pratique le vocabulaire et les collocations dans des tâches d'expression orale et écrite guidées.

\begin{experience}[Jeu de rôle : Interview radio « Youth Engagement »]
\textbf{Consigne :} Par binômes. L'un est journaliste (CRTV / BBC Africa), l'autre est coordinateur/trice d'une ONG locale (ex. « Youth for Change », « Girls Rise »). Durée : 3 min.
\textbf{Contraintes lexicales :} Utilisez au moins \textbf{8 mots/collocations} de la leçon (ex. \eng{raise awareness, advocate for, dropout, scholarship, mentoring, volunteer, commitment, gender equality, launch a campaign, role model}).
\textbf{Structure suggérée :}
\begin{enumerate}
    \item Introduction : Nom ONG, localisation, mission (\eng{promote gender equality}).
    \item Actions concrètes : \eng{workshops, mentoring, outreach, campaign}.
    \item Résultats chiffrés : \eng{scholarships, reduced dropout rates, re-enrolment}.
    \item Témoignage bénévole : \eng{make a difference, empowerment, skills}.
    \item Appel à l'action : \eng{join, sign up, volunteer, support}.
\end{enumerate}
\end{experience}

\begin{exercice}[Production écrite guidée : Article de blog / Discours]
\textbf{Sujet :} « Write a 150-word blog post for your school magazine entitled “Why I Volunteer for Gender Equality”. »
\textbf{Plan imposé :}
\begin{enumerate}
    \item \textbf{Hook} : Accroche (statistique, citation, anecdote personnelle).
    \item \textbf{Context} : Présentez l'association/ONG (\eng{NGO, volunteer coordinator, commitment}).
    \item \textbf{Actions} : Décrivez 2 activités (\eng{conduct workshops, raise awareness, mentoring, advocacy}).
    \item \textbf{Impact} : Résultats (\eng{scholarships, dropout rates, role models, empowerment}).
    \item \textbf{Call to action} : Incitez les lecteurs (\eng{join us, sign up, make a difference, volunteer}).
\end{enumerate}
\textbf{Mots-outils (connecteurs) :} \eng{Furthermore, Moreover, Consequently, For instance, In conclusion.}
\end{exercice}

\begin{aretenir}[Check-list de révision — Vocabulaire Leçon 37]
Avant de passer à l'évaluation, assurez-vous de savoir :
\begin{itemize}
    \item \textbf{Définir} en anglais : \eng{NGO, volunteer, empowerment, advocacy, dropout, scholarship}.
    \item \textbf{Traduire} sans hésiter les 10 collocations de la section 4.4 (\eng{raise awareness, promote equality...}).
    \item \textbf{Former} la famille de \eng{equal} (equality, equally, unequal, inequality) et de \eng{volunteer} (voluntary, voluntarily).
    \item \textbf{Éviter} les 4 faux amis de la boîte \textbf{Attention} (\eng{formation/training, sensible/sensitive, assist/attend, make/launch a campaign}).
    \item \textbf{Rédiger} 5 phrases personnelles en réemployant 5 mots/collocations différents de cette leçon (ex. : \eng{Last summer, I volunteered for an NGO that promotes gender equality.}).
    \item \textbf{Appliquer} les 4 étapes de la méthode d'écoute sur un nouvel extrait audio/vidéo.
\end{itemize}
\end{aretenir}

\begin{savaistu}[Le volontariat au Cameroun : Contexte réel]
Le Cameroun compte des milliers d'associations de jeunesse et d'ONG locales (ex. \eng{RELUFA}, \eng{CAMENA}, \eng{AJESH}) actives sur l'égalité de genre. Le \textbf{Ministère de la Promotion de la Femme et de la Famille (MINPROFF)} pilote la stratégie nationale. Le « \eng{Volontariat National} » (programme gouvernemental) permet aux jeunes diplômés de servir dans des structures publiques ou associatives pour une indemnité modique, acquérant ainsi une première \eng{professional experience}. Dans les zones anglophones (Nord-Ouest, Sud-Ouest), les \eng{Community Based Organisations (CBOs)} jouent un rôle crucial dans la scolarisation des filles déplacées par la crise sociopolitique. Connaître ce contexte enrichit vos exemples à l'oral et à l'écrit (Writing/Composition).
\end{savaistu}

\coursec{Listening strategies and worked example}

Dans la section précédente, nous avons constitué notre boîte à outils lexicale : \eng{volunteer}, \eng{gender equality}, \eng{empowerment}, \eng{awareness campaign}, \eng{certificate}... Mais connaître les mots ne suffit pas : il faut savoir les \emph{entendre}. Un enregistrement ne s'arrête pas pour vous attendre, et au Bac, vous ne l'entendrez que deux ou trois fois. Cette section vous donne une méthode complète, en cinq stratégies, pour transformer l'écoute en réflexe — puis un exemple traité pas à pas sur notre script sur le volontariat.

\courssub{Why you need a strategy, not just good ears}

Écouter en langue étrangère, ce n'est pas « ouvrir les oreilles et espérer ». C'est une activité \cle{active} : on prédit, on sélectionne, on note, on vérifie. Les bons élèves ne comprennent pas plus de mots que les autres — ils savent simplement \emph{quoi} chercher et \emph{quand}.

\begin{popfigure}
\begin{tikzpicture}[
  etape/.style={rectangle, rounded corners=4pt, draw=popblue, fill=popblueL, align=center, minimum width=2.6cm, minimum height=1.1cm, font=\small},
  fleche/.style={-{Stealth[length=2.5mm]}, thick, poporange}
]
\node[etape, align=center] (s1) at (0,0){\textbf{1. Predict}\\ read questions};
\node[etape, align=center] (s2) at (3.4,0){\textbf{2. Listen 1}\\ gist + tone};
\node[etape, align=center] (s3) at (6.8,0){\textbf{3. Listen 2}\\ targeted notes};
\node[etape, align=center] (s4) at (10.2,0){\textbf{4. Listen 3}\\ check};
\node[etape, draw=popgreen, fill=popgreenL, align=center] (s5) at (13.6,0){\textbf{5. Answer}\\ + justification};
\draw[fleche] (s1) -- (s2);
\draw[fleche] (s2) -- (s3);
\draw[fleche] (s3) -- (s4);
\draw[fleche] (s4) -- (s5);
\end{tikzpicture}
\legende{La méthode d'écoute en cinq étapes : Predict $\rightarrow$ Listen 1 (gist) $\rightarrow$ Listen 2 (notes) $\rightarrow$ Listen 3 (check) $\rightarrow$ Answer with justification.}
\end{popfigure}

\courssub{Strategy 1 — Before listening: read the questions and predict}

\begin{propriete}[Prédire avant d'écouter]
Avant la première écoute, lisez \textbf{toutes} les questions. Pour chacune, demandez-vous : \emph{quel type d'information dois-je capter ?} Un nombre ? Un nom propre ? Un lieu ? Une date ? Cette prédiction oriente votre attention : votre cerveau saura reconnaître l'information quand elle passera.
\end{propriete}

\begin{exemplebox}[Prédiction à partir des questions]
Question : \eng{How many families benefited from the project?} $\rightarrow$ je dois capter un \textbf{nombre} suivi de \eng{families}.

Question : \eng{Where did the volunteers work?} $\rightarrow$ je dois capter un \textbf{lieu} (ville, quartier, région).

Question : \eng{True or False: The project started last year.} $\rightarrow$ je dois capter une \textbf{expression de temps} (\eng{last year}, \eng{six months ago}, \eng{in January}...).
\end{exemplebox}

\courssub{Strategy 2 — First listening: no pen, catch the gist}

\begin{propriete}[Première écoute = vue d'ensemble]
Lors de la première écoute, \textbf{posez votre stylo}. Ne notez rien. Concentrez-vous sur trois questions : \emph{Who is speaking? What is the general topic? What is the tone} (optimistic, critical, informative)\emph{?} C'est ce qu'on appelle le \cle{gist} (l'idée générale).
\end{propriete}

Pourquoi sans stylo ? Parce qu'écrire et écouter en même temps divise l'attention : pendant que vous notez un détail, vous ratez la phrase suivante. Le gist, lui, vous donne la « carte » du texte — les détails viendront à la deuxième écoute.

\courssub{Strategy 3 — Second listening: targeted note-taking}

\begin{propriete}[Prise de notes ciblée]
À la deuxième écoute, notez \textbf{uniquement} ce que les questions exigent : les \cle{chiffres}, les \cle{noms propres}, les \cle{lieux}, les \cle{dates}. Écrivez-les en marge, à côté du numéro de la question correspondante. Ignorez le reste.
\end{propriete}

\courssub{Strategy 4 — Use abbreviations}

Vous n'avez pas le temps d'écrire des mots entiers. Créez un code personnel :

\begin{center}
\begin{tabularx}{\linewidth}{|X|X|X|}\hline
\rowcolor{popblueL} \textbf{Abréviation} & \textbf{Mot complet} & \textbf{Français} \\ \hline
vol. & volunteers & volontaires \\ \hline
fam. & families & familles \\ \hline
N / S / E / W & North / South / East / West & Nord / Sud / Est / Ouest \\ \hline
certif. & certificates & attestations \\ \hline
gvt & government & gouvernement \\ \hline
$\uparrow$ / $\downarrow$ & increase / decrease & augmenter / diminuer \\ \hline
b/c & because & parce que \\ \hline
\end{tabularx}
\end{center}

\courssub{Strategy 5 — Spot the signal words}

\begin{propriete}[Les mots-outils annonciateurs]
Certaines expressions annoncent qu'une information importante arrive. Entraînez votre oreille à les reconnaître instantanément :
\begin{itemize}
\item les \textbf{nombres} : \eng{two hundred}, \eng{fifty}, \eng{three-quarters} ;
\item les \textbf{repères temporels} : \eng{so far} (jusqu'à présent), \eng{last month}, \eng{since January}, \eng{by the end of the year} ;
\item les \textbf{citations} : \eng{she said}, \eng{according to the coordinator}, \eng{in her own words} ;
\item les \textbf{conclusions} : \eng{as a result}, \eng{that is why}, \eng{finally}.
\end{itemize}
\end{propriete}

\begin{attention}[Piège : fifteen ou fifty ?]
À l'écoute, les nombres en \eng{-teen} et en \eng{-ty} se confondent facilement. La clé est l'\textbf{accent tonique} : on dit fif\emph{TEEN} (15, accent sur la fin) mais \emph{FIF}ty (50, accent sur le début). Entraînez-vous avec les paires 13/30, 14/40, 15/50, 16/60, 17/70, 18/80, 19/90. Une erreur ici peut faire perdre deux points sur un vrai/faux !
\end{attention}

\begin{methode}[Réussir l'épreuve de listening au Bac]
\begin{enumerate}
\item Lis \textbf{toutes} les questions avant la première écoute.
\item Souligne les \textbf{mots interrogatifs} (\eng{how many}, \eng{where}, \eng{when}, \eng{who}) : ils te disent quoi écouter.
\item Note les \textbf{chiffres dès que tu les entends} — ils ne seront pas répétés.
\item Ne reste \textbf{jamais bloqué} sur un mot inconnu : devine son sens grâce au contexte et continue d'écouter.
\end{enumerate}
\end{methode}

\courssub{Worked example — step by step}

Appliquons la méthode à un extrait de notre script sur le volontariat pour l'égalité des sexes.

\begin{textebox}[Listening script — extract]
“So far, our volunteers have helped more than two hundred families in the North Region. Last month, fifty girls received certificates after completing our training programme in tailoring and computer skills. According to the coordinator, Mrs Amina Bello, the project will reach three hundred families by the end of the year.”
\end{textebox}

\begin{definition}[Glossaire de l'extrait]
\begin{itemize}
\item \eng{so far} (locution) : jusqu'à présent
\item \eng{tailoring} (nom) : couture
\item \eng{to reach} (verbe) : atteindre
\item \eng{coordinator} (nom) : coordinateur / coordinatrice
\end{itemize}
\end{definition}

\begin{exoresolu}[Vrai/Faux justifié]
Énoncé — \eng{True or False: Two hundred girls received certificates last month.} Justify with a quotation from the script. \tcblower
Solution détaillée — Étape 1 (Predict) : la question porte sur un \textbf{nombre} et sur \textbf{qui} a reçu des attestations. Étape 2 (Listen 2) : je note « 200 fam. » et « 50 girls certif. ». Étape 3 (Check) : le script dit \eng{two hundred families} mais \eng{fifty girls received certificates}. La phrase mélange les deux informations : c'est donc \textbf{FALSE}. Réponse rédigée : \eng{False. The script says “fifty girls received certificates”, and the two hundred concerns families, not girls.} (Faux. Le script dit « cinquante filles ont reçu des attestations », et les deux cents concernent des familles, pas des filles.)
\end{exoresolu}

\begin{exemplebox}[Question à réponse courte — modèle]
Question : \eng{How many girls received certificates?} $\rightarrow$ Indice sonore repéré : \eng{“Last month, fifty girls received certificates...”} $\rightarrow$ Réponse : \eng{Fifty girls.} « Cinquante filles ont reçu des attestations. » Remarque : une réponse courte et exacte suffit ; inutile de recopier toute la phrase.
\end{exemplebox}

\begin{savaistu}[Des points faciles au Bac]
Au Bac, la compréhension orale et écrite rapporte des points « faciles » à condition de \textbf{justifier systématiquement} chaque réponse par une courte citation entre guillemets : \eng{True — “fifty girls received certificates”}. Un vrai/faux non justifié ne vaut souvent que la moitié des points, voire zéro. La citation prouve que tu as réellement compris, et le correcteur n'a aucune raison d'hésiter.
\end{savaistu}

\begin{exercice}[À ton tour]
Réécoute mentalement l'extrait du script ci-dessus et réponds en anglais, avec justification : 1) \eng{How many families will the project reach by the end of the year?} 2) \eng{True or False: Mrs Amina Bello is a volunteer tailor.} 3) \eng{In which region did the volunteers work?}
\end{exercice}

\begin{corrige}[Exercice — Listening strategies]
1) \eng{Three hundred families — “the project will reach three hundred families by the end of the year.”} (Trois cents familles.) 2) \eng{False — she is “the coordinator”, not a tailor.} (Faux : elle est coordinatrice.) 3) \eng{In the North Region — “in the North Region”.} (Dans la région du Nord.)
\end{corrige}

\coursec{Guided production: reuse the lexicon}

Dans la section précédente, nous avons appris à \cle{écouter efficacement} : repérer le thème, noter les mots-clés, les chiffres et les noms propres, puis vérifier nos réponses. Mais une compétence d'écoute n'est vraiment acquise que si l'on sait \cle{réemployer} ce que l'on a entendu. C'est exactement l'objectif de cette dernière étape : faire passer le lexique du volontariat pour l'égalité des sexes de la \emph{compréhension} à la \emph{production}, d'abord écrite (phrases à trous, \emph{present perfect}, résumé guidé), puis orale (appel à la radio scolaire, écoute active en binômes).

\courssub{Sentence frames: your safety net}

Avant de produire, observons les \cle{trames de phrases} (\eng{sentence frames}) qui vous permettront de parler du script sans hésiter. Ces structures figées sont des « échafaudages » : vous ne changez que les éléments variables (le nom, le lieu, le chiffre).

\begin{textebox}[SENTENCE FRAMES — talking about the recording]
“The recording is about a volunteer who promotes gender equality.”

“The volunteer's name is Amina, and she works in the Far North region.”

“So far, her team has trained more than two hundred girls.”

“If you want to join, call 677-45-32-10 or visit our office.”

“Volunteering changes lives because it gives hope and skills to young people.”
\end{textebox}

\begin{propriete}[Forme et emploi des trames]
\begin{itemize}
\item \eng{The recording is about + nom / nominal en -ing} : \eng{The recording is about volunteering.} « L'enregistrement parle du volontariat. » Attention : \eng{about} est suivi d'un nom ou d'un \emph{gérondif}, jamais d'un infinitif (\eng{about to volunteer} est faux ici).
\item \eng{So far + present perfect} : \eng{So far, her team has helped 200 families.} « Jusqu'ici, son équipe a aidé 200 familles. » \eng{So far} (« jusqu'à présent ») exige le \cle{present perfect} : \textbf{have/has + participe passé}.
\item \eng{If you want to join, + impératif} : \eng{If you want to join, call us today.} « Si tu veux participer, appelle-nous aujourd'hui. » Après \eng{if} au présent, la principale est à l'impératif ou au futur (first conditional).
\end{itemize}
\end{propriete}

\begin{attention}[Faux amis et pièges des francophones]
\begin{itemize}
\item Ne dites jamais \eng{I want to make voluntary}. Dites \eng{I want to volunteer} (« je veux faire du bénévolat ») ou \eng{I want to do voluntary work} (« je veux faire du travail bénévole »). \eng{Voluntary} est un \emph{adjectif}, pas un nom d'activité.
\item Le verbe \eng{to sensitise} n'existe pas en anglais ! Pour « sensibiliser », dites \eng{to raise awareness} : \eng{We raise awareness about girls' education.} « Nous sensibilisons à l'éducation des filles. »
\item « Un volontaire » se dit \eng{a volunteer} (prononcé « vo-lon-tir ») ; \eng{voluntary work} est le travail, \eng{a voluntary organisation} l'organisation.
\end{itemize}
\end{attention}

\courssub{Réemploi 1 : compléter des phrases à trous}

\begin{exoresolu}[Gap-filling with the theme lexicon]
\textbf{Énoncé} — Complete the sentences with: \emph{volunteer, raise awareness, workshop, gender equality, leadership}.

1. Amina decided to \ldots\ldots\ldots for an NGO after secondary school.
2. The organisation fights for \ldots\ldots\ldots between boys and girls.
3. Every Saturday, the team organises a \ldots\ldots\ldots for young mothers.
4. Posters and radio talks help to \ldots\ldots\ldots about early marriage.
5. The training programme teaches girls \ldots\ldots\ldots skills.
\tcblower
\textbf{Solution détaillée}
1. \textbf{volunteer} — après \eng{decided to}, il faut un infinitif ; « faire du bénévolat » = \eng{to volunteer}.
2. \textbf{gender equality} — « l'égalité entre garçons et filles » ; \eng{between} confirme l'idée d'égalité.
3. \textbf{workshop} — « un atelier » ; l'article \eng{a} exige un nom singulier comptable.
4. \textbf{raise awareness} — « sensibiliser » ; structure fixe \eng{raise awareness about + nom}.
5. \textbf{leadership} — « des compétences en leadership » ; adjectif épithète devant \eng{skills}.
\end{exoresolu}

\courssub{Réemploi 2 : du tableau de notes au present perfect}

Pendant l'écoute, vous avez pris des notes sous forme de tableau. Transformons-les maintenant en \cle{phrases complètes au present perfect}, le temps des bilans d'actions (« jusqu'ici, l'équipe a… »).

\begin{center}
\begin{tabularx}{\linewidth}{|X|X|X|}
\hline
\rowcolor{popblueL} Notes (script) & Phrase au present perfect & Traduction \\
\hline
team / train / 200 girls & The team has trained 200 girls. & L'équipe a formé 200 filles. \\
\hline
volunteers / organise / 12 workshops & The volunteers have organised 12 workshops. & Les volontaires ont organisé 12 ateliers. \\
\hline
Amina / visit / 30 villages & Amina has visited 30 villages. & Amina a visité 30 villages. \\
\hline
the NGO / build / a centre & The NGO has built a training centre. & L'ONG a construit un centre de formation. \\
\hline
\end{tabularx}
\end{center}

\begin{propriete}[Rappel : le present perfect]
\textbf{Forme} : sujet + \textbf{have/has} + participe passé (\eng{has trained}, \eng{have organised}). \textbf{Emploi} : bilan d'actions passées dont le résultat compte maintenant, souvent avec \eng{so far}, \eng{already}, \eng{just}, \eng{since}, \eng{for}. \textbf{Comparaison} : le français utilise le passé composé (« a formé »), mais l'anglais refuse tout complément de temps \emph{terminé} (\eng{yesterday}, \eng{last year}) avec ce temps : on dira \eng{Last year, we helped 200 families} (prétérit), mais \eng{So far, we have helped 200 families} (present perfect).
\end{propriete}

\begin{exercice}[Your turn]
Transformez ces notes en phrases au present perfect : 1. she / talk / to many parents ; 2. we / raise / awareness in schools ; 3. they / help / 200 families so far ; 4. the centre / welcome / 50 girls this year.
\end{exercice}

\begin{corrige}[Exercice « Your turn »]
1. She has talked to many parents. 2. We have raised awareness in schools. 3. So far, they have helped 200 families. 4. The centre has welcomed 50 girls this year. (Attention : \eng{this year} inachevé accepte le present perfect ; \eng{last year} l'interdirait.)
\end{corrige}

\courssub{Réemploi 3 : résumé guidé du script}

Résumez le script d'écoute en \cle{trois phrases} à l'aide des trames. Modèle attendu :

\begin{exemplebox}[Guided summary — model answer]
\eng{The speaker presents Amina, a young volunteer from Maroua. She works with an NGO that promotes girls' education and gender equality in the Far North. Her dream is to see every girl finish secondary school and become a leader in her community.}

« La présentatrice présente Amina, une jeune volontaire de Maroua. Elle travaille avec une ONG qui promeut l'éducation des filles et l'égalité des sexes dans l'Extrême-Nord. Son rêve est de voir chaque fille terminer le secondaire et devenir un leader dans sa communauté. »
\end{exemplebox}

\courssub{Final production: a radio appeal for volunteers}

\begin{methode}[How to write a volunteer appeal in four moves]
\begin{enumerate}
\item \textbf{Hook} : commencez par une \cle{question accrocheuse} — \eng{Do you want to make a difference?}
\item \textbf{Problem} : présentez le \cle{problème en une phrase} — \eng{In the Far North, many girls still leave school early.}
\item \textbf{Action + chiffre fort} : montrez ce que fait l'ONG avec un \cle{nombre précis} — \eng{Last year, we helped 200 families.}
\item \textbf{Call to join} : terminez par une \cle{invitation claire} (téléphone) et un \cle{slogan} — \eng{Call 677-45-32-10 today. Volunteering changes lives!}
\end{enumerate}
\end{methode}

\begin{popfigure}
\begin{center}
\begin{tikzpicture}[node distance=0.55cm,
box/.style={draw=popblue, fill=popblueL, rounded corners, align=center, text width=3.1cm, minimum height=1.1cm, font=\small},
arr/.style={-{Stealth[length=2.5mm]}, thick, poporange}]
\node[box, align=center] (hook){\textbf{1. Hook}\\ “Do you want to make a difference?”};
\node[box, right=of hook, align=center] (prob){\textbf{2. Problem}\\ “Many girls leave school early.”};
\node[box, right=of prob, align=center] (act){\textbf{3. Action}\\ “We trained 200 girls last year.”};
\node[box, right=of act, align=center] (call){\textbf{4. Call to join}\\ “Call 677-45-32-10!”};
\draw[arr] (hook) -- (prob);
\draw[arr] (prob) -- (act);
\draw[arr] (act) -- (call);
\end{tikzpicture}
\end{center}
\legende{Structure d'un appel au volontariat : accroche, problème, action chiffrée, invitation.}
\end{popfigure}

\begin{textebox}[MODEL APPEAL — read and imitate]
“Do you want to make a difference? Join Equal Chances! In the Far North, many girls still leave school early. Our volunteers talk to parents, organise workshops and train young girls in leadership. Last year, we helped 200 families. You can help too. Call 677-45-32-10 today. Volunteering changes lives — starting with yours!”

\textbf{Glossaire} : \emph{to make a difference} (locution, faire une différence / changer les choses) ; \emph{to join} (verbe, rejoindre) ; \emph{early} (adverbe, tôt, prématurément) ; \emph{workshop} (nom, atelier) ; \emph{leadership} (nom, leadership, capacité à diriger).
\end{textebox}

\begin{exercice}[Write your own appeal]
Rédigez un appel au volontariat de 5 à 6 lignes pour la radio scolaire de votre lycée. Respectez les quatre étapes du schéma, réemployez au moins quatre mots du lexique (\eng{volunteer, raise awareness, workshop, gender equality, leadership}) et un chiffre au present perfect ou au prétérit.
\end{exercice}

\courssub{Pair work: active listening}

\begin{experience}[Radio appeal and keyword hunt]
Travaillez en binômes. L'élève A lit son appel à voix haute, lentement, comme à la radio. L'élève B, sans regarder le texte, relève sur une feuille les \cle{mots-clés} entendus : le nom de l'ONG, le problème, le chiffre, le numéro de téléphone et le slogan. Puis B vérifie avec le texte de A : a-t-il tout capté ? Échangez ensuite les rôles. Critère de réussite : B a retrouvé au moins 4 informations sur 5 — sinon, A doit articuler davantage ou mieux choisir ses mots-clés !
\end{experience}

\begin{aretenir}[Ce qu'il faut retenir]
\begin{itemize}
\item Réemployer le lexique entendu, c'est le mémoriser durablement : trous, present perfect, résumé, production libre forment une progression du guidé vers l'autonome.
\item Un appel au volontariat suit toujours quatre mouvements : \textbf{Hook → Problem → Action → Call to join}.
\item Dites \eng{to volunteer} / \eng{voluntary work} (jamais \eng{to make voluntary}) et \eng{to raise awareness} (jamais \eng{to sensitise}).
\item En écoute active, on note les mots-clés : noms propres, chiffres, actions, slogan.
\end{itemize}
\end{aretenir}

\newpage

\coursec{Activité d'intégration}

\courssub{Activité d'intégration : Mini-reportage radio « Community Voices »}

\courssub{Objectifs de l'activité}

À l'issue de cette activité d'intégration, l'élève sera capable de :
\begin{itemize}
\item \cle{réinvestir} les stratégies d'écoute étudiées dans la leçon (repérer le thème, les mots-clés, les chiffres et les noms propres) dans une situation de communication authentique ;
\item \cle{produire oralement} un mini-reportage radio de deux minutes sur un volontaire engagé pour l'égalité des sexes, en réemployant le lexique du volontariat (\eng{volunteer, to empower, awareness campaign, gender gap, role model}...) ;
\item \cle{écouter activement} les productions des autres groupes et remplir une grille de prise de notes (\eng{name, place, numbers, actions}) ;
\item \cle{formuler} des questions vrai/faux pertinentes et y répondre en justifiant ses choix ;
\item \cle{coopérer} en groupe de quatre en répartissant les rôles (présentateur, journaliste, volontaire, lecteur de l'appel à l'action).
\end{itemize}

\courssub{Situation}

\begin{textebox}[Community Voices — Call for contributions]
\textbf{Radio Buea 95.1 FM — “Community Voices”}

Are you a student with a story to tell? Our weekly programme “Community Voices” is looking for short radio reports about \textbf{local volunteers} who promote \textbf{gender equality} in our region.

Your report must last about \textbf{two minutes} and include:
an interview with the volunteer (name, age, village or town), the actions he or she carries out, at least \textbf{two figures} (number of people trained, meetings held, girls sent back to school...), a memorable \textbf{slogan}, and a short \textbf{call to join} the association.

The best report, chosen by the listeners, will be broadcast on Saturday at 6 p.m. Send us your voices — your community is listening!
\end{textebox}

Votre classe d'anglais a décidé de répondre à cet appel à contributions. Par groupes de quatre, vous allez préparer un mini-reportage radio sur un(e) volontaire \textbf{imaginaire} mais réaliste, engagé(e) pour l'égalité des sexes dans votre région (par exemple : une jeune femme de Bamenda qui forme des filles aux métiers du numérique, un jeune homme de Maroua qui sensibilise les chefs traditionnels à la scolarisation des filles, une retraitée de Douala qui aide des veuves à connaître leurs droits). Les autres groupes seront vos auditeurs : ils prendront des notes et répondront à vos questions.

\courssub{Consignes}

\begin{enumerate}
\item \textbf{Préparation en groupe (15 min).} Inventez votre volontaire : nom, âge, village ou quartier, association, actions concrètes, au moins deux chiffres, un slogan. Répartissez les rôles : un présentateur (\eng{host}), un journaliste (\eng{interviewer}), le/la volontaire (\eng{interviewee}), un lecteur de l'appel à rejoindre l'association (\eng{call to action}).
\item \textbf{Rédaction du script (15 min).} Écrivez le script en anglais correct : présentation courte, interview de cinq ou six questions-réponses, appel à l'action. Utilisez au moins \textbf{huit mots du lexique} de la leçon et le \eng{present perfect} pour les actions accomplies (\eng{She has trained fifty girls.}).
\item \textbf{Préparation des questions (5 min).} Rédigez \textbf{deux questions vrai/faux} sur votre reportage, dont une portant sur un chiffre, pour tester l'écoute de la classe.
\item \textbf{Enregistrement ou représentation (2 min par groupe).} Jouez votre reportage devant la classe. Parlez clairement, à voix haute, en articulant.
\item \textbf{Écoute active.} Pendant chaque reportage, remplissez la grille de prise de notes ci-dessous, puis répondez aux deux questions vrai/faux du groupe en justifiant votre réponse par une phrase (\eng{It is true because the volunteer said that...}).
\item \textbf{Vote final.} La classe vote pour le reportage le plus convaincant : clarté, exactitude de la langue, richesse du lexique, réalisme.
\end{enumerate}

\begin{center}
\begin{tabularx}{\linewidth}{|X|X|X|X|}
\hline
\rowcolor{popblueL} \textbf{Name and place} & \textbf{Numbers} & \textbf{Actions} & \textbf{Slogan} \\
\hline
 & & & \\
\hline
 & & & \\
\hline
 & & & \\
\hline
 & & & \\
\hline
\end{tabularx}
\end{center}

\courssub{Ressources à mobiliser}

\begin{aretenir}[Boîte à outils linguistique]
\begin{itemize}
\item \textbf{Lexique du volontariat et de l'égalité :} \eng{volunteer} (bénévole), \eng{to volunteer} (faire du bénévolat), \eng{to empower women} (autonomiser les femmes), \eng{awareness campaign} (campagne de sensibilisation), \eng{gender equality} (égalité des sexes), \eng{gender gap} (écart entre les sexes), \eng{role model} (modèle de rôle), \eng{to fight discrimination} (lutter contre la discrimination), \eng{community} (communauté), \eng{to raise funds} (collecter des fonds), \eng{gender-based violence} (violences basées sur le genre), \eng{advocate} (défenseur/se).
\item \textbf{Structures de l'interview :} \eng{Can you tell us about...? — What inspired you to...? — How many people have you trained? — What is your message to young girls? — What challenges do you face?}
\item \textbf{Grammaire :} \eng{present perfect} pour les bilans d'actions (\eng{She has organised twelve workshops.}) avec \eng{since} (date) ou \eng{for} (durée) ; \eng{present simple} pour les habitudes (\eng{Every Saturday, she visits schools.}) ; \eng{going to / will} pour les projets (\eng{We are going to open a new centre.}).
\item \textbf{Stratégies d'écoute :} avant d'écouter, lire la grille ; pendant l'écoute, noter d'abord les noms propres et les chiffres ; ne pas chercher à tout comprendre mot à mot.
\end{itemize}
\end{aretenir}

\begin{attention}[Pièges à éviter]
Ne dites pas \eng{“the equality of genders”} mais \eng{gender equality}. Attention à la prononciation de \eng{women} (« ouim-ine ») et \eng{woman} (« wou-mane »). N'oubliez pas la troisième personne au présent simple : \eng{she organises}, \eng{he works}. Enfin, \eng{volunteer} se prononce « vol-on-ti-euh » (accent sur la dernière syllabe), pas « vo-lan-tir ».
\end{attention}

\courssub{Proposition de production et corrigé commenté}

\begin{exoresolu}[Modèle de reportage : “Community Voices”, Groupe 3]
Énoncé : rédiger et jouer un reportage de deux minutes sur un volontaire imaginaire de la région du Nord-Ouest.
\tcblower
\textbf{Script modèle :}

\textbf{Host:} Good evening, dear listeners, and welcome to “Community Voices”. Today, our reporters take you to Bamenda, in the North-West Region, to meet a remarkable volunteer.

\textbf{Journalist:} Good afternoon, Madam. Can you introduce yourself to our listeners?

\textbf{Volunteer:} With pleasure. My name is \textbf{Grace Nfor}, I am twenty-eight years old, and I live in Nkwen, Bamenda. I have been a volunteer for the association “Girls Can Lead” for five years.

\textbf{Journalist:} What inspired you to fight for gender equality?

\textbf{Volunteer:} When I was fifteen, my parents wanted me to leave school and get married. A teacher stood up for me, and I finished my studies. I decided that no girl in my community should face the same situation alone.

\textbf{Journalist:} What actions have you carried out so far?

\textbf{Volunteer:} Since 2020, we have organised \textbf{forty awareness campaigns} in villages, and we have helped \textbf{one hundred and twenty girls} go back to school. We have also trained \textbf{sixty women} in sewing and small business management.

\textbf{Journalist:} What is your message to young girls?

\textbf{Volunteer:} Education is your right, not a favour. Our slogan is: \textbf{“Educate a girl, empower a nation!”}

\textbf{Call to action:} Do you want to make a difference? Join “Girls Can Lead” every Saturday at 10 a.m. at the Nkwen community hall. Together, we can close the gender gap!

\textbf{Host:} Thank you, Grace, and thank you for listening. See you next week on “Community Voices”!

\textbf{True or false questions prepared by the group:}
1. Grace has been a volunteer for ten years. (\emph{False — she said five years.})
2. The association has trained sixty women. (\emph{True.})

\textbf{Commentaire des choix de langue :}
\begin{itemize}
\item Le \eng{present perfect} (\eng{have organised, have helped, have trained}) est utilisé pour le bilan des actions, conformément à la règle : action passée dont le résultat compte aujourd'hui. L'association avec \eng{since 2020} et \eng{for five years} ancre la durée.
\item Les chiffres (\eng{forty, one hundred and twenty, sixty}) sont placés dans des phrases courtes pour être facilement repérés par les auditeurs — c'est la stratégie d'écoute des nombres mise en pratique par les émetteurs eux-mêmes.
\item Le slogan utilise l'impératif (\eng{Educate, empower}), forme typique des slogans et des appels à l'action.
\item L'appel final combine une question rhétorique (\eng{Do you want to make a difference?}) et des indications pratiques (jour, heure, lieu), comme dans un vrai message radio.
\end{itemize}
\end{exoresolu}

\courssub{Bilan de l'activité}

Cette activité a permis de boucler le cycle complet de la compétence : les élèves sont passés du statut d'\textbf{auditeurs} d'un texte sur le volontariat à celui de \textbf{producteurs} d'un message authentique, puis à nouveau d'auditeurs actifs grâce à la grille de notes et aux questions vrai/faux. Le lexique du thème (\eng{volunteer, to empower, gender equality, awareness campaign}) a été réemployé dans un contexte de communication réel et motivant, ancré dans la réalité camerounaise. Le vote final valorise la qualité de la langue et de la prise de parole, et prépare naturellement à l'épreuve orale du baccalauréat. Pour aller plus loin, les meilleurs scripts peuvent être réellement enregistrés sur téléphone portable et diffusés lors de la semaine des langues du lycée.

\newpage

\coursec{Méthodes et exercices résolus}

\courssub{Méthodes et exercices résolus}

\begin{textebox}[Script d'écoute : « Voices for Change » — Atelier de sensibilisation à Buea]
\textbf{Speaker 1 (Moderator):} Good morning everyone. Welcome to the “Voices for Change” workshop organized by the \cle{NGO} “Women’s Horizon” here in Buea. Today, we focus on \cle{volunteering} as a driver for \cle{gender equality}. Our first speaker is Mr. Nfor, the program coordinator.

\textbf{Speaker 2 (Mr. Nfor):} Thank you. Since 2021, we have recruited over \cle{three hundred} volunteers across the Southwest Region. \cle{Sixty percent} of them are young men who act as \cle{male allies}. They visit schools and markets to challenge harmful \cle{stereotypes}. Last month, we reached \cle{five thousand} students in \cle{twenty} secondary schools. Our goal is \cle{empowerment} through education.

\textbf{Speaker 3 (Achuo, Beneficiary):} I am Achuo, a Form 5 student. Before this program, I thought leadership was only for boys. The volunteers taught us about \cle{consent} and \cle{gender-based violence}. Now, I lead the \cle{Girls’ Club} in my school. We organize \cle{sensitization} walks every \cle{International Women’s Day}.

\textbf{Speaker 4 (James, Male Ally Volunteer):} As a \cle{volunteer}, I used to think \cle{feminism} was anti-men. The training changed my \cle{mindset}. Now I advocate for \cle{shared responsibilities} at home. My father was skeptical, but when he saw me cooking and helping my sisters with homework, he joined our \cle{community dialogue} sessions. Real change starts in the family.

\textbf{Moderator:} Thank you. Remember, volunteering is not charity; it is \cle{solidarity}. Join us next Saturday at the Buea Town Hall for the \cle{advocacy} march.
\end{textebox}

\begin{methode}[Identifier le thème et les mots-clés principaux (Compréhension globale)]
\textbf{Objectif :} Saisir le sujet général, le ton et les acteurs dès la première écoute.
\begin{enumerate}
    \item \textbf{Avant l'écoute :} Lire le titre (\eng{Voices for Change}) et les consignes. Prédire le vocabulaire : \eng{volunteering, gender, equality, workshop, NGO, stereotypes}.
    \item \textbf{Première écoute (global) :} Ne pas écrire. Identifier : \textbf{Qui parle ?} (Modérateur, Coordinateur, Bénéficiaire, Volontaire), \textbf{Où ?} (Buea), \textbf{Quel thème ?} (Volontariat / Égalité hommes-femmes).
    \item \textbf{Repérer les \cle{mots-clés} récurrents :} \eng{volunteering, gender equality, stereotypes, empowerment, male allies, advocacy, sensitization}. Les entourer sur le sujet d'examen.
    \item \textbf{Deuxième écoute :} Confirmer l'idée principale : \eng{“Volunteering engages men and women to fight stereotypes and promote gender equality in schools and communities.”}
\end{enumerate}
\textbf{Réflexe :} Un mot compris sur trois suffit pour valider l'hypothèse de thème. Ne pas bloquer sur un mot inconnu.
\end{methode}

\begin{exoresolu}[Examen type : Vrai ou Faux + Justification par citation]
\textbf{Consigne :} Écoutez le texte. Dites si les affirmations sont \textbf{True} ou \textbf{False}. Justifiez chaque réponse par une phrase exacte du texte.
\begin{enumerate}
    \item The NGO “Women’s Horizon” organized the workshop in Yaoundé.
    \item Mr. Nfor states that 60\% of the volunteers are young women.
    \item Achuo leads the Girls’ Club in her school.
    \item James’s father refused to participate in community dialogue sessions.
    \item The moderator invites listeners to an advocacy march next Saturday.
\end{enumerate}
\tcblower
\textbf{Solution détaillée :}
\begin{enumerate}
    \item \textbf{False.} \eng{“Welcome to the ‘Voices for Change’ workshop organized by the NGO ‘Women’s Horizon’ here in \textbf{Buea}.”} (Lieu : Buea, pas Yaoundé).
    \item \textbf{False.} \eng{“\textbf{Sixty percent} of them are young \textbf{men} who act as male allies.”} (Ce sont des hommes, pas des femmes).
    \item \textbf{True.} \eng{“Now, I lead the \textbf{Girls’ Club} in my school.”} (Citation directe d'Achuo).
    \item \textbf{False.} \eng{“My father was skeptical, but when he saw me cooking… he \textbf{joined our community dialogue sessions}.”} (Il a rejoint, pas refusé).
    \item \textbf{True.} \eng{“Join us next \textbf{Saturday} at the Buea Town Hall for the \textbf{advocacy march}.”} (Invitation explicite du modérateur).
\end{enumerate}
\textbf{Conclusion :} La justification doit être une \textbf{citation textuelle} (guillemets anglais “ ”), pas une reformulation. Copier exactement les mots entendus/lus.
\end{exoresolu}

\begin{methode}[Repérer les informations précises : chiffres, noms propres, dates (Compréhension détaillée)]
\textbf{Objectif :} Extraire des données factuelles (quantités, noms, temps, lieux) souvent pièges au Bac.
\begin{enumerate}
    \item \textbf{Préparer un tableau de prise de notes} avant la 2ème écoute :
    \begin{center}
    \begin{tabularx}{\linewidth}{|X|X|X|}\hline
    \rowcolor{popblueL} \textbf{Information cherchée} & \textbf{Mot-clé à écouter} & \textbf{Ma note} \\ \hline
    Date de début du programme & \eng{Since / In / Year} & 2021 \\ \hline
    Nombre total de volontaires & \eng{Number / Over / Recruited} & 300+ \\ \hline
    Nombre d'écoles atteintes & \eng{Schools / Reached} & 20 \\ \hline
    Nombre d'élèves touchés & \eng{Students / Thousand} & 5000 \\ \hline
    Événement à venir / Date & \eng{Next / Saturday / March} & Samedi prochain / Buea Town Hall \\ \hline
    \end{tabularx}
    \end{center}
    \item \textbf{Écoute active :} Se concentrer \textbf{uniquement} sur ces cases. Ignorer le reste.
    \item \textbf{Vigilance nombres :} \eng{thirteen (13) vs thirty (30)}, \eng{sixteen (16) vs sixty (60)}. Écrire en chiffres (300) pour éviter l'orthographe.
    \item \textbf{Noms propres :} Les épeler phonétiquement si nécessaire (\eng{N-f-o-r, A-c-h-u-o, B-u-e-a}).
\end{enumerate}
\textbf{Astuce :} Les chiffres sont souvent donnés \textbf{deux fois} (ex: \eng{“over three hundred”} puis \eng{“300 volunteers”}). La 2ème confirmation valide la note.
\end{methode}

\begin{exoresolu}[Examen type : QCM et Réponses courtes (Informations factuelles)]
\textbf{Consigne :} Répondez aux questions en vous basant sur le texte.
\begin{enumerate}
    \item Since when has the NGO recruited volunteers? \quad A) 2020 \quad B) 2021 \quad C) 2022
    \item How many secondary schools were reached last month? \quad A) 12 \quad B) 20 \quad C) 50
    \item What is Achuo’s current role? \quad A) Volunteer \quad B) Program Coordinator \quad C) Girls’ Club Leader
    \item Where will the advocacy march take place? \quad A) Buea Market \quad B) Buea Town Hall \quad C) Women’s Horizon Office
    \item Complete: James advocates for \_\_\_\_\_\_\_\_\_\_\_\_ at home. (2 mots)
\end{enumerate}
\tcblower
\textbf{Solution détaillée :}
\begin{enumerate}
    \item \textbf{B) 2021.} \eng{“Since \textbf{2021}, we have recruited…”} (Présent perfect + \eng{since} = point de départ).
    \item \textbf{B) 20.} \eng{“we reached \textbf{five thousand} students in \textbf{twenty} secondary schools.”}
    \item \textbf{C) Girls’ Club Leader.} \eng{“Now, I \textbf{lead the Girls’ Club} in my school.”}
    \item \textbf{B) Buea Town Hall.} \eng{“Join us next Saturday at the \textbf{Buea Town Hall}…”}
    \item \textbf{shared responsibilities.} \eng{“Now I advocate for \textbf{shared responsibilities} at home.”} (Collocation forte : \eng{advocate for + nom}).
\end{enumerate}
\textbf{Justifications grammaticales :} Q1 : \eng{Since} + date → \eng{Present Perfect}. Q5 : \eng{Advocate for} (verbe + préposition) suivi d'un nom pluriel. Attention à l'orthographe : \eng{responsibilities} (pluriel en \eng{-ies}).
\end{exoresolu}

\begin{methode}[Mobiliser le lexique thématique : collocations, familles de mots, faux amis]
\textbf{Objectif :} Maîtriser le vocabulaire spécifique pour le réemploi (Speaking/Writing) et éviter les pièges francophones.
\begin{enumerate}
    \item \textbf{Regrouper par champs lexicaux (Mindmap mental) :}
    \begin{itemize}
        \item \textbf{Acteurs :} \eng{volunteer (n/v), coordinator, beneficiary, male ally, advocate (n/v)}.
        \item \textbf{Actions :} \eng{recruit, challenge (stereotypes), empower, sensitize, advocate for, join, lead}.
        \item \textbf{Concepts clés :} \eng{gender equality, gender equity, stereotypes, gender-based violence (GBV), consent, empowerment, mindset, solidarity, shared responsibilities}.
        \item \textbf{Événements :} \eng{workshop, sensitization walk, advocacy march, community dialogue, International Women’s Day}.
    \end{itemize}
    \item \textbf{Travailler les \cle{collocations} (mots qui vont ensemble) :}
    \begin{center}
    \begin{tabular}{|l|l|}\hline
    \rowcolor{popblueL} \textbf{Verbe + Nom / Préposition} & \textbf{Traduction} \\ \hline
    \eng{challenge stereotypes} & remettre en question les stéréotypes \\ \hline
    \eng{advocate \textbf{for} equality} & plaider pour / militer pour l'égalité \\ \hline
    \eng{raise \textbf{awareness} / \textbf{sensitize} people \textbf{on/about}} & sensibiliser les gens à \\ \hline
    \eng{empower women / girls} & autonomiser / donner du pouvoir aux femmes \\ \hline
    \eng{join \textbf{in} / take \textbf{part in}} & participer à \\ \hline
    \end{tabular}
    \end{center}
    \item \textbf{Familles de mots (Word formation) — \cle{Bac écrit} :}
    \begin{center}
    \begin{tabular}{|l|l|l|l|}\hline
    \rowcolor{popblueL} \textbf{Verbe} & \textbf{Nom (action/état)} & \textbf{Nom (agent)} & \textbf{Adjectif} \\ \hline
    \eng{volunteer} & \eng{volunteering} & \eng{volunteer} & \eng{voluntary} \\ \hline
    \eng{empower} & \eng{empowerment} & — & \eng{empowered} \\ \hline
    \eng{sensitize} & \eng{sensitization} & — & \eng{sensitive} (\eng{$\neq$ sensible}) \\ \hline
    \eng{advocate} & \eng{advocacy} & \eng{advocate} & — \\ \hline
    \eng{coordinate} & \eng{coordination} & \eng{coordinator} & \eng{coordinated} \\ \hline
    \end{tabular}
    \end{center}
    \item \textbf{\cle{Faux amis} à neutraliser :}
    \begin{itemize}
        \item \eng{Sensible} = \textbf{raisonnable / sensible (person)}. \eng{Sensitive} = \textbf{sensible (émotif) / délicat (sujet)}. \eng{Sensitize} = \textbf{sensibiliser} (pas \eng{sensibilize}).
        \item \eng{Gender} = \textbf{genre (social)}. \eng{Sex} = \textbf{sexe (biologique)}. Ne pas confondre.
        \item \eng{Equity} = \textbf{équité (justice, donner selon les besoins)}. \eng{Equality} = \textbf{égalité (mêmes droits/traitement)}.
    \end{itemize}
\end{enumerate}
\end{methode}

\begin{exoresolu}[Examen type : Complétion, Association, Traduction]
\textbf{Consigne :}
\textbf{A.} Complétez avec la forme correcte du mot entre parenthèses.
\begin{enumerate}
    \item The \_\_\_\_\_\_\_\_\_\_\_\_ (coordinate) of the program spoke first.
    \item We need more \_\_\_\_\_\_\_\_\_\_\_\_ (volunteer) for the march.
    \item Her \_\_\_\_\_\_\_\_\_\_\_\_ (empower) is visible in her confidence.
    \item The \_\_\_\_\_\_\_\_\_\_\_\_ (sensitize) campaign targets rural areas.
\end{enumerate}
\textbf{B.} Associez le verbe à sa préposition correcte.
\begin{center}
\begin{tabular}{ll}
1. Advocate & A. on / about \\
2. Sensitize & B. for \\
3. Challenge & C. (none) \\
4. Join & D. in \\
\end{tabular}
\end{center}
\textbf{C.} Traduisez en anglais : « Les volontaires masculins luttent contre les stéréotypes de genre. »
\tcblower
\textbf{Solution détaillée :}
\textbf{A. Word Formation}
\begin{enumerate}
    \item \textbf{coordinator} (Nom agent : \eng{-or} après \eng{coordin-}). \eng{“The \textbf{coordinator} of the program…”}
    \item \textbf{volunteers} (Nom pluriel : \eng{volunteer} ne change pas, ajoute \eng{s}). \eng{“We need more \textbf{volunteers}…”}
    \item \textbf{empowerment} (Nom abstrait en \eng{-ment}). \eng{“Her \textbf{empowerment} is visible…”}
    \item \textbf{sensitization} (Nom action en \eng{-tion}). Attention : \eng{sensitize} $\to$ \eng{sensitization} (pas \eng{sensibilization}). \eng{“The \textbf{sensitization} campaign…”}
\end{enumerate}
\textbf{B. Collocations (Verbe + Préposition)}
1–\textbf{B} (\eng{advocate \textbf{for}}). 2–\textbf{A} (\eng{sensitize \textbf{on/about}}). 3–\textbf{C} (\eng{challenge} \textbf{transitif direct}, pas de préposition : \eng{challenge stereotypes}). 4–\textbf{D} (\eng{join \textbf{in}} une activité / \eng{join} un groupe).
\textbf{C. Traduction}
\eng{“Male volunteers \textbf{challenge gender stereotypes}.”} ou \eng{“Male volunteers \textbf{fight against gender stereotypes}.”}
\textbf{Pièges évités :} \eng{Male} (pas \eng{masculine} pour humains), \eng{challenge} (pas \eng{fight} qui demande \eng{against}), \eng{gender stereotypes} (ordre adjectif + nom, \eng{gender} adjectif classificateur).
\end{exoresolu}

\begin{methode}[Réemployer le lexique dans une production guidée (Résumé oral / Écrit court)]
\textbf{Objectif :} Produire un résumé structuré (80–100 mots) ou une prise de parole courte (1–2 min) en réutilisant le vocabulaire ciblé.
\begin{enumerate}
    \item \textbf{Structure type « Résumé de document audio » :}
    \begin{itemize}
        \item \textbf{Introduction (1 phrase) :} Source + Thème + Thèse.
        \eng{“The audio document ‘Voices for Change’ reports on a workshop in Buea highlighting the role of volunteering in promoting gender equality.”}
        \item \textbf{Développement (3–4 phrases) :} Acteurs + Actions + Chiffres clés + Témoignages. Utiliser \textbf{connecteurs logiques} : \eng{Firstly, Moreover, For instance, As a result, Finally}.
        \eng{“Firstly, Mr Nfor, the coordinator, explains that since 2021, over 300 volunteers, mostly male allies, have challenged stereotypes in 20 schools. Moreover, Achuo, a student, testifies that she now leads a Girls’ Club. Furthermore, James, a volunteer, describes how his mindset changed and influenced his father.”}
        \item \textbf{Conclusion (1 phrase) :} Portée / Perspective / Appel à l'action.
        \eng{“Finally, the moderator calls for solidarity and invites everyone to an advocacy march next Saturday.”}
    \end{itemize}
    \item \textbf{Contrôle :} Compter les mots. Vérifier : \cle{Passif} (\eng{volunteers \textbf{are recruited}}), \cle{Discours rapporté} (\eng{He said \textbf{that} they \textbf{had reached}}), \cle{Modaux} (\eng{We \textbf{must/should} join}).
    \item \textbf{À l'oral (Speaking) :} S'enregistrer. Vérifier la prononciation des mots-clés : \eng{volunteering} (« vol-un-ti-ring »), \eng{gender} (« jen-der »), \eng{advocacy} (« ad-vo-ka-si »), \eng{empowerment} (« em-pau-eu-ment »).
\end{enumerate}
\end{methode}

\begin{exoresolu}[Examen type : Production écrite courte (Résumé structuré)]
\textbf{Consigne :} Rédigez un résumé du document audio en \textbf{80 à 100 mots}. Utilisez au moins \textbf{6 mots ou expressions} du vocabulaire de la leçon (surlignés dans votre copie).
\tcblower
\textbf{Proposition de réponse (92 mots) :}
“The audio document \textbf{“Voices for Change”} deals with a \textbf{workshop} held in Buea by the \textbf{NGO} “Women’s Horizon”. It highlights \textbf{volunteering} as a key tool for \textbf{gender equality}. The \textbf{coordinator}, Mr Nfor, states that since 2021, over 300 \textbf{volunteers}, including many \textbf{male allies}, have \textbf{challenged stereotypes} in twenty schools, reaching 5,000 students. \textbf{For instance}, Achuo, a \textbf{beneficiary}, now \textbf{leads} a \textbf{Girls’ Club} and organizes \textbf{sensitization} walks. James, a \textbf{volunteer}, explains how the training changed his \textbf{mindset} regarding \textbf{shared responsibilities}, convincing his father to join \textbf{community dialogues}. The moderator concludes by calling for \textbf{solidarity} and announcing an \textbf{advocacy march} next Saturday.”
\textbf{Analyse :} 16 mots-clés réemployés (en gras). Structure Intro/Développement/Conclusion respectée. Connecteurs : \eng{It highlights, states that, including, for instance, explains how, concludes by}. Grammaire : Passif (\eng{is held}), Discours rapporté (\eng{states that... have challenged}), Présent perfect (\eng{has changed}).
\end{exoresolu}

\begin{methode}[Identifier les structures grammaticales clés à l'oral : Passif, Discours rapporté, Modaux]
\textbf{Objectif :} Reconnaître et transformer les structures fréquentes dans les discours institutionnels / associatifs (reporting, recommendations).
\begin{enumerate}
    \item \textbf{Voix Passive (Passive Voice) — Focus sur l'action, pas l'agent :}
    \eng{Active: “We \textbf{have recruited} 300 volunteers.”} $\to$ \eng{Passive: “300 volunteers \textbf{have been recruited}.”}
    \eng{Active: “They \textbf{reached} 5000 students.”} $\to$ \eng{Passive: “5000 students \textbf{were reached}.”}
    \textbf{Règle :} \eng{Be (temps du verbe actif) + Participe Passé (V3)}. Agent introduit par \eng{by} si nécessaire.
    \item \textbf{Discours Rapporté (Reported Speech) — Résumer ce qu'ont dit les intervenants :}
    \begin{itemize}
        \item \eng{Direct: “I \textbf{lead} the Girls’ Club.” (Achuo)} $\to$ \eng{Reported: She \textbf{said (that)} she \textbf{led} the Girls’ Club.} (Backshift : Present $\to$ Preterite).
        \item \eng{Direct: “We \textbf{have recruited} volunteers.” (Mr Nfor)} $\to$ \eng{Reported: He \textbf{stated (that)} they \textbf{had recruited} volunteers.} (Present Perfect $\to$ Past Perfect).
        \item \eng{Direct: “Join us next Saturday.” (Moderator)} $\to$ \eng{Reported: He \textbf{urged / invited} us \textbf{to join} them \textbf{the following Saturday}.} (Impératif $\to$ Infinitif ; \eng{next} $\to$ \eng{the following}).
    \end{itemize}
    \item \textbf{Modaux de recommandation / obligation (Speaking/Writing) :}
    \eng{We \textbf{must / have to / should} challenge stereotypes.} (Obligation forte / conseil).
    \eng{Men \textbf{can / could} be allies.} (Capacité / Possibilité).
    \eng{Volunteers \textbf{need to / ought to} sensitize communities.} (Nécessité).
\end{enumerate}
\textbf{Réflexe Bac :} Transformer une phrase active en passive (sujet = chose/personne affectée). Passer du direct au rapporté (changements de temps, pronoms, indicateurs spatio-temporels).
\end{methode}

\begin{exoresolu}[Examen type : Transformation de phrases (Grammaire en contexte)]
\textbf{Consigne :} Réécrivez les phrases comme indiqué sans en changer le sens.
\begin{enumerate}
    \item Active to Passive: “The NGO \textbf{organizes} workshops every year.” $\to$ Workshops \rule{3cm}{0.4pt} every year by the NGO.
    \item Direct to Reported: “We \textbf{are challenging} harmful stereotypes now,” Mr Nfor said. $\to$ Mr Nfor said that they \rule{3cm}{0.4pt} harmful stereotypes \rule{3cm}{0.4pt}.
    \item Direct to Reported (Imperative): “\textbf{Join} the advocacy march,” the moderator told the crowd. $\to$ The moderator \rule{3cm}{0.4pt} the crowd \rule{3cm}{0.4pt} the advocacy march.
    \item Modal choice: “It is necessary for volunteers to understand consent.” $\to$ Volunteers \rule{3cm}{0.4pt} understand consent.
    \item Word Order (French interference): “They fight for gender equality.” (Traduire en français puis ré-écrire en anglais si nécessaire — ici vérifier l'ordre). $\to$ \textit{Ordre correct SVOMPT respecté.}
\end{enumerate}
\tcblower
\textbf{Solution détaillée :}
\begin{enumerate}
    \item \textbf{are organized} (Present Simple Passive : \eng{am/is/are + V3}). Sujet pluriel \eng{Workshops} $\to$ \eng{are}.
    \item \textbf{were challenging / then} (ou \textbf{at that time}). Backshift : Present Continuous (\eng{are challenging}) $\to$ Past Continuous (\eng{were challenging}). \eng{now} $\to$ \eng{then / at that time}.
    \item \textbf{urged / invited / encouraged / asked} \textbf{to join}. Verbe de parole + Objet + Infinitif. \eng{urged} (fort), \eng{invited} (neutre), \eng{asked} (standard).
    \item \textbf{must / have to / need to / ought to}. \eng{It is necessary} $\leftrightarrow$ \eng{must / have to}. \eng{Need to} (besoin objectif). \eng{Ought to} (recommandation morale).
    \item \textbf{They fight for gender equality.} Ordre anglais strict \textbf{Sujet + Verbe + Complément (SVO)}. Pas de préposition devant \eng{gender equality} après \eng{fight for}. \textbf{Erreur fréquente FR :} *“They fight for the gender equality”* (pas de \eng{the} devant concept abstrait non défini).
\end{enumerate}
\textbf{Conclusion :} La maîtrise du \textbf{backshift} (recul des temps) et de la \textbf{passive} est indispensable pour le \textbf{résumé} (Writing) et le \textbf{compte-rendu} (Speaking).
\end{exoresolu}

\begin{attention}[Les 5 erreurs qui coûtent des points au Bac]
\begin{enumerate}
    \item \textbf{Faux ami \eng{Sensible / Sensitive} :} Écrire \eng{“We must be sensible about gender issues.”} $\to$ \eng{Sensible} = \textbf{raisonnable}. Correct : \eng{“We must be \textbf{sensitive} about gender issues.”} (\textbf{délicat, qui demande du tact}). \eng{To sensitize} = \textbf{sensibiliser}.
    \item \textbf{Article défini devant concepts abstraits :} \eng{“The gender equality is important.”} $\to$ \textbf{Zéro article} : \eng{“Gender equality is important.”} (\eng{The} seulement si défini : \eng{The gender equality \textbf{we promote}}).
    \item \textbf{Ordre des mots (Calque FR) :} \eng{“Volunteers male challenge stereotypes gender.”} $\to$ \textbf{Adjectif AVANT nom} : \eng{Male volunteers challenge gender stereotypes.} (\eng{Gender} = adjectif classificateur, invariable).
    \item \textbf{Préposition après \eng{Advocate} :} \eng{“They advocate gender equality.”} $\to$ \textbf{Advocate FOR} + nom. \eng{“They advocate \textbf{for} gender equality.”} (Mais \eng{advocate} verbe transitif direct au sens juridique : \eng{He advocated the cause}).
    \item \textbf{Concordance des temps (Backshift) oublié :} \eng{“He said they \textbf{have} recruited 300 volunteers.”} $\to$ Après \eng{said} (passé) : \eng{they \textbf{had} recruited} (Past Perfect). Règle : Present Perfect $\to$ Past Perfect.
\end{enumerate}
\end{attention}

\newpage

\coursec{Exercices}

\coursec{Lesson 37 — Listening: Texts on volunteering in gender equality promotion activities}

\courssub{Warm-up: Entering the theme}

\begin{experience}[Brainstorming \& Prediction]
\emph{Étape 1 — Mots-clés.} Au tableau, l'enseignant écrit : \eng{volunteer}, \eng{gender equality}, \eng{workshop}, \eng{campaign}, \eng{empowerment}. En binômes, les élèves associent chaque mot à une définition courte en anglais (ex. \eng{“A person who works for free”}).
\par\emph{Étape 2 — Contexte.} L'enseignant projette ou décrit une image : un groupe de jeunes (filles et garçons) parlant à des parents dans un village camerounais. Question orale : \eng{“What are they doing? Why is it important?”}.
\par\emph{Étape 3 — Prédiction.} Avant d'écouter le premier script, lire le titre : \eng{“A girls' education campaign in the North West”}. Deviner : lieu, acteurs, chiffre clé, objectif.
\end{experience}

\courssub{Listening Scripts (Support Texts)}

\begin{textebox}[Listening script 1 — A girls' education campaign in the North West]
Good morning, listeners. Today we report on “Girls Back to School”, a campaign launched in Bamenda by the association Women in Action. Since January, forty-five volunteers have visited sixty villages in the region. Their goal is to convince parents to enrol their daughters in secondary school. So far, the campaign has helped about 1,200 girls return to the classroom. The project leader, Mrs. Agnes Tabi, says that volunteers organise two workshops every week. The campaign will continue until December, and the association hopes to reach three thousand girls by the end of the year.
\end{textebox}

\begin{textebox}[Listening script 2 — A volunteer in Ghana]
My name is Kofi Mensah and I am a volunteer in Accra, Ghana. I joined the organisation Equal Future two years ago. Every Saturday, I run workshops for teenage boys about gender equality. We discuss why girls should stay in school and why domestic work should be shared at home. At first, only ten boys attended my workshops, but today there are more than fifty. Last month, our organisation also trained thirty women to become community leaders. I am proud because volunteering has changed my own vision of women's rights.
\end{textebox}

\courssub{Listening Strategies}

\begin{methode}[How to approach a listening text on volunteering]
\begin{enumerate}
\item \textbf{Before listening (Prédire) :} Lire le titre, les questions et les mots-clés du vocabulaire. Identifier le thème (\eng{volunteering}, \eng{gender equality}) et le type de document (reportage, témoignage).
\item \textbf{First listening (Compréhension globale) :} Repérer le \cle{thème principal}, le \cle{ton} (formel, personnel), les \cle{noms propres} (lieux, personnes, associations) et les \cle{chiffres} (dates, quantités).
\item \textbf{Second listening (Compréhension détaillée) :} Noter les informations précises pour répondre aux questions (V/F, QCM, réponses courtes). Écouter les \cle{mots de liaison} (\eng{however, because, so far, last month}) qui structurent le discours.
\item \textbf{After listening (Vérifier) :} Relever les citations exactes pour justifier les réponses (V/F justifié). Vérifier l'orthographe des noms propres et des chiffres.
\end{enumerate}
\end{methode}

\begin{exemplebox}[Worked example — Script 1]
\emph{Question :} “How many villages have been visited?”
\par\emph{Stratégie :} Écouter le chiffre + le mot \eng{villages}. \eng{“forty-five volunteers have visited sixty villages”}.
\par\emph{Réponse :} \eng{Sixty villages (60)}.
\par\emph{Piège :} Ne pas confondre \eng{forty-five} (volunteers) et \eng{sixty} (villages). Entraînez-vous à discriminer \eng{16 / 60} (voir exercice phonétique).
\end{exemplebox}

\courssub{Thematic Vocabulary}

\begin{aretenir}[Lexique essentiel : Volunteering \& Gender Equality]
\begin{center}
\begin{tabularx}{\linewidth}{|l|X|l|}
\hline
\rowcolor{popblueL} \textbf{English} & \textbf{Français} & \textbf{Exemple / Collocation} \\
\hline
\eng{volunteer} (n.) & bénévole, volontaire & \eng{a volunteer / to volunteer} \\
\eng{volunteering} (n.) & volontariat & \eng{volunteering activities} \\
\eng{campaign} (n.) & campagne & \eng{launch a campaign / awareness campaign} \\
\eng{workshop} (n.) & atelier & \eng{run / organise a workshop} \\
\eng{awareness} (n.) & sensibilisation, conscience & \eng{raise awareness (about)} \\
\eng{to enrol} (v.) & inscrire (à l'école) & \eng{enrol girls in school} (UK spelling) \\
\eng{empowerment} (n.) & autonomisation, émancipation & \eng{women's empowerment} \\
\eng{discrimination} (n.) & discrimination & \eng{gender discrimination / fight discrimination} \\
\eng{gender equality} (n.) & égalité des sexes & \eng{promote gender equality} \\
\eng{community leader} (n.) & leader communautaire & \eng{train community leaders} \\
\eng{sensitise} (v.) & sensibiliser & \eng{sensitise parents (UK)} \\
\eng{beneficiary} (n.) & bénéficiaire & \eng{the beneficiaries of the project} \\
\hline
\end{tabularx}
\end{center}
\emph{Note :} L'orthographe britannique est de rigueur (\eng{organise, enrol, sensitise}). \eng{Enrol} $\rightarrow$ \eng{enrolled, enrolling} (un seul \eng{l} en UK, double \eng{ll} en US).
\end{aretenir}

\courssub{Je vérifie mes connaissances}

\begin{exercice}[QCM de compréhension orale — Script 1]
Votre enseignant vous lit (ou vous fait écouter) l'enregistrement \textbf{Script 1} ci-dessus. Écoutez deux fois, puis choisissez la bonne réponse (a, b ou c) pour chaque question.

\begin{enumerate}
\item The campaign takes place in: \quad a) Douala \quad b) Bamenda \quad c) Buea
\item The number of volunteers is: \quad a) 15 \quad b) 40 \quad c) 45
\item The volunteers have visited: \quad a) 16 villages \quad b) 60 villages \quad c) 600 villages
\item About how many girls have returned to school? \quad a) 1,200 \quad b) 2,100 \quad c) 12,000
\item The project leader is called: \quad a) Mrs. Agnes Tabi \quad b) Mrs. Agnes Tambe \quad c) Miss Agnes Tabi
\end{enumerate}
\end{exercice}

\begin{exercice}[Vrai ou faux justifié — Script 2]
Écoutez le \textbf{Script 2} ci-dessus (lu par l'enseignant), puis dites si chaque affirmation est \eng{true} ou \eng{false}. Justifiez chaque réponse par une courte citation du texte en anglais.

\begin{enumerate}
\item Kofi has been a volunteer for three years.
\item His workshops are for teenage girls.
\item More than fifty boys now attend the workshops.
\item Thirty women were trained as community leaders.
\item Volunteering has not changed Kofi's personal opinions.
\end{enumerate}
\end{exercice}

\begin{exercice}[Discrimination phonétique : les nombres en \eng{-teen} vs \eng{-ty}]
Votre enseignant lit huit phrases. Dans chaque phrase, entourez le nombre que vous entendez.

\begin{enumerate}
\item 13 / 30
\item 14 / 40
\item 15 / 50
\item 16 / 60
\item 17 / 70
\item 18 / 80
\item 19 / 90
\item 15 / 50
\end{enumerate}

\emph{Rappel : les nombres en \eng{-teen} portent l'accent tonique sur la dernière syllabe (ex. \eng{fif-TEEN} « fi-f-tine »), les nombres en \eng{-ty} sur la première (ex. \eng{FIF-ty} « fif-ti »). Le \eng{/t/} de \eng{-ty} est souvent une flap \eng{[ɾ]} (son \eng{d} rapide) en anglais américain, mais reste \eng{[t]} en anglais britannique.}
\end{exercice}

\begin{exercice}[Phrases à trous : le lexique du thème]
Complétez les phrases avec les mots de la banque suivante (chaque mot ne sert qu'une fois) : \eng{volunteer} — \eng{awareness} — \eng{workshop} — \eng{to enrol} — \eng{empowerment} — \eng{discrimination}.

\begin{enumerate}
\item The NGO organised a \dots\ on women's rights in our school last Friday.
\item Many parents refuse \dots\ their daughters in secondary school.
\item A \dots\ is a person who works for free to help a community.
\item The campaign raises \dots\ about the importance of girls' education.
\item Gender \dots\ is still a serious problem in some workplaces.
\item Education is the key to women's \dots\ in our society.
\end{enumerate}
\end{exercice}

\courssub{Je m'entraîne}

\begin{exercice}[Traduction vers l'anglais]
Traduisez les phrases suivantes en anglais, en utilisant le lexique du thème.

\begin{enumerate}
\item Les volontaires ont organisé trois ateliers dans notre village.
\item Beaucoup de filles ne vont pas à l'école à cause de la pauvreté.
\item Cette association lutte contre la discrimination envers les femmes depuis 2015.
\item Nous devons sensibiliser les parents à l'importance de l'éducation des filles.
\item Le volontariat a changé ma vision de l'égalité des sexes.
\end{enumerate}
\end{exercice}

\begin{exercice}[Transformation de phrases]
Réécrivez chaque phrase en suivant la consigne entre parenthèses, sans changer le sens.

\begin{enumerate}
\item The volunteers have visited sixty villages. (mettre à la forme négative)
\item Mrs. Tabi organises two workshops every week. (poser la question sur \eng{two workshops} avec \eng{how many})
\item “We will train thirty women next month,” the leader said. (mettre au discours indirect avec \eng{the leader said that})
\item People should respect women's rights. (mettre à la voix passive : \eng{Women's rights\dots})
\item The campaign started in January. It is still active today. (relier avec \eng{since} au present perfect)
\end{enumerate}
\end{exercice}

\begin{exercice}[Texte à trous]
Complétez le texte suivant avec les mots de la banque : \eng{gender equality} — \eng{volunteers} — \eng{community} — \eng{rights} — \eng{school} — \eng{awareness} — \eng{training} — \eng{girls}.

\begin{textebox}[Volunteering in Yaoundé]
In Yaoundé, a group of young (1) \dots\ works every weekend to promote (2) \dots\ . They visit families and explain that (3) \dots\ have the same (4) \dots\ as boys. Thanks to their (5) \dots\ campaign, many parents now send their daughters to (6) \dots\ . The group also offers (7) \dots\ sessions to women who want to become leaders in their (8) \dots\ .
\end{textebox}
\end{exercice}

\begin{exercice}[Appariements : mots et définitions]
Associez chaque mot anglais (1--8) à sa définition (a--h).

\begin{center}
\begin{tabularx}{\linewidth}{|l|X|}
\hline
\rowcolor{popblueL} \textbf{Word} & \textbf{Definition} \\
\hline
1. volunteer & a) to make people conscious of a problem \\
\hline
2. awareness & b) a meeting where people learn and practise skills \\
\hline
3. workshop & c) a person who works without payment to help others \\
\hline
4. to enrol & d) unfair treatment because of sex, race or origin \\
\hline
5. empowerment & e) to officially register someone in a school \\
\hline
6. discrimination & f) knowledge or understanding of a social issue \\
\hline
7. to raise awareness & g) the process of giving people power and confidence \\
\hline
8. campaign & h) a series of organised actions to achieve a goal \\
\hline
\end{tabularx}
\end{center}
\end{exercice}

\courssub{Je me prépare au Bac}

\begin{exercice}[Compréhension de l'écrit et questions de langue — type Bac]
Lisez le texte suivant, puis répondez aux questions.

\begin{textebox}[Reading text — Volunteering for gender equality in Cameroon]
In many parts of Cameroon, girls still leave school earlier than boys. Poverty, early marriages and traditional beliefs often prevent them from completing their education. To fight this situation, several organisations now rely on young volunteers.

One of them is “Bright Future”, an association created in Buea in 2018. Its fifty volunteers, most of them university students, travel to rural areas to sensitise parents about the importance of girls' education. They organise workshops in schools, churches and community halls. During these sessions, they explain that educating a girl benefits the whole family: educated women earn better salaries, take better care of their children's health, and send their own daughters to school.

The results are encouraging. Since its creation, “Bright Future” has helped more than two thousand girls return to school. The association has also trained eighty women as community leaders, and it has opened three reading clubs for teenage girls.

However, the work is not always easy. “Some fathers refuse to listen to us,” says Clarisse, a twenty-two-year-old volunteer. “But when we come back again and again, attitudes slowly change.” For Clarisse, volunteering is more than an activity: “It has taught me patience, courage and public speaking. I have become a different person.”
\end{textebox}

\textbf{A. Comprehension questions} — Répondez en anglais, par des phrases complètes.

\begin{enumerate}
\item Why do many girls in Cameroon leave school early? Give two reasons.
\item Who are most of the volunteers of “Bright Future”?
\item Where do the volunteers organise their workshops? Mention two places.
\item According to the text, what are two benefits of educating a girl?
\item What has volunteering taught Clarisse?
\end{enumerate}

\textbf{B. True or false} — Justifiez chaque réponse par une citation du texte.

\begin{enumerate}
\item “Bright Future” was created in Douala.
\item The association has trained eighty women as community leaders.
\item All fathers immediately accept the volunteers' message.
\end{enumerate}

\textbf{C. Language work}

\begin{enumerate}
\item Trouvez dans le texte l'équivalent de : \emph{sensibiliser} (paragraphe 2) et \emph{encourageants} (paragraphe 3).
\item Mettez à la voix passive : \eng{The association has opened three reading clubs.}
\item Mettez au discours indirect : \eng{“Some fathers refuse to listen to us,” says Clarisse.}
\end{enumerate}
\end{exercice}

\begin{exercice}[Traduction — type Bac]
Traduisez en anglais le passage suivant.

\begin{textebox}[Texte à traduire]
Depuis 2020, une association camerounaise envoie des volontaires dans les villages pour promouvoir l'égalité des sexes. Ces jeunes organisent des ateliers et expliquent aux parents que les filles ont les mêmes droits que les garçons. Grâce à leur travail, plus de mille filles sont retournées à l'école. Le volontariat a aussi transformé les volontaires eux-mêmes : ils sont devenus plus confiants et plus responsables.
\end{textebox}
\end{exercice}

\begin{exercice}[Composition — type Bac]
Choisissez \textbf{un} des deux sujets suivants et rédigez une composition en anglais de \textbf{150 à 180 mots}.

\textbf{Sujet 1 — Article pour le journal du lycée.} Your school magazine has asked you to write an article calling on students to volunteer for gender equality. Describe the problem, present the actions volunteers can do, and encourage your schoolmates to join. Use at least five words from the lesson's vocabulary (\eng{volunteer, awareness, workshop, empowerment, campaign, to enrol, discrimination}) and at least one figure.

\textbf{Sujet 2 — Essai argumentatif.} “Boys should also be involved in the promotion of gender equality.” Do you agree with this statement? Give your opinion and support it with arguments and examples from your community or your country.

\emph{Conseils : faites un plan (introduction, deux ou trois paragraphes, conclusion) ; utilisez des connecteurs (\eng{first, moreover, however, in conclusion}) ; vérifiez les temps verbaux et l'orthographe.}
\end{exercice}

\begin{savaistu}[Le volontariat au Cameroun et dans le Commonwealth]
Le Cameroun, pays bilingue (français/anglais), partage une forte culture du volontariat avec les pays du Commonwealth (Nigeria, Ghana, Kenya, Afrique du Sud). Le \eng{National Youth Council} et le \eng{Ministry of Youth Affairs and Civic Education (MINJEC)} encadrent de nombreux programmes de \eng{volunteering} (ex. \eng{Programme National de Volontariat}). Au Royaume-Uni, le \eng{National Citizen Service (NCS)} propose aux 16--17 ans des missions d'intérêt général. Au Canada, le \eng{Canada Service Corps} finance des projets jeunesse. Ces expériences valorisent le \eng{CV} et développent les \eng{soft skills} (leadership, communication, \eng{teamwork}) essentielles pour l'insertion professionnelle.
\end{savaistu}

\begin{attention}[Erreurs fréquentes des francophones — Pièges à éviter]
\begin{itemize}
\item \textbf{Enrol vs Enroll :} UK = \eng{enrol, enrolled, enrolling} (1 \eng{l}) ; US = \eng{enroll, enrolled, enrolling} (2 \eng{l}). Au Cameroun : \textbf{UK}.
\item \textbf{Sensibiliser :} Ne pas dire \eng{“sensibilize”} (anglicisme). Verbe UK : \eng{sensitise} ; US : \eng{sensitize}. Nom : \eng{awareness} (\eng{raise awareness}).
\item \textbf{Discrimination :} \eng{Discrimination against} qqn (pas \eng{discrimination towards}).
\item \textbf{Chiffres :} \eng{1,200} = \eng{one thousand two hundred} (virgule en anglais, point en français). \eng{3,000} = \eng{three thousand}.
\item \textbf{Since / For :} \eng{Since 2020} (point de départ) vs \eng{For two years} (durée). Avec \eng{since} $\rightarrow$ Present Perfect (\eng{has been}).
\item \textbf{Discours indirect :} \eng{“I will go”} $\rightarrow$ \eng{He said he would go} (changement \eng{will} $\rightarrow$ \eng{would}, \eng{I} $\rightarrow$ \eng{he}, \eng{now} $\rightarrow$ \eng{then}, \eng{next month} $\rightarrow$ \eng{the following month}).
\end{itemize}
\end{attention}

\newpage

\coursec{Corrigés des exercices}

\begin{corrige}[Exercice 1 — QCM de compréhension orale (Script 1)]
\textbf{Réponses :}
\begin{enumerate}
\item \textbf{b) Bamenda} — \eng{“a campaign launched in Bamenda by the association Women in Action”}.
\item \textbf{c) 45} — \eng{“forty-five volunteers have visited sixty villages”}.
\item \textbf{b) 60 villages} — même phrase ; attention à ne pas confondre avec \eng{sixteen} (16).
\item \textbf{a) 1,200} — \eng{“the campaign has helped about 1,200 girls return to the classroom”}. \eng{1,200} se dit \eng{one thousand two hundred}.
\item \textbf{a) Mrs. Agnes Tabi} — \eng{“The project leader, Mrs. Agnes Tabi, says that…”}. \eng{Mrs.} (et non \eng{Miss}) est le titre d'une femme mariée.
\end{enumerate}
\emph{Points de vigilance :} en compréhension orale, justifiez toujours mentalement votre choix par une citation entendue ; méfiez-vous des distracteurs phonétiques (\eng{45 / 15}, \eng{60 / 16}) et des noms propres voisins (\eng{Tabi / Tambe}).
\end{corrige}

\begin{corrige}[Exercice 2 — Vrai ou faux justifié (Script 2)]
\begin{enumerate}
\item \textbf{False.} — \eng{“I joined the organisation Equal Future two years ago.”} (deux ans, pas trois).
\item \textbf{False.} — \eng{“Every Saturday, I run workshops for teenage boys about gender equality.”} (des garçons, pas des filles).
\item \textbf{True.} — \eng{“At first, only ten boys attended my workshops, but today there are more than fifty.”}
\item \textbf{True.} — \eng{“Last month, our organisation also trained thirty women to become community leaders.”}
\item \textbf{False.} — \eng{“I am proud because volunteering has changed my own vision of women's rights.”} (sa vision a changé).
\end{enumerate}
\emph{Méthode :} une réponse V/F non justifiée ne rapporte généralement que la moitié des points ; citez toujours le texte entre guillemets anglais.
\end{corrige}

\begin{corrige}[Exercice 3 — Discrimination phonétique \eng{-teen} / \eng{-ty}]
Exercice purement oral : la correction dépend de ce que lit l'enseignant. Proposition de série de dictée : \eng{13, 40, 50, 16, 70, 18, 90, 15}.
\par\emph{Critère de discrimination :}
\begin{itemize}
\item \eng{-teen} : accent tonique sur la \textbf{dernière} syllabe, voyelle longue : \eng{thir-TEEN} « ceur-tine », \eng{fif-TEEN} « fif-tine » ; le \eng{t} est net.
\item \eng{-ty} : accent tonique sur la \textbf{première} syllabe, voyelle courte : \eng{THIR-ty} « ceur-ti », \eng{FIF-ty} « fif-ti ».
\end{itemize}
\emph{Astuce :} si vous hésitez, demandez \eng{“Did you say fifteen or fifty?”} — reformuler est une stratégie d'écoute légitime.
\end{corrige}

\begin{corrige}[Exercice 4 — Phrases à trous (lexique)]
\begin{enumerate}
\item \eng{The NGO organised a \textbf{workshop} on women's rights in our school last Friday.}
\item \eng{Many parents refuse \textbf{to enrol} their daughters in secondary school.} — Après \eng{refuse}, on emploie l'infinitif avec \eng{to} : \eng{refuse to do something}.
\item \eng{A \textbf{volunteer} is a person who works for free to help a community.}
\item \eng{The campaign raises \textbf{awareness} about the importance of girls' education.} — Collocation : \eng{to raise awareness about/of}.
\item \eng{Gender \textbf{discrimination} is still a serious problem in some workplaces.}
\item \eng{Education is the key to women's \textbf{empowerment} in our society.}
\end{enumerate}
\emph{Point de vigilance :} \eng{awareness} est un nom indénombrable : jamais \eng{an awareness} ni \eng{awarenesses}.
\end{corrige}

\begin{corrige}[Exercice 5 — Traduction vers l'anglais]
\begin{enumerate}
\item \eng{The volunteers organised three workshops in our village.} (prétérit, action datée) ; \eng{have organised} est possible si l'on insiste sur le résultat présent.
\item \eng{Many girls do not go to school because of poverty.} — \eng{because of} + nom ; \eng{because} seul introduit une proposition (\eng{because they are poor}).
\item \eng{This association has been fighting discrimination against women since 2015.} — \eng{since} + date de départ $\rightarrow$ present perfect (continuous) ; \eng{discrimination against} (préposition \eng{against}).
\item \eng{We must sensitise parents to the importance of girls' education.} ou \eng{We must raise parents' awareness of the importance of girls' education.}
\item \eng{Volunteering has changed my vision of gender equality.} — \eng{volunteering} en sujet (gérondif) ; \eng{gender equality}, pas de faux ami avec « égalité des genres ».
\end{enumerate}
\end{corrige}

\begin{corrige}[Exercice 6 — Transformation de phrases]
\begin{enumerate}
\item \eng{The volunteers have not (haven't) visited sixty villages.} — Négation du present perfect : \eng{have/has + not + participe passé}.
\item \eng{How many workshops does Mrs. Tabi organise every week?} — Question au present simple : \eng{How many + nom pluriel + does + sujet + base verbale} (le \eng{-s} de la 3\up{e} personne passe sur \eng{does}).
\item \eng{The leader said that they would train thirty women the following month.} — Discours indirect : \eng{will} $\rightarrow$ \eng{would} ; \eng{we} $\rightarrow$ \eng{they} ; \eng{next month} $\rightarrow$ \eng{the following month}.
\item \eng{Women's rights should be respected (by people).} — Voix passive avec modal : \eng{should + be + participe passé}.
\item \eng{The campaign has been active since January.} (ou \eng{has been running since January}). — \eng{since} + point de départ $\rightarrow$ present perfect.
\end{enumerate}
\end{corrige}

\begin{corrige}[Exercice 7 — Texte à trous]
(1) \eng{volunteers} — (2) \eng{gender equality} — (3) \eng{girls} — (4) \eng{rights} — (5) \eng{awareness} — (6) \eng{school} — (7) \eng{training} — (8) \eng{community}.
\par\emph{Texte reconstitué :} \eng{In Yaoundé, a group of young volunteers works every weekend to promote gender equality. They visit families and explain that girls have the same rights as boys. Thanks to their awareness campaign, many parents now send their daughters to school. The group also offers training sessions to women who want to become leaders in their community.}
\par\emph{Points de vigilance :} \eng{the same rights as boys} (comparaison avec \eng{the same\dots as}) ; \eng{thanks to} + nom (« grâce à »).
\end{corrige}

\begin{corrige}[Exercice 8 — Appariements mots / définitions]
1-c (\eng{volunteer} : personne qui travaille sans paiement) ; 2-f (\eng{awareness} : connaissance d'un enjeu social) ; 3-b (\eng{workshop} : séance d'apprentissage pratique) ; 4-e (\eng{to enrol} : inscrire officiellement) ; 5-g (\eng{empowerment} : donner pouvoir et confiance) ; 6-d (\eng{discrimination} : traitement injuste) ; 7-a (\eng{to raise awareness} : rendre conscient) ; 8-h (\eng{campaign} : série d'actions organisées).
\par\emph{Astuce de mémorisation :} apprenez chaque mot avec sa collocation : \eng{launch a campaign}, \eng{run a workshop}, \eng{raise awareness}, \eng{fight discrimination}.
\end{corrige}

\begin{corrige}[Exercice 9 — Compréhension de l'écrit, type Bac]
\textbf{Barème indicatif (sur 20) :} A. Compréhension : 5 $\times$ 2 = 10 pts ; B. V/F justifié : 3 $\times$ 2 = 6 pts ; C. Langue : 2 + 1 + 1 = 4 pts.

\textbf{A. Comprehension questions}
\begin{enumerate}
\item \eng{Many girls leave school early because of poverty and early marriages.} (\eng{traditional beliefs} accepté comme deuxième raison.)
\item \eng{Most of the volunteers of “Bright Future” are university students.}
\item \eng{They organise their workshops in schools and churches.} (\eng{community halls} accepté ; deux lieux suffisent.)
\item \eng{Educated women earn better salaries and take better care of their children's health.} (\eng{they send their own daughters to school} accepté.)
\item \eng{Volunteering has taught Clarisse patience, courage and public speaking.}
\end{enumerate}

\textbf{B. True or false}
\begin{enumerate}
\item \textbf{False.} — \eng{“an association created in Buea in 2018”} (Buea, pas Douala).
\item \textbf{True.} — \eng{“The association has also trained eighty women as community leaders.”}
\item \textbf{False.} — \eng{“Some fathers refuse to listen to us”} et \eng{“attitudes slowly change”} (le changement est lent, pas immédiat).
\end{enumerate}

\textbf{C. Language work}
\begin{enumerate}
\item \emph{sensibiliser} = \eng{to sensitise} (paragraphe 2) ; \emph{encourageants} = \eng{encouraging} (paragraphe 3).
\item \eng{Three reading clubs have been opened by the association.} — Voix passive au present perfect : \eng{have been + participe passé}.
\item \eng{Clarisse says that some fathers refuse to listen to them.} (le verbe d'introduction au present permet de garder les temps) ; accepté aussi : \eng{Clarisse said that some fathers refused to listen to them.}
\end{enumerate}
\emph{Points de vigilance :} répondez par des phrases complètes avec sujet + verbe ; recopiez les citations sans les modifier ; au Bac, une justification fausse ou absente fait perdre la moitié du point.
\end{corrige}

\begin{corrige}[Exercice 10 — Traduction, type Bac]
\textbf{Traduction de référence :}

\eng{Since 2020, a Cameroonian association has been sending volunteers to villages to promote gender equality. These young people organise workshops and explain to parents that girls have the same rights as boys. Thanks to their work, more than a thousand girls have returned to school. Volunteering has also transformed the volunteers themselves: they have become more confident and more responsible.}

\textbf{Barème indicatif (sur 10) :} 2 pts par phrase (1 pt sens, 1 pt correction grammaticale).

\emph{Points de vigilance :}
\begin{itemize}
\item \eng{Since 2020} $\rightarrow$ present perfect (continuous) : \eng{has been sending}, jamais de présent simple.
\item « promouvoir l'égalité des sexes » = \eng{to promote gender equality} (pas \eng{gender equity} ici).
\item « les mêmes droits que » = \eng{the same rights as}.
\item « grâce à » = \eng{thanks to} ; « plus de mille » = \eng{more than a thousand}.
\item « eux-mêmes » = \eng{themselves} ; « confiants » = \eng{confident} (faux ami : \eng{confidential} = confidentiel).
\end{itemize}
\end{corrige}

\begin{corrige}[Exercice 11 — Composition, type Bac]
\textbf{Barème indicatif (sur 20) :} respect de la tâche et du nombre de mots (4 pts) ; organisation et connecteurs (4 pts) ; richesse du lexique du thème (4 pts) ; correction grammaticale (4 pts) ; orthographe et ponctuation (4 pts).

\textbf{Proposition de rédaction modèle — Sujet 1 (article, 170 mots) :}

\emph{\eng{Join the Movement for Gender Equality!}}

\eng{In many villages of our country, girls still leave school earlier than boys because of poverty and traditional beliefs. This discrimination is unfair, and we, the students of this school, can help to end it.}

\eng{Last year, the association Women in Action launched a campaign in the North West region. Forty-five volunteers organised workshops in sixty villages and helped about 1,200 girls return to the classroom. Their example shows that young people can really make a difference.}

\eng{You too can act! Join a volunteering club, raise awareness among your friends, run a workshop on women's empowerment, or help parents enrol their daughters in school. Every small action counts.}

\eng{Gender equality is not only a women's issue: it concerns all of us. So do not wait — become a volunteer today, and together let us build a fairer society!}

\textbf{Points de vigilance (Sujet 1) :} titre accrocheur ; au moins cinq mots du lexique (\eng{campaign, volunteers, workshops, awareness, empowerment, enrol, discrimination}) ; au moins un chiffre ; impératifs pour l'appel à l'action (\eng{join, do not wait}).

\textbf{Proposition de plan — Sujet 2 (essai) :}
\begin{enumerate}
\item \emph{Introduction :} présenter le débat — \eng{Gender equality is often seen as a women's fight, but is it really?}
\item \emph{Paragraphe 1 :} les garçons transmettent les attitudes (\eng{Boys become fathers, husbands and leaders; if they respect girls at school, they will respect women at home.}) — exemple : les ateliers de Kofi au Ghana.
\item \emph{Paragraphe 2 :} l'égalité profite aussi aux garçons (\eng{Sharing domestic work and emotions makes boys freer and happier.}).
\item \emph{Conclusion :} \eng{In conclusion, gender equality can only be achieved if boys and girls work together.}
\end{enumerate}
\emph{Points de vigilance (Sujet 2) :} opinion claire dès l'introduction (\eng{I strongly agree that\dots}) ; connecteurs variés (\eng{first, moreover, however, in conclusion}) ; exemples concrets du contexte camerounais valorisés.
\end{corrige}

\newpage

\coursec{Fiche bilan}

\begin{popfigure}
\begin{tikzpicture}[mindmap, grow cyclic, every node/.style={concept, font=\small},
  root concept/.append style={concept color=popblue, font=\bfseries\small, text=white},
  level 1/.append style={level distance=3.4cm, sibling angle=51.4},
  level 1 concept/.append style={font=\scriptsize}]
\node[root concept]{Listening: Volunteering for Gender Equality}
  child[concept color=poporange]{ node {Theme: volunteering \& equality} }
  child[concept color=popgreen]{ node {Key words: volunteer, empower, raise awareness} }
  child[concept color=poppurple]{ node {Figures \& proper nouns} }
  child[concept color=poppink]{ node {Strategies: predict, listen, check} }
  child[concept color=popgold]{ node {Grammar: present perfect, gerunds} }
  child[concept color=popblue!70]{ node {Questions: T/F, MCQ, short answers} }
  child[concept color=popgreen!70]{ node {Reuse: short sentences, speaking} };
\end{tikzpicture}
\legende{Carte mentale de la leçon : écouter un texte sur le volontariat pour l'égalité des sexes.}
\end{popfigure}

\begin{bilan}
\textbf{1. Lexique clé du thème (à connaître par cœur)}
\begin{itemize}
  \item \eng{volunteering} (n.) : le volontariat, le bénévolat ; \eng{a volunteer} : un(e) bénévole ; \eng{to volunteer} : se porter volontaire.
  \item \eng{gender equality} : l'égalité des sexes ; \eng{gender gap} : l'écart entre les sexes ; \eng{discrimination} : la discrimination.
  \item \eng{to empower women} : autonomiser les femmes ; \eng{empowerment} : l'autonomisation.
  \item \eng{to raise awareness (of/about)} : sensibiliser (à) ; \eng{an awareness campaign} : une campagne de sensibilisation.
  \item \eng{to promote} : promouvoir ; \eng{to advocate for} : plaider en faveur de ; \eng{to fight against} : lutter contre.
  \item \eng{equal rights} : l'égalité des droits ; \eng{equal opportunities} : l'égalité des chances ; \eng{access to education} : l'accès à l'éducation.
  \item \eng{a community} : une communauté ; \eng{an NGO (non-governmental organisation)} : une ONG ; \eng{a project / an initiative} : un projet / une initiative.
  \item \eng{to take part in / to get involved in} : participer à, s'engager dans ; \eng{to make a difference} : changer les choses.
\end{itemize}

\textbf{2. Grammaire à revoir pour cette écoute}
\begin{center}
\begin{tabularx}{\linewidth}{|X|X|}\hline
\rowcolor{popblueL} \textbf{Règle} & \textbf{Exemple} \\ \hline
Present perfect (expérience, bilan) : \cle{have/has + participe passé} & \eng{She has volunteered for three years.} « Elle est bénévole depuis trois ans. » \\ \hline
\cle{for} + durée / \cle{since} + point de départ & \eng{since 2020} / \eng{for two years} \\ \hline
Gérondif après \eng{be involved in, be interested in, succeed in} & \eng{They are involved in promoting girls' education.} \\ \hline
\eng{to} + base verbale (but) & \eng{She joined the NGO to empower women.} « … pour autonomiser les femmes. » \\ \hline
Passif : \cle{be + participe passé} & \eng{Girls are encouraged to stay at school.} \\ \hline
\end{tabularx}
\end{center}

\textbf{3. Méthode express d'écoute (les 4 étapes)}
\begin{enumerate}
  \item \textbf{Avant} : lire les questions, repérer les mots-clés, prédire le thème et le type d'informations attendues (chiffres, noms, dates).
  \item \textbf{1\iere{} écoute} : saisir le thème général et les idées principales ; ne pas s'acharner sur un mot inconnu.
  \item \textbf{2\ieme{} écoute} : noter les détails précis (chiffres, noms propres, lieux) et répondre aux questions.
  \item \textbf{Après} : vérifier la cohérence des réponses, orthographe et grammaire ; réemployer le lexique dans deux ou trois phrases.
\end{enumerate}

\textbf{4. Pièges fréquents des francophones}
\begin{itemize}
  \item \eng{actually} = « en fait » (et non « actuellement » = \eng{currently}).
  \item \eng{to assist} = « aider » ; « assister à » = \eng{to attend}.
  \item \eng{a formation} n'existe pas : « une formation » = \eng{training}.
  \item Prononciation : \eng{volunteer} « vo-lon-tire », \eng{women} « oui-mine », \eng{awareness} « a-wère-nèss ».
\end{itemize}
\end{bilan}

\begin{aretenir}[Checklist avant le Bac]
\begin{itemize}
  \item $\square$ Je sais définir \eng{volunteering}, \eng{gender equality} et \eng{empowerment} en anglais.
  \item $\square$ Je connais au moins 10 mots du lexique du volontariat et de l'égalité des sexes.
  \item $\square$ Je lis les questions \textbf{avant} l'écoute et je souligne les mots-clés.
  \item $\square$ Je sais repérer les chiffres, les dates et les noms propres dans un enregistrement.
  \item $\square$ Je distingue \eng{for} (durée) et \eng{since} (point de départ) avec le present perfect.
  \item $\square$ J'utilise correctement le gérondif après \eng{be involved in} et \eng{be interested in}.
  \item $\square$ Je sais répondre aux vrai/faux en justifiant par une phrase du texte.
  \item $\square$ Je méfie des faux amis : \eng{actually}, \eng{assist}, \eng{training}.
  \item $\square$ Je peux réemployer le lexique dans 3 phrases personnelles (\eng{I would like to volunteer…}).
  \item $\square$ Je vérifie l'orthographe des noms propres et la grammaire de mes réponses courtes.
\end{itemize}
\end{aretenir}

\findecours

\end{document}
